% GENERATED by tools/symbols2tex.py -- do not hand-edit.
\(<\) & canonical strict ordering & P1-P4.2 & \hyperref[entry:symbol-registry-00085]{p.~\pageref{entry:symbol-registry-00085}} & P1.2 & Ambient strict order on declared numerical and order-theoretic objects; P4.3 is excluded because bare \(<\) there is object-language arithmetic syntax.  \\
\(<\) & canonical strict ordering & P4.4-P14 & \hyperref[entry:symbol-registry-00087]{p.~\pageref{entry:symbol-registry-00087}} & P4.4 & Ambient strict order on declared numerical and order-theoretic objects; not arithmetic object-language syntax.  \\
\(<\) & object-language arithmetic strict-order abbreviation & P4.3 & \hyperref[entry:symbol-registry-00086]{p.~\pageref{entry:symbol-registry-00086}} & P4.3 & Eliminable abbreviation \(r<t\defeq\exists z\,(r+Sz=t)\), with the least variable fresh for the displayed terms.  \\
\(>\) & strict ordering & global & \hyperref[entry:symbol-registry-00089]{p.~\pageref{entry:symbol-registry-00089}} & P1.2 & Contract list; strict dual of \(\leq\); ``positive'' means \(>0\) by the natural-numbers convention.  \\
\(0\) & CLL additive false unit & P8.6 & \hyperref[entry:symbol-registry-00566]{p.~\pageref{entry:symbol-registry-00566}} & P8.6 & Additive unit dual to top in the displayed CLL grammar; not the relevance base point or an ordinary numeral.  \\
\(0\) & System-T ground type and zero constant & P7.6 & \hyperref[entry:symbol-registry-00463]{p.~\pageref{entry:symbol-registry-00463}} & P7.6 & Syntactic category distinguishes the ground type from the term constant; both are fixed by sheet:system-t.  \\
\(0(c)=\varnothing,\ 1(c)=\{*\}\) & constant initial and terminal presheaves & P10.9 & \hyperref[entry:symbol-registry-00845]{p.~\pageref{entry:symbol-registry-00845}} & P10.9 & These denote presheaves; numeric object names of a particular base category have their separately declared base-object role.  \\
\(1_A,1_{\mathcal C}\) & identity arrow and identity functor & P10.7-P10.9 & \hyperref[entry:symbol-registry-00777]{p.~\pageref{entry:symbol-registry-00777}} & P10.7 & The subscript states the categorical level: an object for an arrow and a category for its identity functor.  \\
\(1_{\mathbb P}\) & weakest forcing condition & P6.3 & \hyperref[entry:symbol-registry-00362]{p.~\pageref{entry:symbol-registry-00362}} & P6.3 & Top element under the convention that q less-than-or-equal p means q is stronger.  \\
\(2^{<\omega},2^\omega\) & finite and infinite binary-sequence spaces & P6.2-P6.4 & \hyperref[entry:symbol-registry-00356]{p.~\pageref{entry:symbol-registry-00356}} & P6.2 & Finite strings and omega-sequences over two; the first carries prefix order and the second supplies branches or reals as locally declared.  \\
\(3\mathsf{SAT}\) & 3-CNF satisfiability language & P4.6 & \hyperref[entry:symbol-registry-00295]{p.~\pageref{entry:symbol-registry-00295}} & P4.6 & Encodings of satisfiable 3-CNF formulas whose clauses contain at most three literals.  \\
\(A^*\) & Girard phase interpretation of a CLL formula & P8.6 & \hyperref[entry:symbol-registry-00572]{p.~\pageref{entry:symbol-registry-00572}} & P8.6 & Fact assigned to formula A in Girard's source notation; phase validity is membership of the monoid unit 1 in A-star.  \\
\(!A\) & bang exponential and its phase interpretation & P8.6 & \hyperref[entry:symbol-registry-00105]{p.~\pageref{entry:symbol-registry-00105}} & P8.6 & Prefix connective on the source-bounded IMLL extension and full CLL; the phase clause for \(!F\) interprets that same connective, and neither use is factorial.  \\
\(?A\) & why-not exponential & P8.6 & \hyperref[entry:symbol-registry-00567]{p.~\pageref{entry:symbol-registry-00567}} & P8.6 & CLL exponential dual to bang; prefix unary and the only side supporting the displayed structural rules.  \\
\(a\Rightarrow b\) & Heyting residual on algebra elements & P10.3 & \hyperref[entry:symbol-registry-00763]{p.~\pageref{entry:symbol-registry-00763}} & P10.3 & The greatest c with c meet a below b; source notation a-arrow-b denotes this same operation when operands are algebra elements.  \\
\(A\mathbin{\&}B\) & additive conjunction (with) & P8.6 & \hyperref[entry:symbol-registry-00564]{p.~\pageref{entry:symbol-registry-00564}} & P8.6 & CLL additive conjunction; distinct from metalinguistic ampersands.  \\
\(A\oplus B\) & additive disjunction (plus) & P8.6 & \hyperref[entry:symbol-registry-00565]{p.~\pageref{entry:symbol-registry-00565}} & P8.6 & CLL additive disjunction; exact P8.6 scope is disjoint from the P4.4 computable join role.  \\
\(A\oplus B\) & computable join representative & P4.4 & \hyperref[entry:symbol-registry-00262]{p.~\pageref{entry:symbol-registry-00262}} & P4.4 & Even-odd coded union.  \\
\(a\leq b\) & derived relevance-frame preorder & P8.6 & \hyperref[entry:symbol-registry-00559]{p.~\pageref{entry:symbol-registry-00559}} & P8.6 & Metalinguistic abbreviation for R(0,a,b) in the displayed pointed frame; distinct from the algebraic orders later in the chapter.  \\
\(A\approx B\) & equipotence of sets & P6.2 & \hyperref[entry:symbol-registry-00384]{p.~\pageref{entry:symbol-registry-00384}} & P6.2 & There is a bijection from A to B; distinct from isomorphism of structures and from formula equivalence.  \\
\(A\bindnasrepma B\) & multiplicative disjunction (par) & P8.6 & \hyperref[entry:symbol-registry-00563]{p.~\pageref{entry:symbol-registry-00563}} & P8.6 & CLL connective dual to tensor; shares the disjunction rank only in this exact scope.  \\
\(A\mathbin{\&}B\) & strong t-norm conjunction & P8.5 & \hyperref[entry:symbol-registry-00596]{p.~\pageref{entry:symbol-registry-00596}} & P8.5 & Primitive multiplicative conjunction of the MTL and BL-family sheets; distinct from weak lattice conjunction and CLL additive with.  \\
\(a\wedge b,a\vee b,0,1,a\Rightarrow b,\neg a\) & Heyting operations and bounds & P8.1 & \hyperref[entry:symbol-registry-00490]{p.~\pageref{entry:symbol-registry-00490}} & P8.1 & Source operations on elements of the locally declared Heyting algebra; the operand type, not a house subscript, distinguishes them from formula connectives.  \\
\(A\backslash B,B/A\) & Lambek residuals & P8.6 & \hyperref[entry:symbol-registry-00575]{p.~\pageref{entry:symbol-registry-00575}} & P8.6 & Left and right divisions of the ordered Lambek product; neither admits exchange weakening or contraction.  \\
\(A\mathbin{\cdot}B\) & Lambek product & P8.6 & \hyperref[entry:symbol-registry-00574]{p.~\pageref{entry:symbol-registry-00574}} & P8.6 & Associative unit-free product in Lambek L; operand type distinguishes it from phase and arithmetic multiplication.  \\
\(a\mathrel{:\equiv}b\) & definitional abbreviation & global & \hyperref[entry:symbol-registry-00032]{p.~\pageref{entry:symbol-registry-00032}} & P1.2 & The left expression is defined by the right expression; this sign is metalinguistic.  \\
\(A\leq_m^{\log}B\) & logspace many-one reducibility & P4.6 & \hyperref[entry:symbol-registry-00312]{p.~\pageref{entry:symbol-registry-00312}} & P4.6 & A deterministic logarithmic-work-space translator preserves membership.  \\
\(A\leq_m^pB\) & polynomial-time many-one reducibility & P4.6 & \hyperref[entry:symbol-registry-00292]{p.~\pageref{entry:symbol-registry-00292}} & P4.6 & A total deterministic polynomial-time translator preserves membership.  \\
\(A\alpha,\mathcal F\alpha\) & Davies--Humberstone actuality and fixedly operators & P8.2 & \hyperref[entry:symbol-registry-00508]{p.~\pageref{entry:symbol-registry-00508}} & P8.2 & Exact source unary operators in the two-index S5 presentation: A resets the lower evaluation index to the upper actual-world index and script F universally varies the upper index.  \\
\(A\mathbin{\mathsf S}B\) & temporal since operator & P9.4 & \hyperref[entry:symbol-registry-00646]{p.~\pageref{entry:symbol-registry-00646}} & P9.4 & Past-time binary operator with a witness at or before the current position and A at every strictly intervening position.  \\
\(A\mathbin{\mathsf U}B\) & temporal until & P8.3-P9.5 & \hyperref[entry:symbol-registry-00095]{p.~\pageref{entry:symbol-registry-00095}} & P8.3 & Binds tighter than every binary connective and looser than the unary rank; licenses the U-unfolding axiom's written form.  \\
\(a\leq_{\N}b\) & metalinguistic natural-number order & P4.3-P5.1 & \hyperref[entry:symbol-registry-00084]{p.~\pageref{entry:symbol-registry-00084}} & P4.3 & Used where bare \(\leq\) simultaneously denotes arithmetic-language syntax; operands are standard naturals or codes, never arithmetic terms.  \\
\(A\Box\!\!\to B\) & Lewis variably strict conditional & P8.4 & \hyperref[entry:symbol-registry-00591]{p.~\pageref{entry:symbol-registry-00591}} & P8.4 & Evaluated by nested similarity spheres; distinct from material implication and the unique-selection conditional.  \\
\(A{:}a\) & matrix value-labelled formula & P8.5 & \hyperref[entry:symbol-registry-00549]{p.~\pageref{entry:symbol-registry-00549}} & P8.5 & Tableau occurrence asserting that formula A has exactly matrix value a; distinct from term typing and world labels.  \\
\(A:\alpha\) & type ascription & P2.3-P7.6 & \hyperref[entry:symbol-registry-00204]{p.~\pageref{entry:symbol-registry-00204}} & P2.3 & The active typed calculus fixes the displayed term and type categories; distinct from the P7.7 labelled-formula colon.  \\
\(A>B\) & conditional connective & P8.4 & \hyperref[entry:symbol-registry-00590]{p.~\pageref{entry:symbol-registry-00590}} & P8.4 & Formula operands type this glyph as the primitive CK or Stalnaker conditional; numerical operands retain the strict-order role.  \\
\(A,B,C\) & default formula metavariables in inherited Part-8 prose & P8 & \hyperref[entry:symbol-registry-00118]{p.~\pageref{entry:symbol-registry-00118}} & P8.1 & Default only where the inherited or cited presentation uses these letters.  A named source presentation's own locally declared syntax overrides this row; this row never licenses renaming source metavariables.  \\
\(A,B,C,F,P,R\) & typed-term metavariables & P2.3 & \hyperref[entry:symbol-registry-00203]{p.~\pageref{entry:symbol-registry-00203}} & P2.3 & Used only under local type declarations in P2.3; disjoint from P8 formula metavariables and algebraic carriers.  \\
\(a,b,\langle a\rangle A,[a]A\) & Kozen program constants and indexed modalities & P8.3 & \hyperref[entry:symbol-registry-00530]{p.~\pageref{entry:symbol-registry-00530}} & P8.3 & Program constants index Kozen's primitive diamond; the box is its exact defined dual.  Automaton input letters use sigma to keep this role distinct.  \\
\(A^{\Box}\) & Goedel--McKinsey--Tarski modal translation & P8.2 & \hyperref[entry:symbol-registry-00502]{p.~\pageref{entry:symbol-registry-00502}} & P8.2 & Recursive IPL-to-S4 formula translation; distinct from applying the object-language box to an already translated formula.  \\
\(A^\circ\) & da Costa C1 consistency formula & P8.7 & \hyperref[entry:symbol-registry-00600]{p.~\pageref{entry:symbol-registry-00600}} & P8.7 & Eliminable abbreviation \(\neg(A\land\neg A)\) expressing formula-level well-behavedness in C1; not a matrix truth value or a metalinguistic theory-consistency claim.  \\
\(a<_{\N}b\) & metalinguistic natural-number strict order & P4.3 & \hyperref[entry:symbol-registry-00088]{p.~\pageref{entry:symbol-registry-00088}} & P4.3 & Used where bare \(<\) simultaneously denotes arithmetic-language syntax; operands are standard naturals or codes, never arithmetic terms.  \\
\(A^\perp\) & formula-level linear negation & P8.6 & \hyperref[entry:symbol-registry-00568]{p.~\pageref{entry:symbol-registry-00568}} & P8.6 & Involutive De Morgan operation on CLL formulas; distinct from phase orthogonality on subsets.  \\
\([A]_T\) & Turing-equivalence class & P4.4 & \hyperref[entry:symbol-registry-00261]{p.~\pageref{entry:symbol-registry-00261}} & P4.4 & Alternate display of the degree of \(A\).  \\
\(\aleph_\alpha,\beth_\alpha\) & aleph and beth cardinal hierarchies & P6.3-P6.8 & \hyperref[entry:symbol-registry-00351]{p.~\pageref{entry:symbol-registry-00351}} & P6.3 & Aleph enumerates infinite initial ordinals and beth iterates powerset cardinality; limit stages take suprema.  \\
\(\alpha\) & Davies--Humberstone formula metavariable & P8.2 & \hyperref[entry:symbol-registry-00510]{p.~\pageref{entry:symbol-registry-00510}} & P8.2 & Ranges over formulas of the source two-index S5 language; disjoint from the CTL-star path-formula role in P8.3.  \\
\(\alpha\to_{\mathrm{ty}}\beta\) & object-language function-type former & P2.3-P7.6 & \hyperref[entry:symbol-registry-00202]{p.~\pageref{entry:symbol-registry-00202}} & P2.3 & Right-associative type constructor in the active typed calculus; distinct from object-language implication.  \\
\(\alpha\to_{\mathrm{ty}}\beta\) & object-language function-type former & P12.5 & \hyperref[entry:symbol-registry-00630]{p.~\pageref{entry:symbol-registry-00630}} & P12.5 & Right-associative arrow of the exact reused STLC syntax.  \\
\(\alpha\mathbin{\times_{\mathrm{ty}}}\beta,\alpha\mathbin{+_{\mathrm{ty}}}\beta\) & product and sum type formers in the categorical STLC extension & P7.3 & \hyperref[entry:symbol-registry-00433]{p.~\pageref{entry:symbol-registry-00433}} & P7.3 & Exact added type formers used to construct finite products and coproducts; distinct from arithmetic operations.  \\
\([\alpha]A\) & PDL universal program modality & P9.6 & \hyperref[entry:symbol-registry-00668]{p.~\pageref{entry:symbol-registry-00668}} & P9.6 & True when every alpha-successor satisfies A; the index is a regular program rather than a modal frame label.  \\
\(\alpha,\beta\) & CTL-star path-formula metavariables & P8.3 & \hyperref[entry:symbol-registry-00532]{p.~\pageref{entry:symbol-registry-00532}} & P8.3 & Range over Reynolds's path-formula category after the mutually generated CTL-star state/path syntax is declared; disjoint from ordinal and type roles in earlier parts.  \\
\(\alpha,\beta\) & finite-type metavariables & P7.6 & \hyperref[entry:symbol-registry-00459]{p.~\pageref{entry:symbol-registry-00459}} & P7.6 & Finite types generated from ground type zero by the registered type arrow; disjoint from P7.5 ordinals.  \\
\(\alpha,\beta\) & ordinal metavariables & P6 & \hyperref[entry:symbol-registry-00019]{p.~\pageref{entry:symbol-registry-00019}} & P6.1 & Ordinals of the ambient set theory; recurs under a local declaration in P7.5 ordinal analysis; disjoint from the type-metavariable scope.  \\
\(\alpha,\beta\) & ordinal metavariables & P7.5 & \hyperref[entry:symbol-registry-00450]{p.~\pageref{entry:symbol-registry-00450}} & P7.5 & Ordinals in the local proof-theoretic calibration; disjoint from the P7.3-P7.4 and P7.6 type scopes.  \\
\(\alpha;\beta,\alpha\cup\beta,\alpha^*,A?\) & PDL program constructors & P9.6 & \hyperref[entry:symbol-registry-00667]{p.~\pageref{entry:symbol-registry-00667}} & P9.6 & Relational composition choice reflexive-transitive closure and identity test in the mutually generated PDL grammar.  \\
\(\alpha+\beta,\alpha\cdot\beta,\alpha^\beta\) & ordinal arithmetic operations & P6.1-P7.5 & \hyperref[entry:symbol-registry-00383]{p.~\pageref{entry:symbol-registry-00383}} & P6.1 & Right-recursive ordinal sum product and exponentiation; distinct from cardinal arithmetic by the declared operands and chapter scope.  \\
\(\alpha,\beta,\gamma\) & type metavariables & P7.3-P7.4 & \hyperref[entry:symbol-registry-00020]{p.~\pageref{entry:symbol-registry-00020}} & P7.3 & Scope is disjoint from the P7.5 ordinal role and the P7.6 finite-type role.  \\
\(\alpha,\beta,\gamma\) & type metavariables & P2.3 & \hyperref[entry:symbol-registry-00200]{p.~\pageref{entry:symbol-registry-00200}} & P2.3 & Local to the simple-type grammar; disjoint from ordinal metavariables and the separate P7.4 type scope.  \\
\(\alpha,\beta,\gamma\) & type metavariables & P12.5 & \hyperref[entry:symbol-registry-00629]{p.~\pageref{entry:symbol-registry-00629}} & P12.5 & Types in the reused STLC and lambda-cube examples; this scope is disjoint from ordinal and path-formula roles.  \\
\(\alpha_i\) & function variable in the analytical hierarchy & P4.5 & \hyperref[entry:symbol-registry-00272]{p.~\pageref{entry:symbol-registry-00272}} & P4.5 & Scope disjoint from ordinal and type metavariables.  \\
\(\alpha^{\mathrm d}\) & NNF dual operation & P2.1 & \hyperref[entry:symbol-registry-00135]{p.~\pageref{entry:symbol-registry-00135}} & P2.1 & Metalinguistic involution fixing literals and exchanging conjunction with disjunction and falsum with verum.  \\
\(\ast,\Box;(s_1,s_2)\) & Barendregt lambda-cube sorts and product-sort pair & P12.5 & \hyperref[entry:symbol-registry-00623]{p.~\pageref{entry:symbol-registry-00623}} & P12.5 & Exact source notation: the pair abbreviates the triple (s1,s2,s2); the locally declared PTS role distinguishes Box and star from modal and algebraic uses.  \\
\(B'\) & Turing jump & P4.4 & \hyperref[entry:symbol-registry-00053]{p.~\pageref{entry:symbol-registry-00053}} & P4.4 & Postfix prime on an oracle set.  \\
\(b\vee_{\mathbb B}c\) & Boolean-algebra join & P6.3 & \hyperref[entry:symbol-registry-00327]{p.~\pageref{entry:symbol-registry-00327}} & P6.3 & Operation in the declared algebra; distinct from the globally registered object-language disjunction.  \\
\(b\wedge_{\mathbb B}c\) & Boolean-algebra meet & P6.3 & \hyperref[entry:symbol-registry-00326]{p.~\pageref{entry:symbol-registry-00326}} & P6.3 & Operation in the declared algebra; distinct from the globally registered object-language conjunction.  \\
\(B^A,\operatorname{ev}_{A,B},\Lambda(f)\) & exponential object evaluation and transpose & P10.7-P10.9 & \hyperref[entry:symbol-registry-00783]{p.~\pageref{entry:symbol-registry-00783}} & P10.7 & Evaluation has domain \(B^A\times A\); transposition sends \(f:X\times A\to B\) to the unique corresponding arrow \(X\to B^A\).  \\
\(B[A/x]\) & capture-avoiding typed substitution & P2.3 & \hyperref[entry:symbol-registry-00209]{p.~\pageref{entry:symbol-registry-00209}} & P2.3 & Substitutes the typed term A for free occurrences of x in B after the required alpha-renaming.  \\
\(b^{\complement_{\mathbb B}}\) & Boolean-algebra complement & P6.3 & \hyperref[entry:symbol-registry-00325]{p.~\pageref{entry:symbol-registry-00325}} & P6.3 & Postfix operation in the declared algebra; the subscript prevents collision with object-language negation.  \\
\(B_iA\) & agent-indexed belief operator & P8.3 & \hyperref[entry:symbol-registry-00534]{p.~\pageref{entry:symbol-registry-00534}} & P8.3 & Doxastic modality for agent i on the registered \(\mathsf{KD45}_{I}\) sheet; distinct from factive \(K_i\) knowledge.  \\
\(B^{+}_{\mathrm{at}},\mathcal M_{B^{+}_{\mathrm{at}}}\) & positive atomic branch and tableau countermodel & P2.2 & \hyperref[entry:symbol-registry-00164]{p.~\pageref{entry:symbol-registry-00164}} & P2.2 & The positive ground atoms on a fair saturated branch and the induced active-language Herbrand structure.  \\
\(B^{(n)}\) & n-fold Turing jump & P4.4-P5 & \hyperref[entry:symbol-registry-00257]{p.~\pageref{entry:symbol-registry-00257}} & P4.4 & Iteration from an arbitrary oracle set.  \\
\(\bar E\) & NNF complement of an ALC concept & P9.3 & \hyperref[entry:symbol-registry-00694]{p.~\pageref{entry:symbol-registry-00694}} & P9.3 & The recursive involution complements literals and constants and exchanges Boolean and modal duals; its extension is the complement in the interpretation domain.  \\
\(\bar n\) & numeral & P4.3-P5 & \hyperref[entry:symbol-registry-00056]{p.~\pageref{entry:symbol-registry-00056}} & P4.3 & Closed arithmetic term \(S^n0\); bars are reserved here for numerals, while tuples use vectors.  \\
\(\bar t,\bar\sigma\) & one-component term and basic-operation arrows & P10.7 & \hyperref[entry:symbol-registry-00806]{p.~\pageref{entry:symbol-registry-00806}} & P10.7 & An \(n\)-variable term determines an arrow \(n\to1\); a symbol of arity \(k\) gives its basic term arrow \(k\to1\).  \\
\(\bar v\) & complemented valuation & P2.1 & \hyperref[entry:symbol-registry-00136]{p.~\pageref{entry:symbol-registry-00136}} & P2.1 & Defined by \(\bar v(p)=1-v(p)\); used only in the propositional duality theorem.  \\
\(\beta_B\) & Boolean representation isomorphism & P10.2 & \hyperref[entry:symbol-registry-00720]{p.~\pageref{entry:symbol-registry-00720}} & P10.2 & Maps a Boolean element a to its basic clopen hat-a in the Stone space X-B.  \\
\(\beta_H(a)\) & Heyting embedding into upsets of prime filters & P10.3 & \hyperref[entry:symbol-registry-00768]{p.~\pageref{entry:symbol-registry-00768}} & P10.3 & The set of prime filters containing a; the implication witness proves preservation of the residual.  \\
\(\beta(w,i)\) & Goedel beta function & P4.3 & \hyperref[entry:symbol-registry-00252]{p.~\pageref{entry:symbol-registry-00252}} & P4.3 & Numerical finite-sequence decoder.  \\
\(\beth_n(\kappa)\) & finite beth iteration above a cardinal & P6.4 & \hyperref[entry:symbol-registry-00366]{p.~\pageref{entry:symbol-registry-00366}} & P6.4 & Starts at kappa and iterates powerset cardinality n times; distinct from the global beth hierarchy indexed from aleph-zero.  \\
\(\bigvee\Delta\) & finite context disjunction & P7.1 & \hyperref[entry:symbol-registry-00393]{p.~\pageref{entry:symbol-registry-00393}} & P7.1 & Metalinguistic fold of a finite context with empty value bottom; distinct from the P2.4 infinitary object-language connective.  \\
\(\bigvee_{i\in I}\varphi_i\) & countable infinitary disjunction & P2.4 & \hyperref[entry:symbol-registry-00222]{p.~\pageref{entry:symbol-registry-00222}} & P2.4 & I is countable; the connective retains the ordinary disjunction precedence rank.  \\
\(\bigwedge\Gamma\) & finite context conjunction & P7.1 & \hyperref[entry:symbol-registry-00392]{p.~\pageref{entry:symbol-registry-00392}} & P7.1 & Metalinguistic fold of a finite context with empty value top; distinct from the P2.4 infinitary object-language connective.  \\
\(\bigwedge_{i\in I}\varphi_i\) & countable infinitary conjunction & P2.4 & \hyperref[entry:symbol-registry-00221]{p.~\pageref{entry:symbol-registry-00221}} & P2.4 & I is countable; the connective retains the ordinary conjunction precedence rank.  \\
\(\boldsymbol{\Sigma}^{0}_{\alpha},\boldsymbol{\Pi}^{0}_{\alpha},\boldsymbol{\Delta}^{0}_{\alpha}\) & boldface Borel pointclasses & P6.7 & \hyperref[entry:symbol-registry-00334]{p.~\pageref{entry:symbol-registry-00334}} & P6.7 & Countable-ordinal Borel levels on a Polish space; distinct from the P4-P5 bold arithmetical-definability classes.  \\
\(\boldsymbol{\Sigma}^0_n,\boldsymbol{\Pi}^0_n,\boldsymbol{\Delta}^0_n\) & arithmetical definability classes of relations & P4-P5 & \hyperref[entry:symbol-registry-00018]{p.~\pageref{entry:symbol-registry-00018}} & P4.3 & Bold symbols classify relations definable in \(\mathcal N_{\mathrm{ar}}\), not formula strings or Borel pointclasses.  \\
\(\boldsymbol{\Sigma}^{1}_{n},\boldsymbol{\Pi}^{1}_{n},\boldsymbol{\Delta}^{1}_{n}\) & boldface projective pointclasses & P6.6-P6.8 & \hyperref[entry:symbol-registry-00335]{p.~\pageref{entry:symbol-registry-00335}} & P6.6 & Real-parameter projective classes; distinct from lightface analytical classes and from the Levy hierarchy.  \\
\(\bot,\top\) & falsum and verum or declared CLL nullary units & global & \hyperref[entry:symbol-registry-00038]{p.~\pageref{entry:symbol-registry-00038}} & P2.1 & Outside the P8.6 CLL grammar these are bottom and the defined top abbreviation; in that grammar \(\bot\) is the par unit and \(\top\) is the with unit, with the ambient syntax fixing the role.  \\
\(\Box\) & necessity operator & P5-P8.3 & \hyperref[entry:symbol-registry-00091]{p.~\pageref{entry:symbol-registry-00091}} & P5.3 & Prefix unary at the negation rank; modal and provability readings are fixed by the active system.  \\
\(\Box A\) & Anderson--Belnap relevance abbreviation & P8.6 & \hyperref[entry:symbol-registry-00092]{p.~\pageref{entry:symbol-registry-00092}} & P8.6 & Exact eliminable abbreviation \((A\to A)\to A\) on the E and R sheets; it is not a primitive modality.  \\
\(\Box_A,\Diamond_A\) & dual operations of a normal modal algebra & P10.3 & \hyperref[entry:symbol-registry-00740]{p.~\pageref{entry:symbol-registry-00740}} & P10.3 & Box preserves finite meets and one; diamond is its Boolean dual and preserves finite joins and zero.  \\
\(C\sqsubseteq D\) & description-logic concept inclusion & P9.3 & \hyperref[entry:symbol-registry-00639]{p.~\pageref{entry:symbol-registry-00639}} & P9.3 & An object-language TBox axiom interpreted by subset inclusion of concept extensions; not metalinguistic formula consequence.  \\
\(c_f,u,\mathsf{pair}_{A,B}\) & named arrow terminal and pairing constants & P10.8 & \hyperref[entry:symbol-registry-00790]{p.~\pageref{entry:symbol-registry-00790}} & P10.8 & Closed constants in the expanded internal language; their declared applications denote the corresponding arrows and product pairing.  \\
\(C_G\) & common knowledge & P8.3 & \hyperref[entry:symbol-registry-00101]{p.~\pageref{entry:symbol-registry-00101}} & P8.3 & Prefix unary at the negation rank; satisfaction via reflexive-transitive reachability under the union of the \(\sim_i\).  \\
\(c_\Gamma(\beta),\Delta^\Gamma\) & curried type and context exponential & P10.8 & \hyperref[entry:symbol-registry-00787]{p.~\pageref{entry:symbol-registry-00787}} & P10.8 & Right-associated arrows over Gamma's type list; the exponential list has one such curried type per component of Delta.  \\
\(C_{M,w},\mathcal T_{M,s}\) & ALC concept and TBox encoding alternating acceptance & P9.3 & \hyperref[entry:symbol-registry-00702]{p.~\pageref{entry:symbol-registry-00702}} & P9.3 & The concept fixes the input configuration and maximal binary rank; the finite TBox enforces configurations branching transitions and countdown.  \\
\(c_{\mathrm{init}}(x),c_{\mathrm{acc}}\) & initial and canonical accepting configuration blocks & P9.8 & \hyperref[entry:symbol-registry-00688]{p.~\pageref{entry:symbol-registry-00688}} & P9.8 & Boolean constant blocks substituted into the final closed QBF; not first-order individual constants.  \\
\(c_n^M\) & canonical model power comparison & P10.7 & \hyperref[entry:symbol-registry-00805]{p.~\pageref{entry:symbol-registry-00805}} & P10.7 & The bijection \(M(n)\to M(1)^n\) has coordinates \(M(p_i^n)\); the comparison at one uses singleton tuples.  \\
\(C_V(x)\) & plain Kolmogorov complexity & P4.5 & \hyperref[entry:symbol-registry-00284]{p.~\pageref{entry:symbol-registry-00284}} & P4.5 & Relative to decompressor \(V\).  \\
\(c(\varphi)\) & formula complexity & P2.1-P2.2 & \hyperref[entry:symbol-registry-00128]{p.~\pageref{entry:symbol-registry-00128}} & P2.1 & Recursively counts constructor occurrences under the active propositional or first-order grammar.  \\
\((\cdot)^\sharp\) & incidence translation on formulas contexts structures and assignments & P2.3 & \hyperref[entry:symbol-registry-00196]{p.~\pageref{entry:symbol-registry-00196}} & P2.3 & A metalinguistic tagging and translation operation; the operand determines its eliminable clause.  \\
\(\check x\) & check name & P6.3 & \hyperref[entry:symbol-registry-00070]{p.~\pageref{entry:symbol-registry-00070}} & P6.3 & Canonical forcing name of a ground-model set.  \\
\(\chi_A\) & characteristic function of a set & P4.1-P4.4 & \hyperref[entry:symbol-registry-00238]{p.~\pageref{entry:symbol-registry-00238}} & P4.1 & Source-bound notation: one exactly on members of \(A\).  \\
\([!\chi]A,\langle!\chi\rangle A\) & public-announcement modalities & P8.3 & \hyperref[entry:symbol-registry-00539]{p.~\pageref{entry:symbol-registry-00539}} & P8.3 & Truthful-update box and its executable dual; a false box precondition is vacuously true.  \\
\(\chi_A:X\to\Omega\) & characteristic map of a subpresheaf & P10.9 & \hyperref[entry:symbol-registry-00851]{p.~\pageref{entry:symbol-registry-00851}} & P10.9 & At \(x\in X(c)\), it is the sieve of arrows \(u:d\to c\) whose restriction \(X(u)(x)\) belongs to \(A(d)\).  \\
\(\chi_R\) & characteristic function of a relation & P4.1-P4.4 & \hyperref[entry:symbol-registry-00239]{p.~\pageref{entry:symbol-registry-00239}} & P4.1 & Source-bound notation: one exactly on tuples in \(R\).  \\
\(D\) & designated values & P8.5 & \hyperref[entry:symbol-registry-00112]{p.~\pageref{entry:symbol-registry-00112}} & P8.5 & Subset of the active matrix value set fixing matrix consequence.  \\
\(D_\alpha\) & semantic domain at type alpha & P2.3 & \hyperref[entry:symbol-registry-00207]{p.~\pageref{entry:symbol-registry-00207}} & P2.3 & Nonempty type domain; standard function types contain all functions and Henkin domains may select a closure-stable subset.  \\
\(D_d\) & finite set decoded from d & P4.4 & \hyperref[entry:symbol-registry-00269]{p.~\pageref{entry:symbol-registry-00269}} & P4.4 & Positions of one-bits in \(d\).  \\
\(D,D_1,D_2\) & derivation metavariables & P7.2-P7.6 & \hyperref[entry:symbol-registry-00394]{p.~\pageref{entry:symbol-registry-00394}} & P7.2 & Finite derivation trees in the locally fixed calculus; distinct from the roman Dialectica superscript.  \\
\(D,D^{\mathrm{in}},D^{\mathrm{out}}\) & free-logic outer and inner domains & P2.4-P8.4 & \hyperref[entry:symbol-registry-00217]{p.~\pageref{entry:symbol-registry-00217}} & P2.4 & Local outer-domain presentations; the inner existing domain is nonempty and contained in the outer domain.  \\
\(D_GA\) & distributed-knowledge operator & P8.3 & \hyperref[entry:symbol-registry-00535]{p.~\pageref{entry:symbol-registry-00535}} & P8.3 & Truth at every successor of the intersection of the group relations; distinct from everyone-knows and common knowledge.  \\
\(D,S_n,E_n,\widehat P\) & incidence-translation signature data & P2.3 & \hyperref[entry:symbol-registry-00195]{p.~\pageref{entry:symbol-registry-00195}} & P2.3 & \(D\) is the object sort; \(S_n\) the n-ary relation sort; \(E_n\) incidence; and \(\widehat P\) the relation-sort constant representing predicate \(P\).  \\
\(d(\varphi),d(\Gamma)\) & defined-negation formula and context translation & P7.1 & \hyperref[entry:symbol-registry-00391]{p.~\pageref{entry:symbol-registry-00391}} & P7.1 & Recursive equality-free translation from canonical LK syntax to the defined-negation language; context action is occurrencewise.  \\
\(D_w\) & world-indexed intuitionistic domain & P8.1 & \hyperref[entry:symbol-registry-00493]{p.~\pageref{entry:symbol-registry-00493}} & P8.1 & Nonempty domain at world w in expanding-domain IQC semantics; monotone along the intuitionistic preorder.  \\
\(\deg_T(A)\) & Turing degree of a set & P4.4 & \hyperref[entry:symbol-registry-00260]{p.~\pageref{entry:symbol-registry-00260}} & P4.4 & Equivalence class under Turing equivalence.  \\
\(\delta\) & Turing-machine transition function & P4.1-P4.2 & \hyperref[entry:symbol-registry-00232]{p.~\pageref{entry:symbol-registry-00232}} & P4.1 & Typed by the fixed machine tuple.  \\
\(\Delta^0_1\) & Delta-zero-one arithmetic formula class & P7.5 & \hyperref[entry:symbol-registry-00455]{p.~\pageref{entry:symbol-registry-00455}} & P7.5 & Formulas both Sigma-zero-one and Pi-zero-one under the fixed arithmetic syntax; used in RCA-zero comprehension.  \\
\(\Delta_0,\Sigma_n,\Pi_n,\Delta_n\) & Levy-hierarchy formula classes & P6.6 & \hyperref[entry:symbol-registry-00373]{p.~\pageref{entry:symbol-registry-00373}} & P6.6 & Set-theoretic bounded and unbounded-quantifier classes; distinct from arithmetical and projective hierarchies.  \\
\(\delta_A\) & internal equality predicate & P10.9 & \hyperref[entry:symbol-registry-00819]{p.~\pageref{entry:symbol-registry-00819}} & P10.9 & An element of \(P(A\times A)\) satisfying the full parameterized diagonal adjunction; no subobject representation is assumed.  \\
\(\Delta_A,\tau_A\) & diagonal and exchange of an equality pair & P10.9 & \hyperref[entry:symbol-registry-00822]{p.~\pageref{entry:symbol-registry-00822}} & P10.9 & The arrows \(A\to A\times A\) and \(A\times A\to A\times A\) used to state reflexivity and symmetry.  \\
\(\Delta_{\mathrm{ty}}\) & System-F type-variable context & P7.4 & \hyperref[entry:symbol-registry-00415]{p.~\pageref{entry:symbol-registry-00415}} & P7.4 & Finite list of distinct type variables; distinct from formula and sequent contexts.  \\
\(\Delta_{\mathrm{ty}}\derives_{\mathsf F}\alpha\ \mathsf{type}\) & System-F type-formation judgment & P7.4 & \hyperref[entry:symbol-registry-00420]{p.~\pageref{entry:symbol-registry-00420}} & P7.4 & Formation judgment for the least type set fixed by sheet:system-f.  \\
\(\Delta_{\mathrm{ty}};\Theta_{\mathrm{ty}}\derives_{\mathsf F}M{:}\alpha\) & System-F term-typing judgment & P7.4 & \hyperref[entry:symbol-registry-00421]{p.~\pageref{entry:symbol-registry-00421}} & P7.4 & Church-style typing judgment with separate type-variable and term-variable contexts.  \\
\(\derives_{\mathrm{MTL}},\derives_{\mathrm{BL}},\derives_G,\derives_{\Pi}\) & t-norm Hilbert derivability relations & P8.5 & \hyperref[entry:symbol-registry-00598]{p.~\pageref{entry:symbol-registry-00598}} & P8.5 & The subscript fixes one exact sheet; their standard completeness strengths are stated separately.  \\
\(\derives_{\mathsf{par}},\derives_{\mathsf{tot}}\) & partial and total Hoare derivability tags & P9.6 & \hyperref[entry:symbol-registry-00673]{p.~\pageref{entry:symbol-registry-00673}} & P9.6 & Judgments on sheet:hoare-while; only the total tag licenses the decreasing-variant loop rule.  \\
\(\derives_{\mathsf{PDL}}A\) & PDL theoremhood judgment & P9.6 & \hyperref[entry:symbol-registry-00669]{p.~\pageref{entry:symbol-registry-00669}} & P9.6 & Finite no-premise derivability on sheet:hilbert-pdl.  \\
\(\Diamond\) & possibility operator & P5-P8.3 & \hyperref[entry:symbol-registry-00093]{p.~\pageref{entry:symbol-registry-00093}} & P5.3 & Prefix unary at the negation rank; \(\neg\Box\neg\) where defined.  \\
\(\diamondsuit_{\omega_1}\) & diamond principle at omega-one & P6.4 & \hyperref[entry:symbol-registry-00370]{p.~\pageref{entry:symbol-registry-00370}} & P6.4 & Stationary-guessing sequence defined in the combinatorial-principles entry.  \\
\(\downarrow S,\uparrow x\) & downward closure and principal upset & P10.3 & \hyperref[entry:symbol-registry-00769]{p.~\pageref{entry:symbol-registry-00769}} & P10.3 & Order closures in a Priestley or Esakia space; the downclosure condition applies to every clopen subset.  \\
\(e\mapsto(D_e,\prec_e)\) & effective notation-system enumeration & P7.5 & \hyperref[entry:symbol-registry-00453]{p.~\pageref{entry:symbol-registry-00453}} & P7.5 & Maps each natural-number code to its recursive domain and ordering relation in the relative ordinal calibration.  \\
\(e_{A,B},l_{A,B}\) & comparison constants for a named exponential & P10.8 & \hyperref[entry:symbol-registry-00791]{p.~\pageref{entry:symbol-registry-00791}} & P10.8 & Compare the base type naming \(B^A\) with the arrow type between the named types for \(A\) and \(B\); both unary interpretations are identities.  \\
\(E,EM,EC,EN,EMC,EMCN\) & Chellas compositional modal logic identifiers & P8.2 & \hyperref[entry:symbol-registry-00498]{p.~\pageref{entry:symbol-registry-00498}} & P8.2 & Chellas's exact plain math-italic compositional system names.  M C and N remain schema labels rather than substitute system identifiers; normal K is EMCN.  \\
\(E_G\) & mutual knowledge & P8.3 & \hyperref[entry:symbol-registry-00100]{p.~\pageref{entry:symbol-registry-00100}} & P8.3 & \(E_GA:\equiv\bigwedge_{i\in G}K_iA\), prefix unary at the negation rank.  \\
\(e_{I,A},p_{I,A},q_{I,A},r_{I,A}\) & maps for equality with parameters & P10.9 & \hyperref[entry:symbol-registry-00820]{p.~\pageref{entry:symbol-registry-00820}} & P10.9 & Use \((I\times A)\times A\): duplicate the final variable, select either copy with its parameters, or select the equality pair.  \\
\((E_i,F_i)_{i<k}\) & Rabin acceptance pairs & P9.4 & \hyperref[entry:symbol-registry-00648]{p.~\pageref{entry:symbol-registry-00648}} & P9.4 & A run accepts when one pair has no infinitely recurring E-state and at least one infinitely recurring F-state.  \\
\(E!(t)\) & free-logic existence predicate & P2.4-P8.4 & \hyperref[entry:symbol-registry-00218]{p.~\pageref{entry:symbol-registry-00218}} & P2.4 & True exactly when the term value lies in the declared inner existing domain of the displayed free-logic semantics.  \\
\(E(x)\supset\Phi\) & Fitting object-language conditional & P8.2 & \hyperref[entry:symbol-registry-00517]{p.~\pageref{entry:symbol-registry-00517}} & P8.2 & Exact source conditional inside the actualist universal-quantifier abbreviation; locally typed apart from the handbook's default implication glyph.  \\
\(E(x,y),F_f(\vec x,y)\) & interpreted equality and function-graph formulas & P5.4 & \hyperref[entry:symbol-registry-00319]{p.~\pageref{entry:symbol-registry-00319}} & P5.4 & The target-language equality surrogate and total functional graph data fixed at def:p05-relative-interpretation; constants are nullary functions.  \\
\(\equiv_E^n\) & equational congruence in a finite context & P10.7 & \hyperref[entry:symbol-registry-00801]{p.~\pageref{entry:symbol-registry-00801}} & P10.7 & The least family containing the stated equations and closed under operation compatibility and simultaneous substitutions between all finite contexts.  \\
\(\equiv_m\) & many-one equivalence & P4.4 & \hyperref[entry:symbol-registry-00258]{p.~\pageref{entry:symbol-registry-00258}} & P4.4 & Exact source notation for mutual many-one reducibility.  \\
\(\eta_D\) & Priestley lattice representation & P10.3 & \hyperref[entry:symbol-registry-00758]{p.~\pageref{entry:symbol-registry-00758}} & P10.3 & Maps each lattice element to the clopen upset of prime filters containing it.  \\
\(\eta,\varepsilon\) & unit and counit of an adjunction & P10.7-P10.9 & \hyperref[entry:symbol-registry-00780]{p.~\pageref{entry:symbol-registry-00780}} & P10.7 & Natural transformations from the identity to GF and from FG to the identity; their components satisfy the two triangle identities.  \\
\(\exists_f,\forall_f\) & powerset quantifiers along an arbitrary function & P10.9 & \hyperref[entry:symbol-registry-00826]{p.~\pageref{entry:symbol-registry-00826}} & P10.9 & Direct image and the condition that the entire fibre belongs to the predicate; empty fibres satisfy the universal condition.  \\
\(\exists^{*}\forall^{*}\) & finite quantifier-block pattern & P2.2 & \hyperref[entry:symbol-registry-00181]{p.~\pageref{entry:symbol-registry-00181}} & P2.2 & A finite existential block followed by a finite universal block; stars are metalinguistic repetition rather than object-language operators.  \\
\(\exists_{\pi_{1,A}}\top_{1\times A}=\top_1\) & internal inhabitance of a sort & P10.9 & \hyperref[entry:symbol-registry-00823]{p.~\pageref{entry:symbol-registry-00823}} & P10.9 & Equivalent by Beck--Chevalley to the corresponding condition in every parameter context; it implies reflection of predicate order, not of base-arrow equality.  \\
\(\exists_\pi,\forall_\pi\) & projection quantifier adjoints & P10.9 & \hyperref[entry:symbol-registry-00818]{p.~\pageref{entry:symbol-registry-00818}} & P10.9 & The maps from \(P(I\times A)\) to \(P(I)\) are respectively left and right adjoints to projection reindexing.  \\
\(\exists r.C,\forall r.C\) & ALC role restrictions & P9.3 & \hyperref[entry:symbol-registry-00640]{p.~\pageref{entry:symbol-registry-00640}} & P9.3 & Qualified existential and universal concept constructors interpreted over the declared role relation.  \\
\(\exists x\,\varphi\) & existential quantifier & global & \hyperref[entry:symbol-registry-00012]{p.~\pageref{entry:symbol-registry-00012}} & P2.2 & Binds occurrences of x in its declared scope; prefix at the unary rank.  \\
\(f\times g\) & product of categorical arrows & P10.7-P10.9 & \hyperref[entry:symbol-registry-00782]{p.~\pageref{entry:symbol-registry-00782}} & P10.7 & The arrow with components f after the first projection and g after the second; respects identity and composition.  \\
\(F\dashv G,\Phi_{A,B}\) & adjunction and hom-set transposition & P10.7-P10.9 & \hyperref[entry:symbol-registry-00779]{p.~\pageref{entry:symbol-registry-00779}} & P10.7 & Phi sends an arrow from F(A) to B to an arrow from A to G(B); its bijections are natural in both arguments.  \\
\(F[a]\) & Boolean filter generated by adjoining one element & P10.2 & \hyperref[entry:symbol-registry-00717]{p.~\pageref{entry:symbol-registry-00717}} & P10.2 & Elements above f meet a for some f in F; used in ultrafilter dichotomy and avoiding-element extension.  \\
\(f:A\rightharpoonup B\) & partial-function type arrow & global & \hyperref[entry:symbol-registry-00226]{p.~\pageref{entry:symbol-registry-00226}} & P1.3 & The function may be undefined outside a declared subset of \(A\).  \\
\(f^B\) & abstract Boolean operator & P10.3 & \hyperref[entry:symbol-registry-00731]{p.~\pageref{entry:symbol-registry-00731}} & P10.3 & For positive arity preserves zero and finite joins in each coordinate; the conditions are vacuous for nullary symbols.  \\
\(f_c\) & matrix operation for connective c & P8.5 & \hyperref[entry:symbol-registry-00113]{p.~\pageref{entry:symbol-registry-00113}} & P8.5 & Total n-ary operation on \(\mathsf V\) interpreting the declared connective \(c\).  \\
\(f={}_Dg,\ f\leq_Dg,\ f<_Dg\) & equality and comparison modulo a filter & P6.5 & \hyperref[entry:symbol-registry-00691]{p.~\pageref{entry:symbol-registry-00691}} & P6.5 & The set of coordinates satisfying the indicated comparison belongs to D; for an ultrafilter the quotient order is linear.  \\
\(f_i(\chi)\) & one-coordinate modal formula context & P10.3 & \hyperref[entry:symbol-registry-00749]{p.~\pageref{entry:symbol-registry-00749}} & P10.3 & All other arguments are fixed; used in theorem-level monotonicity and the normality and additivity axiom schemes.  \\
\(F_L\) & normal modal Lindenbaum algebra & P10.3 & \hyperref[entry:symbol-registry-00742]{p.~\pageref{entry:symbol-registry-00742}} & P10.3 & The formula algebra modulo provable equivalence; arbitrary valuations are handled with finite representative selections per formula.  \\
\(F_{\Lambda,L}\) & Lindenbaum algebra of a finitary normal logic & P10.3 & \hyperref[entry:symbol-registry-00750]{p.~\pageref{entry:symbol-registry-00750}} & P10.3 & The formula quotient by L-provable equivalence with all finite-arity and nullary operators interpreted on classes.  \\
\(F_{\mathcal V}(X),\eta_X\) & free algebra and its generator map & P10.1 & \hyperref[entry:symbol-registry-00709]{p.~\pageref{entry:symbol-registry-00709}} & P10.1 & The map from X has a unique homomorphic extension into every member of the variety; it need not be injective in a trivial variety.  \\
\(f,n,b,t\) & Belnap--Dunn value names & P8.5 & \hyperref[entry:symbol-registry-00550]{p.~\pageref{entry:symbol-registry-00550}} & P8.5 & Omori--Wansing's source letters for false-only neither both and true-only values in the BD algebra.  \\
\(f^*:P(I)\to P(J)\) & predicate reindexing along a context map & P10.9 & \hyperref[entry:symbol-registry-00816]{p.~\pageref{entry:symbol-registry-00816}} & P10.9 & For \(f:J\to I\), the contravariant functor sends predicates in \(I\) to predicates in \(J\), preserving every Heyting operation.  \\
\(F(\vec x)\mathord\downarrow\) & convergence of a partial numerical expression & P4.3-P4.5 & \hyperref[entry:symbol-registry-00248]{p.~\pageref{entry:symbol-registry-00248}} & P4.3 & Metalinguistic definedness predicate.  \\
\(F(\vec x)\mathord\uparrow\) & divergence of a partial numerical expression & P4.3-P4.5 & \hyperref[entry:symbol-registry-00249]{p.~\pageref{entry:symbol-registry-00249}} & P4.3 & Metalinguistic undefinedness predicate.  \\
\(f(w,X),F_w(X)\) & singular and set-valued conditional selection functions & P8.4 & \hyperref[entry:symbol-registry-00592]{p.~\pageref{entry:symbol-registry-00592}} & P8.4 & Stalnaker f selects one member of a nonempty proposition; CK \(F_w\) returns an arbitrary set and need not be contained in X.  \\
\(\forall a.\alpha\) & System-F universal type former & P7.4 & \hyperref[entry:symbol-registry-00417]{p.~\pageref{entry:symbol-registry-00417}} & P7.4 & Binds the type variable a in alpha; its syntactic category distinguishes it from the object-language formula quantifier.  \\
\((\forall^E x)\Phi,(\exists^E x)\Phi\) & Fitting actualist quantifier abbreviations & P8.2 & \hyperref[entry:symbol-registry-00516]{p.~\pageref{entry:symbol-registry-00516}} & P8.2 & Exact source relativizations of the primitive possibilist quantifiers by the existence predicate E.  \\
\(\forall_\sigma(\varphi),\exists_\sigma(\varphi)\) & quantifiers on unary formula abstractions & P10.9 & \hyperref[entry:symbol-registry-00832]{p.~\pageref{entry:symbol-registry-00832}} & P10.9 & Pitts's printed quantifier notation; the body \(\varphi(x)\) is interpreted in the context extended by \(x:\sigma\).  \\
\(\forall x\,\varphi\) & universal quantifier & global & \hyperref[entry:symbol-registry-00011]{p.~\pageref{entry:symbol-registry-00011}} & P2.2 & Binds occurrences of x in its declared scope; prefix at the unary rank.  \\
\(\frac{\alpha:\beta_1,\ldots,\beta_n}{\gamma}\) & Reiter default & P8.7 & \hyperref[entry:symbol-registry-00605]{p.~\pageref{entry:symbol-registry-00605}} & P8.7 & Prerequisite alpha justifications beta-i and conclusion gamma; this fraction is not an ordinary inference rule.  \\
\(G(A)\) & Gale--Stewart payoff game & P6.8 & \hyperref[entry:symbol-registry-00338]{p.~\pageref{entry:symbol-registry-00338}} & P6.8 & Length-omega perfect-information game whose payoff set is \(A\subseteq\omega^\omega\).  \\
\(G(A),A^{\mathrm{Ko}},A^{\mathrm N},A_{\mathrm{Ku}},A^{\mathrm{Ku}}\) & negative-translation maps & P8.1 & \hyperref[entry:symbol-registry-00488]{p.~\pageref{entry:symbol-registry-00488}} & P8.1 & Exact Glivenko Kolmogorov Goedel--Gentzen and inner/full Kuroda maps; superscripts act memberwise on premise contexts when printed there.  \\
\(g[\uparrow x]=\uparrow g(x)\) & Esakia order back condition & P10.3 & \hyperref[entry:symbol-registry-00770]{p.~\pageref{entry:symbol-registry-00770}} & P10.3 & Together with continuity and monotonicity this characterizes Esakia morphisms and ensures implication-preserving inverse images.  \\
\(\gamma\) & Ackermann theoremhood metarule & P8.6 & \hyperref[entry:symbol-registry-00557]{p.~\pageref{entry:symbol-registry-00557}} & P8.6 & From theoremhood of A and of not-A-or-B infer theoremhood of B; admissible for E and R but not a primitive or local-consequence rule.  \\
\(\gamma\) & ordinal metavariable & P7.5 & \hyperref[entry:symbol-registry-00451]{p.~\pageref{entry:symbol-registry-00451}} & P7.5 & Ordinal variable in the Veblen hierarchy and the definition of Gamma-zero; locally disjoint from formula and type metavariables.  \\
\(\Gamma\models\varphi\) & semantic consequence & global & \hyperref[entry:symbol-registry-00028]{p.~\pageref{entry:symbol-registry-00028}} & P2.1 & Every interpretation and assignment in the declared class satisfying \(\Gamma\) satisfies \(\varphi\).  \\
\(\Gamma\Rightarrow\Delta\) & sequent & global & \hyperref[entry:symbol-registry-00075]{p.~\pageref{entry:symbol-registry-00075}} & P2.1 & Metalevel punctuation between antecedent and succedent contexts.  \\
\(\Gamma\vdash A\equiv B\;\mathsf{type},\ \Gamma\vdash a\equiv b:A\) & MLTT judgmental type and term equality & P7.4 & \hyperref[entry:symbol-registry-00425]{p.~\pageref{entry:symbol-registry-00425}} & P7.4 & Typed definitional equality generated by the displayed computations and congruence; equality reflection into identity types is absent.  \\
\(\Gamma\vdash a:A\) & MLTT term-typing judgment & P7.4 & \hyperref[entry:symbol-registry-00424]{p.~\pageref{entry:symbol-registry-00424}} & P7.4 & Dependent typing judgment on sheet:mltt-intensional; the colon is typed ascription rather than a labelled modal formula.  \\
\(\Gamma\vdash A\;\mathsf{type}\) & MLTT type-formation judgment & P7.4 & \hyperref[entry:symbol-registry-00423]{p.~\pageref{entry:symbol-registry-00423}} & P7.4 & Type judgment in the selected intensional dependent calculus; distinct from System-F formation and formula derivability.  \\
\(\Gamma\models_{\mathrm{alg}}s\approx t\) & universal identity consequence & P10.1 & \hyperref[entry:symbol-registry-00706]{p.~\pageref{entry:symbol-registry-00706}} & P10.1 & Every nonempty algebra validating all premise identities under all valuations validates the conclusion identity under all valuations.  \\
\(\Gamma\vdash_{\mathrm{eq}}s\approx t\) & finite equational derivability & P10.1 & \hyperref[entry:symbol-registry-00707]{p.~\pageref{entry:symbol-registry-00707}} & P10.1 & Derivability from identity premises by the complete registered reflexivity symmetry transitivity compatibility and substitution rules.  \\
\(\Gamma\models_{\mathrm H}^{\mathsf{CC}}\varphi\) & matched Henkin consequence & P2.3 & \hyperref[entry:symbol-registry-00193]{p.~\pageref{entry:symbol-registry-00193}} & P2.3 & Truth preservation over exactly the matched comprehension-choice Henkin model class and matching assignments.  \\
\(\Gamma\entails_{\mathsf{IPL}}A\) & IPL Kripke consequence & P8.1 & \hyperref[entry:symbol-registry-00486]{p.~\pageref{entry:symbol-registry-00486}} & P8.1 & Truth preservation at one world in every intuitionistic Kripke structure of the declared model class under all premises in Gamma.  \\
\(\Gamma\derives_{\mathsf{IPL}}A\) & finite-support IPL derivability & P8.1-P8.2 & \hyperref[entry:symbol-registry-00485]{p.~\pageref{entry:symbol-registry-00485}} & P8.1 & A premise-set judgment witnessed by a finite multiset subcontext derivable on sheet:lj; P8.2 reuses it only in the exact IPL-to-S4 embedding theorem.  \\
\(\Gamma\vdash_{S}\varphi\) & derivability (per-system turnstile) & global & \hyperref[entry:symbol-registry-00029]{p.~\pageref{entry:symbol-registry-00029}} & P2.1 & A derivation exists in the named calculus \(S\); the subscript is mandatory outside a declared local abbreviation.  \\
\(\Gamma\models_{S}\varphi\) & semantic consequence (per-system turnstile) & global & \hyperref[entry:symbol-registry-00114]{p.~\pageref{entry:symbol-registry-00114}} & P1.3 & The subscript names the semantics (\(\models_{\mathcal M}\) matrix, \(\models_{\mathrm{K3}}\), \(\models_{\mathrm{LP}}\), \(\models_{\mathsf{S4}}\), ...).  \\
\(\Gamma_0\) & Feferman--Schuette ordinal & P7.5 & \hyperref[entry:symbol-registry-00457]{p.~\pageref{entry:symbol-registry-00457}} & P7.5 & Least positive gamma satisfying Veb-gamma of zero equals gamma under the locally defined Veblen hierarchy.  \\
\(\Gamma_0\) & finite-support witness context & P8.1 & \hyperref[entry:symbol-registry-00487]{p.~\pageref{entry:symbol-registry-00487}} & P8.1 & Finite multiset whose members lie in the ambient IPL premise set; disjoint from the P7.5 Feferman--Schuette ordinal.  \\
\(\Gamma,\Delta\) & Fitting quantified-modal worlds & P8.2 & \hyperref[entry:symbol-registry-00513]{p.~\pageref{entry:symbol-registry-00513}} & P8.2 & Source world metavariables in the displayed Fitting model; their local element-of-G type overrides the default context role.  \\
\(\Gamma,\Delta,\Sigma,\Pi\) & formula contexts & global & \hyperref[entry:symbol-registry-00014]{p.~\pageref{entry:symbol-registry-00014}} & P2.1 & Finite sets, multisets, or sequences exactly as declared by the active calculus.  Source-bound locally typed uses such as the Buchi alphabet Sigma are recorded separately and are not contexts.  \\
\(\gamma,\delta,\xi\) & additional ordinal metavariables & P6 & \hyperref[entry:symbol-registry-00348]{p.~\pageref{entry:symbol-registry-00348}} & P6.1 & Ordinals in recursive, limit, club, forcing, and hierarchy clauses; complements the existing alpha-beta ordinal row.  \\
\(\Gamma\;\mathsf{ctx}\) & MLTT context-formation judgment & P7.4 & \hyperref[entry:symbol-registry-00422]{p.~\pageref{entry:symbol-registry-00422}} & P7.4 & Judgment of sheet:mltt-intensional asserting that the ordered dependent context Gamma is well formed.  \\
\(\Gamma_P,L,G\) & stable operator and well-founded approximation bounds & P9.7 & \hyperref[entry:symbol-registry-00682]{p.~\pageref{entry:symbol-registry-00682}} & P9.7 & Gamma-P maps M to the least model of P-superscript-M; L and G are the least and greatest fixed points of its square.  \\
\(\Gamma,\Phi\) & variable context and hypothesis list in Pitts's calculus & P10.9 & \hyperref[entry:symbol-registry-00827]{p.~\pageref{entry:symbol-registry-00827}} & P10.9 & The first is a finite list of distinct sorted variables; the second is a finite list of formulas in that context.  \\
\(H\) & binary application domain set & P4.2 & \hyperref[entry:symbol-registry-00246]{p.~\pageref{entry:symbol-registry-00246}} & P4.2 & Paired indices and inputs with defined computation.  \\
\(h_*\) & bounded morphism dual to a Boolean-operator homomorphism & P10.3 & \hyperref[entry:symbol-registry-00738]{p.~\pageref{entry:symbol-registry-00738}} & P10.3 & Inverse image on ultrafilters; its back condition follows from the simultaneous filter-extension lemma.  \\
\(h_1\perp h_2,h_1\uplus h_2\) & disjoint finite heaps and their union & P9.6 & \hyperref[entry:symbol-registry-00674]{p.~\pageref{entry:symbol-registry-00674}} & P9.6 & Perpendicular denotes disjoint domains and uplus the defined union of compatible finite partial heaps.  \\
\(h(A)\) & Hartogs number of a set & P6.2 & \hyperref[entry:symbol-registry-00385]{p.~\pageref{entry:symbol-registry-00385}} & P6.2 & Least ordinal admitting no injection into A; its existence is a ZF theorem.  \\
\(H_a\) & effective jump-hierarchy set & P4.5 & \hyperref[entry:symbol-registry-00275]{p.~\pageref{entry:symbol-registry-00275}} & P4.5 & Indexed by a Kleene-O notation.  \\
\(H:K\to X_B\) & explicit Cantor-to-Stone homeomorphism & P10.2 & \hyperref[entry:symbol-registry-00728]{p.~\pageref{entry:symbol-registry-00728}} & P10.2 & Using a Boolean isomorphism alpha from B to the cylinder algebra sends x to the elements b for which x belongs to alpha(b).  \\
\(H_L\) & Lindenbaum Heyting algebra of an intermediate logic & P10.3 & \hyperref[entry:symbol-registry-00772]{p.~\pageref{entry:symbol-registry-00772}} & P10.3 & The quotient by L-provable biconditional, with formula implication inducing the algebraic residual.  \\
\(H_{\mathsf{IPC}},[\varphi]\) & IPC Lindenbaum Heyting algebra & P10.3 & \hyperref[entry:symbol-registry-00766]{p.~\pageref{entry:symbol-registry-00766}} & P10.3 & Formula classes under mutual IPL derivability, ordered by one-premise derivability.  \\
\(H^*,\widehat H\) & restriction of models and translation of named terms & P10.7 & \hyperref[entry:symbol-registry-00810]{p.~\pageref{entry:symbol-registry-00810}} & P10.7 & A theory morphism induces precomposition on models; its named-symbol translation commutes with the recovery isomorphisms.  \\
\(H_{x,\varphi,c}\) & Henkin witness implication & P2.2 & \hyperref[entry:symbol-registry-00174]{p.~\pageref{entry:symbol-registry-00174}} & P2.2 & Metalinguistic abbreviation for \(\exists x\,\varphi\to\varphi[c/x]\) with fresh constant \(c\).  \\
\(I\) & idempotent exponential phase set & P8.6 & \hyperref[entry:symbol-registry-00573]{p.~\pageref{entry:symbol-registry-00573}} & P8.6 & Idempotent members of the tensor-unit fact used in the displayed full-CLL exponential clauses.  \\
\(I\) & nonempty agent-index set & P8.3 & \hyperref[entry:symbol-registry-00525]{p.~\pageref{entry:symbol-registry-00525}} & P8.3 & Fixed set indexing the epistemic or doxastic modalities and relations; disjoint from the P8.6 idempotent exponential phase set.  \\
\(@_i\Diamond j\) & hybrid tableau accessibility atom & P8.3 & \hyperref[entry:symbol-registry-00583]{p.~\pageref{entry:symbol-registry-00583}} & P8.3 & Branch notation asserting that the world named by j is accessible from the world named by i.  \\
\(I=(\Delta^I,\mathord{\cdot}^I)\) & description-logic interpretation datum & P9.3 & \hyperref[entry:symbol-registry-00637]{p.~\pageref{entry:symbol-registry-00637}} & P9.3 & Delta-I is the nonempty individual domain and dot-I interprets concept role and individual names with the types fixed in Definition def:p09-alc.  \\
\(I_\Gamma\) & interpreted finite variable context & P10.9 & \hyperref[entry:symbol-registry-00834]{p.~\pageref{entry:symbol-registry-00834}} & P10.9 & A chosen product starting with terminal object \(1\), extended by one sort object for each declaration.  \\
\(i,j\) & hybrid-logic nominals & P8.3 & \hyperref[entry:symbol-registry-00581]{p.~\pageref{entry:symbol-registry-00581}} & P8.3 & Object-language atoms constrained to name exactly one world each; disjoint from metalinguistic world labels.  \\
\(I_{\mathcal C}\) & extensional index set & P4.4 & \hyperref[entry:symbol-registry-00270]{p.~\pageref{entry:symbol-registry-00270}} & P4.4 & Indices \(e\) such that \(W_e\) belongs to the declared class.  \\
\(@_iA\) & hybrid satisfaction operator & P8.3 & \hyperref[entry:symbol-registry-00582]{p.~\pageref{entry:symbol-registry-00582}} & P8.3 & Evaluates A at the unique world named by nominal i; prefix unary and internal to the object language.  \\
\(\in_X\hookrightarrow P(X)\times X\) & membership subobject of a power object & P10.9 & \hyperref[entry:symbol-registry-00853]{p.~\pageref{entry:symbol-registry-00853}} & P10.9 & The pullback of truth along evaluation; under stage classification membership tests whether \((1_c,x)\) lies in the represented subpresheaf.  \\
\(\iota_n\) & at-least-n sentence & P2.2 & \hyperref[entry:symbol-registry-00159]{p.~\pageref{entry:symbol-registry-00159}} & P2.2 & The first-order sentence asserting the existence of \(n\) pairwise distinct elements.  \\
\(\iota,o\) & individual and truth-value base types & P2.3 & \hyperref[entry:symbol-registry-00201]{p.~\pageref{entry:symbol-registry-00201}} & P2.3 & Base types of the representative simple type theory; iota here is a type and not the P2.4 description binder.  \\
\(\iota x\,\varphi\) & definite-description term former & P2.4-P8.4 & \hyperref[entry:symbol-registry-00219]{p.~\pageref{entry:symbol-registry-00219}} & P2.4 & Binds x in the description formula; a local Russellian elimination or primitive partial-term policy fixes its standing.  \\
\(J_C(d,a,b,p)\) & MLTT identity eliminator & P7.4 & \hyperref[entry:symbol-registry-00429]{p.~\pageref{entry:symbol-registry-00429}} & P7.4 & Typed eliminator for the displayed identity motive C with its judgmental refl computation rule.  \\
\(j_{\Gamma,\Delta},\chi_{\Gamma,\Delta}\) & product and exponential comparison isomorphisms & P10.8 & \hyperref[entry:symbol-registry-00796]{p.~\pageref{entry:symbol-registry-00796}} & P10.8 & The product comparison uses interpreted block projections; the exponential comparison transposes interpreted evaluation after the inverse product comparison.  \\
\(j:V\to M\) & elementary class embedding & P6.6-P6.8 & \hyperref[entry:symbol-registry-00375]{p.~\pageref{entry:symbol-registry-00375}} & P6.6 & Metalinguistic embedding read formula by formula; neither V nor M is added as a class object to ZFC.  \\
\(K\) & self-application domain set & P4.2-P4.4 & \hyperref[entry:symbol-registry-00245]{p.~\pageref{entry:symbol-registry-00245}} & P4.2 & Indices \(e\) for which \(\varphi_e(e)\) is defined.  \\
\(K=2^{\mathbb N}\) & Cantor space of binary sequences & P10.2 & \hyperref[entry:symbol-registry-00725]{p.~\pageref{entry:symbol-registry-00725}} & P10.2 & The countable product of the discrete two-element space with coordinates starting at zero.  \\
\(K*A,K+A,K-A\) & AGM revision expansion and contraction & P8.7 & \hyperref[entry:symbol-registry-00603]{p.~\pageref{entry:symbol-registry-00603}} & P8.7 & Operations on the deductively closed belief set K; their standing is fixed by the AGM postulates and Levi and Harper identities.  \\
\(K_i\) & knowledge operator (agent i) & P8.3 & \hyperref[entry:symbol-registry-00098]{p.~\pageref{entry:symbol-registry-00098}} & P8.3 & Prefix unary at the negation rank.  \\
\(K,KD,KT,KTB,KB\) & Chellas normal modal logic identifiers & P8.2 & \hyperref[entry:symbol-registry-00494]{p.~\pageref{entry:symbol-registry-00494}} & P8.2 & Chellas writes KD KT and KTB for the serial reflexive and reflexive-symmetric extensions and writes KB for the symmetry-only extension.  The familiar D T and B system names occur only as Hughes--Cresswell aliases on the corresponding cards.  \\
\((k,s,t),E\) & equations with explicit finite contexts & P10.7 & \hyperref[entry:symbol-registry-00799]{p.~\pageref{entry:symbol-registry-00799}} & P10.7 & Each triple asserts equality of its two terms under every assignment to the declared \(k\)-variable context.  \\
\(K_V(x)\) & prefix-free Kolmogorov complexity & P4.5 & \hyperref[entry:symbol-registry-00285]{p.~\pageref{entry:symbol-registry-00285}} & P4.5 & Relative to prefix-free decompressor \(V\).  \\
\(K4,KT4,S4,KD45\) & Chellas normal modal logic identifiers & P8.2 & \hyperref[entry:symbol-registry-00495]{p.~\pageref{entry:symbol-registry-00495}} & P8.2 & Chellas's exact plain math-italic names; S4 is the source-named system KT4.  \\
\(\kappa^+\) & successor-cardinal operation & P6.3-P6.8 & \hyperref[entry:symbol-registry-00349]{p.~\pageref{entry:symbol-registry-00349}} & P6.3 & Least cardinal strictly above kappa; distinct from ordinal successor alpha+1.  \\
\(\kappa\longrightarrow(\lambda)^\ell_\mu\) & partition relation & P6.4-P6.6 & \hyperref[entry:symbol-registry-00365]{p.~\pageref{entry:symbol-registry-00365}} & P6.4 & Every mu-coloring of the ell-element subsets of kappa has a homogeneous subset of cardinality lambda; not metalinguistic implication.  \\
\(\kappa^\lambda,\lambda^{<\kappa}\) & cardinal exponentiation and bounded-exponent supremum & P6.3-P6.6 & \hyperref[entry:symbol-registry-00350]{p.~\pageref{entry:symbol-registry-00350}} & P6.3 & Cardinal function-space size and the supremum over exponents below kappa; the full entry is Chapter 6.5.  \\
\(\kappa,\lambda,\mu\) & ordinal or cardinal metavariables as explicitly quantified & P6 & \hyperref[entry:symbol-registry-00347]{p.~\pageref{entry:symbol-registry-00347}} & P6.1 & The local quantifier fixes the standing: limit ordinals and conditional hierarchy markers occur before the cardinal-arithmetic chapters.  \\
\(\kappa_T,\iota_T\) & recursive type comparison terms & P10.8 & \hyperref[entry:symbol-registry-00793]{p.~\pageref{entry:symbol-registry-00793}} & P10.8 & Closed inverse comparisons between T and the named type for its interpreted object; applications preserve denotation.  \\
\(\ker h\) & kernel congruence of an algebra homomorphism & P10.1 & \hyperref[entry:symbol-registry-00713]{p.~\pageref{entry:symbol-registry-00713}} & P10.1 & Pairs of domain elements with equal h-images; the kernel of a product map is the intersection of its coordinate kernels.  \\
\(L\) & constructible universe & P6.3-P6.6 & \hyperref[entry:symbol-registry-00387]{p.~\pageref{entry:symbol-registry-00387}} & P6.3 & Proper class union of the constructible levels; distinct from the global signature macro and from a set-sized level.  \\
\(L(A)\) & CTL labelling truth set & P9.5 & \hyperref[entry:symbol-registry-00653]{p.~\pageref{entry:symbol-registry-00653}} & P9.5 & The set of states assigned to subformula A by bottom-up CTL labelling; locally typed apart from languages and well-founded bounds.  \\
\(L_\alpha\) & constructible level & P6.3 & \hyperref[entry:symbol-registry-00067]{p.~\pageref{entry:symbol-registry-00067}} & P6.3 & The \(\alpha\)-th level of the constructible hierarchy.  \\
\(L_e(\alpha)\) & left adjoint for parameterized equality & P10.9 & \hyperref[entry:symbol-registry-00821]{p.~\pageref{entry:symbol-registry-00821}} & P10.9 & Defined by \(p^*\alpha\wedge r^*\delta_A\), with \(L_e\dashv e^*\); its full order adjunction gives equality elimination.  \\
\(L_{\mu}\) & Kozen modal fixed-point language and calculus identifier & P8.3 & \hyperref[entry:symbol-registry-00520]{p.~\pageref{entry:symbol-registry-00520}} & P8.3 & The exact italic identifier used by Kozen for the finitary propositional modal mu-calculus; the same identifier subscripts its equational derivability relation.  \\
\(L_{\omega_1\omega}\) & countably infinitary first-order language & P2.4 & \hyperref[entry:symbol-registry-00220]{p.~\pageref{entry:symbol-registry-00220}} & P2.4 & Allows countable Boolean combinations but only finite quantifier strings.  \\
\(L,Q;b,b_0\) & Reynolds auxiliary atom data & P8.3 & \hyperref[entry:symbol-registry-00533]{p.~\pageref{entry:symbol-registry-00533}} & P8.3 & In the exact CTL-star AA side condition \(L\) and \(Q\) are disjoint atom sets and \(b\) and \(b_0\) are auxiliary atoms; \(Q\) is locally typed as an atom set, unlike the Buchi state-set role.  \\
\(L(Y)\) & schematic extension of a modal logic to another variable set & P10.3 & \hyperref[entry:symbol-registry-00747]{p.~\pageref{entry:symbol-registry-00747}} & P10.3 & The least normal logic over Y containing all substituted instances of the original logic; used to state all-variable-set canonicity.  \\
\(LA,\operatorname{Th}\) & autoepistemic belief and classical closure & P8.7 & \hyperref[entry:symbol-registry-00608]{p.~\pageref{entry:symbol-registry-00608}} & P8.7 & L is the belief operator; Th is ordinary classical propositional closure over the autoepistemic language and adds no modal axiom.  \\
\(\lambda\omega,\lambda\mathrm P\omega=\lambda\mathrm C\) & Barendregt lambda-cube calculus identifiers & P12.5 & \hyperref[entry:symbol-registry-00627]{p.~\pageref{entry:symbol-registry-00627}} & P12.5 & Exact source identifiers for the terms-on-types plus types-on-types vertex and the all-three-axes vertex; lambda P omega equals lambda C.  \\
\(\lambda\underline{\omega},\lambda\mathrm P\underline{\omega}\) & Barendregt lambda-cube calculus identifiers & P12.5 & \hyperref[entry:symbol-registry-00626]{p.~\pageref{entry:symbol-registry-00626}} & P12.5 & Exact source identifiers for the types-on-types vertex and its types-on-terms extension.  \\
\(\lambda 2,\lambda\mathrm P2\) & Barendregt lambda-cube calculus identifiers & P12.5 & \hyperref[entry:symbol-registry-00625]{p.~\pageref{entry:symbol-registry-00625}} & P12.5 & Exact source identifiers for the terms-on-types vertex and its types-on-terms extension.  \\
\(\Lambda a.M\) & System-F type abstraction & P7.4 & \hyperref[entry:symbol-registry-00418]{p.~\pageref{entry:symbol-registry-00418}} & P7.4 & Term former binding type variable a in M; distinct from ordinary term abstraction.  \\
\(\lambda f.\lambda x.f^n x\) & Church numeral encoding of n & P7.4 & \hyperref[entry:symbol-registry-00412]{p.~\pageref{entry:symbol-registry-00412}} & P7.4 & Closed untyped term used as the fixed numerical input-output encoding for lambda definability.  \\
\({}^{\lambda}M\subseteq M\) & closure under ambient lambda-sequences & P6.6-P6.8 & \hyperref[entry:symbol-registry-00376]{p.~\pageref{entry:symbol-registry-00376}} & P6.6 & Every ambient function from lambda into M belongs to M; stronger than transitivity.  \\
\(\Lambda,n_f\) & finitary Boolean-operator signature and arity & P10.3 & \hyperref[entry:symbol-registry-00730]{p.~\pageref{entry:symbol-registry-00730}} & P10.3 & A set of additional symbols with one finite arity per symbol; nullary symbols designate arbitrary Boolean elements.  \\
\(\lambda{\to},\lambda\mathrm P\) & Barendregt lambda-cube calculus identifiers & P12.5 & \hyperref[entry:symbol-registry-00624]{p.~\pageref{entry:symbol-registry-00624}} & P12.5 & Exact source identifiers for the base simply typed vertex and the types-on-terms vertex.  \\
\(\lambda x{:}\alpha.B\) & typed lambda abstraction & P2.3-P7.6 & \hyperref[entry:symbol-registry-00205]{p.~\pageref{entry:symbol-registry-00205}} & P2.3 & Object-language term former binding x in the displayed body; used by STLC System F and System T with capture-avoiding substitution.  \\
\(\lambda x.M\) & untyped lambda abstraction & P7.4 & \hyperref[entry:symbol-registry-00407]{p.~\pageref{entry:symbol-registry-00407}} & P7.4 & Term former binding x in M; separate exact-scope role from typed abstraction.  \\
\(\Lang_2,\mathsf{RelVar}_n\) & relational second-order signature and variable stocks & P2.3 & \hyperref[entry:symbol-registry-00184]{p.~\pageref{entry:symbol-registry-00184}} & P2.3 & For every positive n the second component is the countably infinite stock of n-ary relation variables.  \\
\(\Lang_A\) & first-order language of arithmetic & P4.3-P5 & \hyperref[entry:symbol-registry-00247]{p.~\pageref{entry:symbol-registry-00247}} & P4.3 & Signature containing zero successor addition and multiplication.  \\
\(\Lang_{\Gamma},\Lang_{\Gamma}^{*}\) & active tableau signatures & P2.2 & \hyperref[entry:symbol-registry-00163]{p.~\pageref{entry:symbol-registry-00163}} & P2.2 & The finite root subsignature and its countable fresh-parameter expansion used by the fair first-order tableau construction.  \\
\(\Lang^H\) & Henkin constant expansion & P2.2 & \hyperref[entry:symbol-registry-00175]{p.~\pageref{entry:symbol-registry-00175}} & P2.2 & Signature obtained by adjoining the fresh constants used by the Henkin witness construction.  \\
\(\Lang_{\in}\) & first-order membership signature & P6 & \hyperref[entry:symbol-registry-00378]{p.~\pageref{entry:symbol-registry-00378}} & P6.1 & One-sorted signature with equality and membership as its sole nonlogical relation; classes remain metalinguistic.  \\
\(\Lang(\mathcal M),\mathcal M_{\mathcal M}\) & parameter signature and canonical expansion & P2.2-P3 & \hyperref[entry:symbol-registry-00177]{p.~\pageref{entry:symbol-registry-00177}} & P2.2 & Adds one constant for each domain element and interprets every such constant by the element it names.  \\
\(\Lang_{\mathrm{MS}}^{\flat},Q_{\mathsf s}\) & one-sorted reduct signature and sort predicates & P2.4 & \hyperref[entry:symbol-registry-00214]{p.~\pageref{entry:symbol-registry-00214}} & P2.4 & The flat signature forgets type declarations and \(Q_{\mathsf s}\) cuts out the component for sort \(\mathsf s\).  \\
\(\Lang_{\mathrm{MS}},\mathsf{Sort},\operatorname{tp}\) & strict many-sorted signature data & P2.4 & \hyperref[entry:symbol-registry-00212]{p.~\pageref{entry:symbol-registry-00212}} & P2.4 & The nonempty sort set and symbol-typing map determine the strict many-sorted language.  \\
\(\Lang_R,\delta_R\) & definitional-extension signature and axiom & P2.2 & \hyperref[entry:symbol-registry-00160]{p.~\pageref{entry:symbol-registry-00160}} & P2.2 & The signature obtained by adjoining relation symbol \(R\) and its exact defining biconditional axiom.  \\
\(\langle C\rangle\) & clone generated by truth functions & P2.1 & \hyperref[entry:symbol-registry-00137]{p.~\pageref{entry:symbol-registry-00137}} & P2.1 & Intersection of all positive-arity Boolean clones containing C.  \\
\(\langle C,s\rangle\Downarrow t\) & deterministic while-language big-step evaluation & P9.6 & \hyperref[entry:symbol-registry-00671]{p.~\pageref{entry:symbol-registry-00671}} & P9.6 & Least evaluation relation from a command and initial integer state to a terminating final state.  \\
\(\langle W,V\rangle\models_x^y\alpha\) & Davies--Humberstone two-index truth relation & P8.2 & \hyperref[entry:symbol-registry-00509]{p.~\pageref{entry:symbol-registry-00509}} & P8.2 & Exact source judgment: x is the lower evaluation-world index and y is the upper actual-world perspective.  \\
\(\lbrack A\rbrack\) & IMLL Lindenbaum quotient class & P8.6 & \hyperref[entry:symbol-registry-00579]{p.~\pageref{entry:symbol-registry-00579}} & P8.6 & Equivalence class of formula A in the annex sketch's syntactic algebra; distinct from semantic double brackets.  \\
\(\lbrack\!\lbrack A\rbrack\!\rbrack_{\mathcal M}\) & neighborhood truth-set abbreviation & P8.2 & \hyperref[entry:symbol-registry-00506]{p.~\pageref{entry:symbol-registry-00506}} & P8.2 & Set of worlds satisfying A in the declared neighborhood model; distinct from typed-term and P8.6 semantic brackets.  \\
\(\lbrack\!\lbrack A\rbrack\!\rbrack_s^{\mathcal M}\) & typed-term denotation & P2.3 & \hyperref[entry:symbol-registry-00208]{p.~\pageref{entry:symbol-registry-00208}} & P2.3 & Recursive semantic value of a well-typed term under typed assignment s.  \\
\(\leftarrow A_1,\ldots,A_n\) & SLD goal & P9.7 & \hyperref[entry:symbol-registry-00677]{p.~\pageref{entry:symbol-registry-00677}} & P9.7 & Finite ordered goal sequence; an empty sequence is refutation success rather than a clause with an omitted head.  \\
\(\leq\) & canonical ordering & P1-P4.2 & \hyperref[entry:symbol-registry-00079]{p.~\pageref{entry:symbol-registry-00079}} & P1.2 & Ambient order on declared numerical and order-theoretic objects; P4.3 is excluded because bare \(\leq\) there is object-language arithmetic syntax.  \\
\(\leq\) & canonical ordering & P4.4-P4.6 & \hyperref[entry:symbol-registry-00081]{p.~\pageref{entry:symbol-registry-00081}} & P4.4 & Ambient order on degrees, natural numbers, and resource bounds; not arithmetic object-language syntax.  \\
\(\leq\) & canonical ordering & P5.2-P14 & \hyperref[entry:symbol-registry-00083]{p.~\pageref{entry:symbol-registry-00083}} & P6.1 & Ambient order on declared numerical and order-theoretic objects, including ordinals, cardinals, degrees, and frame orders.  \\
\(\leq\) & object-language arithmetic order abbreviation & P4.3 & \hyperref[entry:symbol-registry-00080]{p.~\pageref{entry:symbol-registry-00080}} & P4.3 & Eliminable abbreviation \(r\leq t\defeq\exists z\,(z+r=t)\), with the least variable fresh for the displayed terms; marked bounded quantifiers inherit it.  \\
\(\leq\) & object-language arithmetic order abbreviation & P5.1 & \hyperref[entry:symbol-registry-00082]{p.~\pageref{entry:symbol-registry-00082}} & P5.1 & The P4.3 abbreviation imported for the Q and Rosser chapters.  \\
\(\leq_1\) & one-one reducibility & P4.4 & \hyperref[entry:symbol-registry-00268]{p.~\pageref{entry:symbol-registry-00268}} & P4.4 & The translator is total computable and injective.  \\
\(\leq_m\) & many-one reducibility & P4-P5 & \hyperref[entry:symbol-registry-00050]{p.~\pageref{entry:symbol-registry-00050}} & P4.2 & A total computable translator preserves membership.  \\
\(\leq_T\) & Turing reducibility & P4-P5 & \hyperref[entry:symbol-registry-00051]{p.~\pageref{entry:symbol-registry-00051}} & P4.4 & An oracle machine for the right-hand set decides the left-hand set.  \\
\(\leq_t,\leq_k\) & Belnap--Dunn truth and knowledge orders & P8.5 & \hyperref[entry:symbol-registry-00551]{p.~\pageref{entry:symbol-registry-00551}} & P8.5 & Omori--Wansing's source subscripts for the distinct truth and knowledge orders on BD; neither is a consequence relation.  \\
\(\llbracket\Gamma\rrbracket,\llbracket M\rrbracket\) & categorical interpretation of a context or term & P10.8 & \hyperref[entry:symbol-registry-00788]{p.~\pageref{entry:symbol-registry-00788}} & P10.8 & Context products begin with the chosen terminal object; a term denotes an arrow from its context object to its interpreted type.  \\
\(\llbracket M\rrbracket_\Gamma,\llbracket\varphi\rrbracket_\Gamma\) & contextual term and predicate interpretations & P10.9 & \hyperref[entry:symbol-registry-00835]{p.~\pageref{entry:symbol-registry-00835}} & P10.9 & A term denotes an arrow from \(I_\Gamma\) to its sort object; a formula denotes an element of \(P(I_\Gamma)\).  \\
\(\llbracket t\rrbracket_n\) & categorical evaluation of a named algebraic term & P10.7 & \hyperref[entry:symbol-registry-00808]{p.~\pageref{entry:symbol-registry-00808}} & P10.7 & An arrow \(n\to1\) in the given theory; substitution of terms is composition with their tuple of interpreted arrows.  \\
\(\longrightarrow\) & abstract reduction relation & P7.3 & \hyperref[entry:symbol-registry-00400]{p.~\pageref{entry:symbol-registry-00400}} & P7.3 & Locally declared reduction relation used only in the abstract normalization definitions; not implication or a sequent arrow.  \\
\(\longrightarrow\) & proof-system function arrow & P7.8 & \hyperref[entry:symbol-registry-00483]{p.~\pageref{entry:symbol-registry-00483}} & P7.8 & Metalinguistic total-function arrow in the Cook-Reckhow definition; disjoint from the P7.3 abstract-reduction scope.  \\
\(\longrightarrow_{\mathcal T}\) & one-step machine transition & P4.1-P4.2 & \hyperref[entry:symbol-registry-00233]{p.~\pageref{entry:symbol-registry-00233}} & P4.1 & Relation on configurations of the displayed machine.  \\
\(\longrightarrow_{\mathsf{NJ}}\) & one-step NJ derivation reduction & P7.3 & \hyperref[entry:symbol-registry-00399]{p.~\pageref{entry:symbol-registry-00399}} & P7.3 & Compatible closure of the exact principal reductions and permutation conversions fixed in the normalization entry.  \\
\(M\) & Slagle ground-interpretation metavariable & P12.3 & \hyperref[entry:symbol-registry-00617]{p.~\pageref{entry:symbol-registry-00617}} & P12.3 & The fixed ground interpretation in semantic resolution; exact source role, disjoint from the P7 term-metavariable scope.  \\
\(M\longrightarrow_{\beta}N,\ M\longrightarrow_{\beta}^{*}N\) & one-step beta reduction and finite reflexive-transitive closure & P7.3-P7.4 & \hyperref[entry:symbol-registry-00408]{p.~\pageref{entry:symbol-registry-00408}} & P7.3 & The active sheet fixes term beta reduction; the star denotes finite reflexive-transitive closure.  \\
\(M\mathrel{\mathsf{Par}}N\) & parallel beta-reduction relation & P7.4 & \hyperref[entry:symbol-registry-00410]{p.~\pageref{entry:symbol-registry-00410}} & P7.4 & Auxiliary relation on alpha-equivalence classes fixed by the four local clauses; its closure is compared with ordinary beta reduction.  \\
\(M\longrightarrow_{\mathsf{SK}}N\) & one-step S-K contraction & P7.4 & \hyperref[entry:symbol-registry-00414]{p.~\pageref{entry:symbol-registry-00414}} & P7.4 & Compatible contextual contraction fixed by sheet:sk-combinatory.  \\
\(M\longrightarrow_{\mathsf T_{\mathrm{tm}}}N\) & one-step System-T term reduction & P7.6 & \hyperref[entry:symbol-registry-00466]{p.~\pageref{entry:symbol-registry-00466}} & P7.6 & Compatible reduction generated by beta contraction and the two registered recursor computation rules.  \\
\(M_A,M_h\) & evaluation model and homomorphism transformation & P10.7 & \hyperref[entry:symbol-registry-00803]{p.~\pageref{entry:symbol-registry-00803}} & P10.7 & The model assigns \(A^n\) to \(n\) and evaluates term tuples; a homomorphism \(h\) gives components \(h^n\).  \\
\(M{:}\alpha\) & type ascription & P12.5 & \hyperref[entry:symbol-registry-00631]{p.~\pageref{entry:symbol-registry-00631}} & P12.5 & The local STLC example fixes the term and type categories; distinct from labelled formulas.  \\
\(M[\alpha],\alpha[\beta/a],M[\alpha/a]\) & System-F type application and capture-avoiding type substitution & P7.4 & \hyperref[entry:symbol-registry-00419]{p.~\pageref{entry:symbol-registry-00419}} & P7.4 & The active syntactic category distinguishes term-level type application from substitution into a type or term.  \\
\(M=_{\beta}N\) & beta conversion & P7.4 & \hyperref[entry:symbol-registry-00409]{p.~\pageref{entry:symbol-registry-00409}} & P7.4 & Least congruence equivalence relation containing beta reduction; not elementary equivalence or formula equivalence.  \\
\(M[G]\) & generic extension & P6.3-P6.6 & \hyperref[entry:symbol-registry-00072]{p.~\pageref{entry:symbol-registry-00072}} & P6.3 & The forcing extension obtained by adjoining the \(M\)-generic filter \(G\).  \\
\(M=(I(V),T(V,V'),F(V))\) & Boolean finite transition-system presentation & P9.5 & \hyperref[entry:symbol-registry-00656]{p.~\pageref{entry:symbol-registry-00656}} & P9.5 & I and F denote initial and bad states and total T relates current to next state variables in the McMillan interpolation procedure.  \\
\(M,N,P\) & lambda-term metavariables & P7.3-P7.6 & \hyperref[entry:symbol-registry-00404]{p.~\pageref{entry:symbol-registry-00404}} & P7.3 & Terms of the locally fixed typed or untyped lambda grammar; the active calculus fixes the exact syntax.  \\
\(M[N/x]\) & capture-avoiding term substitution & P7.3-P7.6 & \hyperref[entry:symbol-registry-00406]{p.~\pageref{entry:symbol-registry-00406}} & P7.3 & Substitutes term N for free occurrences of x in M after required alpha-renaming; distinct from formula substitution.  \\
\((M,\perp);mn;1\) & Girard phase-space datum and monoid notation & P8.6 & \hyperref[entry:symbol-registry-00576]{p.~\pageref{entry:symbol-registry-00576}} & P8.6 & Girard's phase space is the pair of a commutative monoid M and perpendicular pole; multiplication is juxtaposition and the monoid unit is 1.  \\
\(M:\sigma[\Gamma],M=N:\sigma[\Gamma]\) & term formation and external term-equation judgments & P10.9 & \hyperref[entry:symbol-registry-00828]{p.~\pageref{entry:symbol-registry-00828}} & P10.9 & Formation is well-sortedness in the retained context; an equation is interpreted as equality of term arrows.  \\
\(M=_\sigma N\) & equality as a first-order formula & P10.9 & \hyperref[entry:symbol-registry-00830]{p.~\pageref{entry:symbol-registry-00830}} & P10.9 & Reindexes \(\delta_{A_\sigma}\) along the pair of term arrows; its truth need not reflect external equality of those arrows.  \\
\(\mathbb B\) & Boolean truth-value set & P2.1 & \hyperref[entry:symbol-registry-00127]{p.~\pageref{entry:symbol-registry-00127}} & P2.1 & The two-element set \(\{0,1\}\); distinct from a Boolean algebra unless one is separately declared.  \\
\(\mathbb B\) & complete Boolean algebra used for Boolean-valued models & P6.3 & \hyperref[entry:symbol-registry-00321]{p.~\pageref{entry:symbol-registry-00321}} & P6.3 & Scope is disjoint from the P2.1 Boolean truth-value set; completeness supplies all joins and meets used by the Boolean-valued recursion.  \\
\(\mathbb F_2\) & two-element field & P2.1 & \hyperref[entry:symbol-registry-00139]{p.~\pageref{entry:symbol-registry-00139}} & P2.1 & Field used for affine Boolean functions and XOR constraints.  \\
\(\mathbb F,j:\mathbb F^{\mathrm{op}}\to T\) & finite ordinals and the structural theory functor & P10.7 & \hyperref[entry:symbol-registry-00797]{p.~\pageref{entry:symbol-registry-00797}} & P10.7 & The identity-on-objects functor preserves finite products; a finite-set function \(m\to n\) selects an arrow \(n\to m\) in the theory.  \\
\(\mathbb N\) & set of natural numbers & global & \hyperref[entry:symbol-registry-00039]{p.~\pageref{entry:symbol-registry-00039}} & P1.3 & Contains 0 under the Part-1 natural-number convention.  \\
\(\mathbb{P}\) & forcing poset & P6.3-P6.4 & \hyperref[entry:symbol-registry-00068]{p.~\pageref{entry:symbol-registry-00068}} & P6.3 & \(q\leq p\) reads \(q\) is stronger, per the fixed forcing convention.  \\
\(\mathbb Q\) & rational ordered field & global & \hyperref[entry:symbol-registry-00183]{p.~\pageref{entry:symbol-registry-00183}} & P5.1 & The ordinary field of rational numbers; used as an order type in P5 and as the scalar domain of linear rational arithmetic in P12.2.  \\
\(\mathbb R\) & real ordered field & global & \hyperref[entry:symbol-registry-00182]{p.~\pageref{entry:symbol-registry-00182}} & P2.3 & The ordinary field of real numbers; the P2.3 categoricity claim names its ordered-field structure.  \\
\(\mathbb Z\) & set of integers & global & \hyperref[entry:symbol-registry-00225]{p.~\pageref{entry:symbol-registry-00225}} & P3.1 & Integers; used for tape positions and integer-coefficient polynomials.  \\
\(\mathbf 1\) & multiplicative unit & P8.6 & \hyperref[entry:symbol-registry-00104]{p.~\pageref{entry:symbol-registry-00104}} & P8.6 & Syntactic tensor unit in IMLL and CLL; the ambient calculus distinguishes it from a semantic monoid identity.  \\
\(\mathbf B_*\) & descriptive general frame of a Boolean algebra with operators & P10.3 & \hyperref[entry:symbol-registry-00736]{p.~\pageref{entry:symbol-registry-00736}} & P10.3 & The canonical relations on the Stone space together with precisely the basic-clopen image as admissible algebra.  \\
\(\mathbf B^\bullet\) & ultrafilter frame of a Boolean algebra with operators & P10.3 & \hyperref[entry:symbol-registry-00734]{p.~\pageref{entry:symbol-registry-00734}} & P10.3 & The set of Boolean ultrafilters with all canonical operator relations.  \\
\(\mathbf B^\sigma\) & canonical extension of a Boolean algebra with operators & P10.3 & \hyperref[entry:symbol-registry-00735]{p.~\pageref{entry:symbol-registry-00735}} & P10.3 & The full complex algebra of the ultrafilter frame; the Boolean embedding is dense and compact.  \\
\(\mathbf E,\mathbf R,\mathbf{RM}\) & relevance-system identifiers & P8.6 & \hyperref[entry:symbol-registry-00555]{p.~\pageref{entry:symbol-registry-00555}} & P8.6 & Bold Roman system typography reproduced from Dunn--Restall's Anderson--Belnap convention; locally typed as system identifiers.  \\
\(\mathbf E_{\to},\mathbf R_{\to}\) & pure implicational relevance-fragment identifiers & P8.6 & \hyperref[entry:symbol-registry-00556]{p.~\pageref{entry:symbol-registry-00556}} & P8.6 & Dunn--Restall's source subscripts for the pure implicational fragments of bold Roman E and R.  \\
\(\mathbf K_3,\mathsf{LP},\mathbf{FDE}\) & three-valued and four-valued logic identifiers & P8.5 & \hyperref[entry:symbol-registry-00542]{p.~\pageref{entry:symbol-registry-00542}} & P8.5 & The inherited presentation controls bold K3 and sans-serif LP; Omori--Wansing controls bold upright FDE.  The weak-Kleene matrix remains unnamed because Kleene's presentation supplies no symbolic system identifier.  \\
\(\mathbf s\) & finite binary-sequence metavariable & P6.2 & \hyperref[entry:symbol-registry-00358]{p.~\pageref{entry:symbol-registry-00358}} & P6.2 & Node of the canonical binary tree; bold form avoids the globally reserved assignment metavariable s.  \\
\(\mathbf t\) & set-theoretic tree-node metavariable & P6.4 & \hyperref[entry:symbol-registry-00368]{p.~\pageref{entry:symbol-registry-00368}} & P6.4 & Node of the locally declared tree; bold form avoids the globally reserved term metavariable t.  \\
\(\mathbf u\sim_G\mathbf v\) & generic-ultrafilter equivalence of Boolean names & P6.3 & \hyperref[entry:symbol-registry-00329]{p.~\pageref{entry:symbol-registry-00329}} & P6.3 & Defined by Boolean equality having value in \(G\); the quotient relation is not ordinary set equality.  \\
\(\mathbf u,\mathbf v,\mathbf w\) & Boolean-valued name metavariables & P6.3 & \hyperref[entry:symbol-registry-00328]{p.~\pageref{entry:symbol-registry-00328}} & P6.3 & Members of \(V^{\mathbb B}\); bold forms avoid the globally reserved valuation and term metavariables.  \\
\(\mathbf0\) & computable Turing degree & P4.4 & \hyperref[entry:symbol-registry-00263]{p.~\pageref{entry:symbol-registry-00263}} & P4.4 & Degree of the empty set.  \\
\(\mathbf0'\) & jump of the computable degree & P4.4 & \hyperref[entry:symbol-registry-00264]{p.~\pageref{entry:symbol-registry-00264}} & P4.4 & Degree of the empty-set jump.  \\
\(\mathcal A\) & admissible proposition family & P8.2 & \hyperref[entry:symbol-registry-00507]{p.~\pageref{entry:symbol-registry-00507}} & P8.2 & Boolean-and-modal-operation closed proposition family in a general frame; disjoint from the P8.3 automaton and P8.6 pomonoid.  \\
\(\mathcal A\) & auxiliary translated structure & P2.3-P2.4 & \hyperref[entry:symbol-registry-00210]{p.~\pageref{entry:symbol-registry-00210}} & P2.3 & A one-sorted or strict many-sorted model in the translation proofs; scope is disjoint from P8 algebraic structures.  \\
\(\mathcal A=(A,\leq,\otimes,1,\multimap)\) & residuated pomonoid datum & P8.6 & \hyperref[entry:symbol-registry-00119]{p.~\pageref{entry:symbol-registry-00119}} & P8.6 & The inherited source tuple.  Within its definition A is locally typed as the algebra carrier, overriding the default formula-metavariable role.  \\
\(\mathcal A_{\mathrm{in}},\mathcal A_{\mathrm{tape}}\) & input and tape alphabets & P4.1 & \hyperref[entry:symbol-registry-00015]{p.~\pageref{entry:symbol-registry-00015}} & P4.1 & No formula-context standing in this scope.  \\
\(\mathcal A=(Q,\Sigma,\Delta,F),L(\mathcal A)\) & bottom-up finite-tree automaton and accepted language & P9.4 & \hyperref[entry:symbol-registry-00651]{p.~\pageref{entry:symbol-registry-00651}} & P9.4 & Delta contains ranked bottom-up transitions and acceptance requires a run whose root state lies in F; distinct from infinite-word and Rabin infinite-tree tuples.  \\
\(\mathcal A=(Q,\Sigma,\delta,q_0,F)\) & Buchi automaton datum & P8.3-P9.4 & \hyperref[entry:symbol-registry-00540]{p.~\pageref{entry:symbol-registry-00540}} & P8.3 & Finite state set alphabet set-valued transition function initial state and accepting set; a deterministic specialization has singleton transition sets.  \\
\(\mathcal{ALC}\) & description-logic language and system identifier & P9.3 & \hyperref[entry:symbol-registry-00636]{p.~\pageref{entry:symbol-registry-00636}} & P9.3 & The exact Boolean concept language with qualified role restrictions from Definition def:p09-alc; extensions in the description-logic table have different identifiers.  \\
\(\mathcal B_n,\mathcal B_{\rm cyl}\) & finite cylinder algebras and their union & P10.2 & \hyperref[entry:symbol-registry-00727]{p.~\pageref{entry:symbol-registry-00727}} & P10.2 & B-n contains all unions of length-n cylinders; the union is the clopen algebra of Cantor space and the direct limit of the inclusions.  \\
\(\mathcal C\) & formula class for transfinite induction & P7.5 & \hyperref[entry:symbol-registry-00454]{p.~\pageref{entry:symbol-registry-00454}} & P7.5 & Locally fixed family of formulas parameterizing the proof-theoretic calibration; disjoint from the P4 class scope.  \\
\(\mathcal C\) & locally declared class of sets or encoded languages & P4.3-P4.6 & \hyperref[entry:symbol-registry-00255]{p.~\pageref{entry:symbol-registry-00255}} & P4.3 & The active class and member type are declared locally.  \\
\(\mathcal C\) & modal frame-class metavariable & P8.2 & \hyperref[entry:symbol-registry-00500]{p.~\pageref{entry:symbol-registry-00500}} & P8.2 & Locally declared class of modal frames; disjoint from the P4 set or language class and the P7 formula-class use.  \\
\(\mathcal C(A,B)\) & categorical hom-set & P10.7-P10.9 & \hyperref[entry:symbol-registry-00775]{p.~\pageref{entry:symbol-registry-00775}} & P10.7 & Arrows from object A to object B; local smallness places each such set in the fixed universe.  \\
\(\mathcal C,m\) & ALC input closure and its cardinality & P9.3 & \hyperref[entry:symbol-registry-00695]{p.~\pageref{entry:symbol-registry-00695}} & P9.3 & Finite set of normalized input concepts their sub-concepts and NNF complements; m is linear in the total input length.  \\
\(\mathcal C^{\mathrm{op}}\) & opposite category & P10.7-P10.9 & \hyperref[entry:symbol-registry-00776]{p.~\pageref{entry:symbol-registry-00776}} & P10.7 & Reverses the arrows and composition order; a contravariant functor is an ordinary functor out of this category.  \\
\(\mathcal C=(S,R,L),R(s,t),E(s),P(s,t),\pi[i],\pi@t,\tau_i\) & finite CTMC rates embedded jumps and timed paths & P9.5 & \hyperref[entry:symbol-registry-00663]{p.~\pageref{entry:symbol-registry-00663}} & P9.5 & The nonnegative rational rate map is typed on all state pairs and may contain positive self-loops; E sums the full row P is R divided by E at nonabsorbing states and tau-i is the residence time before path position i-plus-one. Pi-at-t uses the locally stated endpoint convention.  \\
\(\mathcal D(\Gamma),\mathcal I(P)(\Gamma),\mathcal I(E)(\Gamma)\) & Fitting local domain and predicate extensions & P8.2 & \hyperref[entry:symbol-registry-00515]{p.~\pageref{entry:symbol-registry-00515}} & P8.2 & The source world-domain and world-relative interpretation notation; the special existence predicate satisfies I(E)(Gamma)=D(Gamma).  \\
\(\mathcal D_L,W_L\) & descriptive canonical frame of a finitary normal logic & P10.3 & \hyperref[entry:symbol-registry-00753]{p.~\pageref{entry:symbol-registry-00753}} & P10.3 & Worlds are ultrafilters of F-Lambda-L; the relations are the full canonical operator relations and the admissible sets are Stone clopens.  \\
\(\mathcal D=(S,P,L),\mathcal M=(S,\mathsf{Act},P,L)\) & finite DTMC and MDP data & P9.5 & \hyperref[entry:symbol-registry-00658]{p.~\pageref{entry:symbol-registry-00658}} & P9.5 & P is row-stochastic for the chain and action-indexed for the MDP; only a scheduler-resolved MDP induces one path measure.  \\
\(\mathcal F_{\mathbb T},\eta_X\) & free Lawvere-model functor and generator map & P10.7 & \hyperref[entry:symbol-registry-00813]{p.~\pageref{entry:symbol-registry-00813}} & P10.7 & The free model has values \(Q_X^n\) after the fixed theory isomorphism; the generator maps \(x\) to the singleton tuple \(([x])\).  \\
\(\mathcal F,\mathcal R\) & function- and relation-symbol sets & global & \hyperref[entry:symbol-registry-00150]{p.~\pageref{entry:symbol-registry-00150}} & P2.2 & The disjoint nonlogical-symbol families of a first-order signature; constants are zero-ary function symbols.  \\
\(\mathcal G^+,\mathcal A\) & admissible algebra of a general frame & P10.3 & \hyperref[entry:symbol-registry-00737]{p.~\pageref{entry:symbol-registry-00737}} & P10.3 & The carrier is the Boolean family of admissible sets closed under complex operators.  \\
\(\mathcal G,\mathfrak F\) & permutation group and normal subgroup filter & P6.3 & \hyperref[entry:symbol-registry-00330]{p.~\pageref{entry:symbol-registry-00330}} & P6.3 & Typed components of a permutation-model datum; \(\mathfrak F\) contains subgroups of \(\mathcal G\).  \\
\(\mathcal H=(H,\wedge,\vee,0,1,\Rightarrow)\) & Heyting algebra datum & P8.1 & \hyperref[entry:symbol-registry-00489]{p.~\pageref{entry:symbol-registry-00489}} & P8.1 & The cited presentation's bounded distributive lattice with residual.  Local element operands type these glyphs as algebraic operations rather than object-language connectives or a sequent arrow.  \\
\(\mathcal H,(\mathcal R_n)_{n\geq1}\) & relational Henkin structure and relation domains & P2.3 & \hyperref[entry:symbol-registry-00190]{p.~\pageref{entry:symbol-registry-00190}} & P2.3 & The object domain is M and each nonempty relation domain is a selected subset of the corresponding powerset.  \\
\(\mathcal I\) & Kozen-model interpretation & P8.3 & \hyperref[entry:symbol-registry-00531]{p.~\pageref{entry:symbol-registry-00531}} & P8.3 & Maps propositional constants to state sets and program constants to binary relations in the exact L-mu semantics; disjoint from the P5.4 relative-interpretation datum.  \\
\(\mathcal I(c)(\Gamma)\) & world-relative interpretation of a quantified-modal constant & P8.2 & \hyperref[entry:symbol-registry-00584]{p.~\pageref{entry:symbol-registry-00584}} & P8.2 & The value of constant c at Fitting world Gamma; rigidity means equality of these values at every pair of worlds and is independent of inner-domain existence.  \\
\(\mathcal I,\varphi_{\mathrm{dom}},\psi^{\mathcal I}\) & relative-interpretation data and formula translation & P5.4 & \hyperref[entry:symbol-registry-00318]{p.~\pageref{entry:symbol-registry-00318}} & P5.4 & A one-dimensional parameter-free interpretation; the domain formula and recursive translation are fixed at def:p05-relative-interpretation.  \\
\(\mathcal K=(S,R,V),\pi\) & total transition structure and path & P8.3-P9.5 & \hyperref[entry:symbol-registry-00521]{p.~\pageref{entry:symbol-registry-00521}} & P8.3 & K is a total transition structure and pi is an infinite path from the declared state; this path role is disjoint from the P7.8 proof-code metavariable.  \\
\(\mathcal L\) & signature & global & \hyperref[entry:symbol-registry-00002]{p.~\pageref{entry:symbol-registry-00002}} & P2.2 & Nonlogical symbols with declared arities or sorts; governs text after a language declaration.  \\
\(\mathcal L(x),\mathcal L(y)\subseteq\mathcal L(x)\) & ALC completion label and subset-blocking comparison & P9.3 & \hyperref[entry:symbol-registry-00643]{p.~\pageref{entry:symbol-registry-00643}} & P9.3 & Labels are subsets of the finite closure; a nonroot y is directly blocked by ancestor x exactly under the displayed subset comparison, re-evaluated after label growth.  \\
\(\mathcal M\models\sigma\) & truth in a structure & global & \hyperref[entry:symbol-registry-00027]{p.~\pageref{entry:symbol-registry-00027}} & P2.2 & Sentence \(\sigma\) is true in \(\mathcal M\).  \\
\(\mathcal M\equiv\mathcal N\) & elementary equivalence & global & \hyperref[entry:symbol-registry-00042]{p.~\pageref{entry:symbol-registry-00042}} & P2.2 & The structures satisfy the same sentences; this is the sole global role of \(\equiv\).  \\
\(\mathcal M\preccurlyeq\mathcal N\) & elementary substructure & global & \hyperref[entry:symbol-registry-00043]{p.~\pageref{entry:symbol-registry-00043}} & P2.2 & Every formula with parameters from \(\mathcal M\) has the same truth value in \(\mathcal M\) and \(\mathcal N\).  \\
\(\mathcal M\cong\mathcal N\) & isomorphism & global & \hyperref[entry:symbol-registry-00048]{p.~\pageref{entry:symbol-registry-00048}} & P2.2 & Structure-preserving bijection; \(\simeq\) remains separately reserved for Kleene equality.  \\
\(\mathcal M\models T\) & model of a theory & global & \hyperref[entry:symbol-registry-00122]{p.~\pageref{entry:symbol-registry-00122}} & P2.2 & \(\mathcal M\) satisfies every member of \(T\).  \\
\(\mathcal M_B\) & Herbrand structure determined by ground atoms & P2.2 & \hyperref[entry:symbol-registry-00165]{p.~\pageref{entry:symbol-registry-00165}} & P2.2 & Ground-term structure whose positive relation atoms are exactly the members of the declared Herbrand-base subset \(B\).  \\
\(\mathcal M,\Gamma\Vdash_v\Phi\) & Fitting quantified-modal truth judgment & P8.2 & \hyperref[entry:symbol-registry-00518]{p.~\pageref{entry:symbol-registry-00518}} & P8.2 & Exact source-shaped truth judgment at world Gamma under a world-independent valuation v.  \\
\(\mathcal M=\langle\mathcal G,\mathcal R,\mathcal D,\mathcal I\rangle\) & Fitting quantified-modal model datum & P8.2 & \hyperref[entry:symbol-registry-00512]{p.~\pageref{entry:symbol-registry-00512}} & P8.2 & Exact source tuple of worlds accessibility world-domain assignment and world-relative interpretation; the local types govern overlapping glyphs.  \\
\(\mathcal M,\mathcal N\) & structures & global & \hyperref[entry:symbol-registry-00022]{p.~\pageref{entry:symbol-registry-00022}} & P2.2 & Structures for the active signature with nonempty domains unless locally declared otherwise; in exact Part V, \(\mathcal N\) is the standard structure for the fixed arithmetic signature.  \\
\(\mathcal M,s\models\varphi\) & satisfaction & global & \hyperref[entry:symbol-registry-00026]{p.~\pageref{entry:symbol-registry-00026}} & P2.2 & \(\varphi\) is satisfied by \(\mathcal M\) under \(s\).  \\
\(\mathcal M,w\models A\) & satisfaction at a world & P8 & \hyperref[entry:symbol-registry-00117]{p.~\pageref{entry:symbol-registry-00117}} & P8.2 & Satisfaction at world \(w\) of the structure \(\mathcal M\).  \\
\(\mathcal N_{\mathrm{ar}}\) & subscripted alias for the standard first-order arithmetic structure & P2.2-P4 & \hyperref[entry:symbol-registry-00178]{p.~\pageref{entry:symbol-registry-00178}} & P2.2 & Earlier separately authored material uses this explicit subscript; exact inherited Part V writes \(\mathcal N\) for the same structure in the fixed arithmetic signature.  \\
\(<_{\mathcal O}\) & Kleene notation order & P4.5 & \hyperref[entry:symbol-registry-00274]{p.~\pageref{entry:symbol-registry-00274}} & P4.5 & Order belonging to the fixed cited notation system.  \\
\(\mathcal O\) & Kleene system of ordinal notations & P4.5 & \hyperref[entry:symbol-registry-00273]{p.~\pageref{entry:symbol-registry-00273}} & P4.5 & The cited standard system fixed in the method entry.  \\
\(\mathcal O(X)\) & Heyting algebra of open subsets & P10.3 & \hyperref[entry:symbol-registry-00764]{p.~\pageref{entry:symbol-registry-00764}} & P10.3 & Finite meets are intersections and implication is the interior of the union of the first operand's complement with the second.  \\
\(\mathcal P(A)\) & power set & global & \hyperref[entry:symbol-registry-00040]{p.~\pageref{entry:symbol-registry-00040}} & P1.3 & Set of all subsets of the set A; set existence comes from the ambient foundation.  \\
\(\mathcal S_{\beta\eta}\) & category of arrow-only typed contexts & P10.8 & \hyperref[entry:symbol-registry-00784]{p.~\pageref{entry:symbol-registry-00784}} & P10.8 & Objects are finite type lists and arrows are typed term tuples modulo componentwise beta-eta equality.  \\
\(\mathcal S_{\mathcal C},E,J\) & internal context category evaluation and quasi-inverse & P10.8 & \hyperref[entry:symbol-registry-00795]{p.~\pageref{entry:symbol-registry-00795}} & P10.8 & Evaluation sends \(\Gamma\) to \(P_\Gamma\); \(J\) names objects and arrows, with \(EJ(A)=1\times A\) for the chosen product.  \\
\(\mathcal S_w,U_w\) & Lewis sphere system and accessible union & P8.4 & \hyperref[entry:symbol-registry-00593]{p.~\pageref{entry:symbol-registry-00593}} & P8.4 & The inclusion-nested union/intersection-closed sphere family at w and its union; no least antecedent-permitting sphere is assumed.  \\
\(\mathcal T\) & deterministic Turing-machine metavariable & P4.1-P4.2 & \hyperref[entry:symbol-registry-00230]{p.~\pageref{entry:symbol-registry-00230}} & P4.1 & Machine under the fixed single-tape convention.  \\
\(\mathcal T\) & tree metavariable & P6.2-P6.4 & \hyperref[entry:symbol-registry-00357]{p.~\pageref{entry:symbol-registry-00357}} & P6.2 & A binary or set-theoretic tree as locally declared; scope is disjoint from the P4.1-P4.2 Turing-machine role.  \\
\(\mathcal T_\alpha,\operatorname{ht}_{\mathcal T}(\mathbf t)\) & tree level and node height & P6.4 & \hyperref[entry:symbol-registry-00367]{p.~\pageref{entry:symbol-registry-00367}} & P6.4 & Height is the predecessor order type and the alpha-level contains exactly nodes of that height.  \\
\(\mathcal T,\mathcal A\) & ALC TBox and ABox & P9.3 & \hyperref[entry:symbol-registry-00638]{p.~\pageref{entry:symbol-registry-00638}} & P9.3 & Finite sets respectively of concept inclusions and concept or role assertions; script A is locally typed apart from automata.  \\
\(\mathcal U\) & kappa-complete ultrafilter witnessing measurability & P6.6 & \hyperref[entry:symbol-registry-00374]{p.~\pageref{entry:symbol-registry-00374}} & P6.6 & Calligraphic form avoids the globally reserved bare theory metavariable U.  \\
\(\mathcal U\) & universe used for categorical size bounds & P10.7-P10.9 & \hyperref[entry:symbol-registry-00837]{p.~\pageref{entry:symbol-registry-00837}} & P10.7 & May be the ambient collection of all sets; a set-sized relativization has the explicitly stated Grothendieck closures, with existence treated as a hypothesis.  \\
\(\mathcal U_i\) & MLTT Russell universe at level i & P7.4 & \hyperref[entry:symbol-registry-00428]{p.~\pageref{entry:symbol-registry-00428}} & P7.4 & Cumulative predicative universe hierarchy on sheet:mltt-intensional; distinct from the P6 measurable-cardinal ultrafilter role of bare calligraphic U.  \\
\(\mathfrak c\) & cardinality of the continuum & P6.4 & \hyperref[entry:symbol-registry-00354]{p.~\pageref{entry:symbol-registry-00354}} & P6.4 & Defined as two to the aleph-zero; used only to bound the Martin-axiom scheme.  \\
\(\mathfrak F\) & declared frame datum & P8.2-P8.6 & \hyperref[entry:symbol-registry-00558]{p.~\pageref{entry:symbol-registry-00558}} & P8.2 & The locally declared tuple fixes whether this is a modal Kripke frame a broad Routley--Meyer frame or the pointed Routley--Meyer frame used by the cited presentation; disjoint from the P6.3 normal-subgroup-filter use.  \\
\(\mathfrak F^+\) & full complex algebra of a relational frame & P10.3 & \hyperref[entry:symbol-registry-00732]{p.~\pageref{entry:symbol-registry-00732}} & P10.3 & Power-set Boolean algebra with existential successor-tuple operations; the carrier may be the power set of the empty frame.  \\
\(\mathfrak F_L,X_L,R_L\) & underlying canonical relational frame & P10.3 & \hyperref[entry:symbol-registry-00745]{p.~\pageref{entry:symbol-registry-00745}} & P10.3 & Same worlds and relation as the canonical general frame; ordinary valuations may use every subset of X-L.  \\
\(\mathfrak G_L,\mathcal A_L\) & canonical general frame and its admissible sets & P10.3 & \hyperref[entry:symbol-registry-00744]{p.~\pageref{entry:symbol-registry-00744}} & P10.3 & Ultrafilters of the modal Lindenbaum algebra with the canonical relation and Stone clopens as admissible variable values.  \\
\(\mathfrak L\) & logic metavariable in a capture statement & P4.6 & \hyperref[entry:symbol-registry-00311]{p.~\pageref{entry:symbol-registry-00311}} & P4.6 & Distinct from \(\Lang\), which globally denotes a signature.  \\
\(\mathfrak L\) & normal modal logic metavariable & P8.2 & \hyperref[entry:symbol-registry-00499]{p.~\pageref{entry:symbol-registry-00499}} & P8.2 & Locally fixed normal modal logic in derivability canonical-model and frame-class constructions; disjoint from the P4.6 capture-logic use.  \\
\(\mathit{HB}_P,T_P,M_P\) & Herbrand base immediate-consequence map and least model & P9.7 & \hyperref[entry:symbol-registry-00676]{p.~\pageref{entry:symbol-registry-00676}} & P9.7 & HB-P is a set of ground atoms T-P is its monotone consequence operator and M-P is the least fixed point.  \\
\(\mathop{\#}\) & Hessenberg natural sum & P7.5 & \hyperref[entry:symbol-registry-00438]{p.~\pageref{entry:symbol-registry-00438}} & P7.5 & Commutative ordinal sum used only to combine the finitely many Cut-node contributions in o(D).  \\
\(\mathop{\triangle}_{\xi<\kappa}C_\xi\) & diagonal intersection & P6.4 & \hyperref[entry:symbol-registry-00363]{p.~\pageref{entry:symbol-registry-00363}} & P6.4 & Set of \(\alpha<\kappa\) belonging to every earlier \(C_\xi\); distinct from symmetric difference.  \\
\(\mathrel{|\!\sim}\) & defeasible consequence & P8.7 & \hyperref[entry:symbol-registry-00601]{p.~\pageref{entry:symbol-registry-00601}} & P8.7 & A relation between information states or formulas whose conclusions may be defeated by additional premises.  \\
\(\mathrm{BD}\) & Belnap--Dunn four-valued algebra & P8.5 & \hyperref[entry:symbol-registry-00547]{p.~\pageref{entry:symbol-registry-00547}} & P8.5 & Omori--Wansing's upright source identifier for the four-valued algebra, represented here by the exact positive/negative support pairs.  \\
\(\mathrm{BL},G,\Pi\) & basic Goedel and Product fuzzy-logic identifiers & P8.5 & \hyperref[entry:symbol-registry-00546]{p.~\pageref{entry:symbol-registry-00546}} & P8.5 & Hajek's source identifiers: upright Roman BL, plain italic G, and plain Pi; local types distinguish G and Pi from metavariables.  \\
\(\mathrm{Eq0}\) & parameter-free equality bi-rule & P10.9 & \hyperref[entry:symbol-registry-00839]{p.~\pageref{entry:symbol-registry-00839}} & P10.9 & Both directions of Pitts's Remark 5.6.2; the displayed variable contexts have no hidden parameter prefix.  \\
\(\mathrm{FO}^{2},\mathrm{GF},\mathrm{EPR}\) & decidable first-order fragment identifiers & P2.2 & \hyperref[entry:symbol-registry-00179]{p.~\pageref{entry:symbol-registry-00179}} & P2.2 & Respectively the two-variable relational fragment guarded relational fragment and Bernays--Schoenfinkel effectively propositional class; equality and function scopes are row-local.  \\
\(\mathrm{MTL}\) & monoidal t-norm logic identifier & P8.5 & \hyperref[entry:symbol-registry-00545]{p.~\pageref{entry:symbol-registry-00545}} & P8.5 & Esteva--Godo's upright Roman system identifier; the source controls both its letters and typography.  \\
\(\mathrm{NAND}(\varphi,\psi),\mathrm{NOR}(\varphi,\psi)\) & metalinguistic Boolean operator names & P2.1 & \hyperref[entry:symbol-registry-00140]{p.~\pageref{entry:symbol-registry-00140}} & P2.1 & Formula abbreviations for negated conjunction and negated disjunction; not active binary glyphs.  \\
\(\mathrm{ND},\mathrm{ADJ}\) & local abbreviations for two contextual rule presentations & P10.9 & \hyperref[entry:symbol-registry-00838]{p.~\pageref{entry:symbol-registry-00838}} & P10.9 & The same theory and structural/equational basis, with either the natural-deduction logical rules or every direction of Pitts's adjoint bi-rules.  \\
\(\mathrm y c,\mathrm y(u)\) & representable presheaf and the Yoneda map & P10.9 & \hyperref[entry:symbol-registry-00842]{p.~\pageref{entry:symbol-registry-00842}} & P10.9 & The value at \(d\) is \(\mathcal C(d,c)\); an arrow between representing objects acts by postcomposition.  \\
\(\mathsf 1_{\mathrm{ty}},\mathsf 0_{\mathrm{ty}}\) & unit and empty types in the categorical STLC extension & P7.3 & \hyperref[entry:symbol-registry-00432]{p.~\pageref{entry:symbol-registry-00432}} & P7.3 & Typed subscripts distinguish these formers from numerical constants and the computable Turing degree.  \\
\(\mathsf{AC}\) & Axiom of Choice & P6 & \hyperref[entry:symbol-registry-00345]{p.~\pageref{entry:symbol-registry-00345}} & P6.1 & Object-language set-theory axiom; occurrences in metatheoretic dependency tables name that exact principle.  \\
\(\mathsf{ACA}\) & full-induction arithmetical-comprehension theory & P5.5 & \hyperref[entry:symbol-registry-00447]{p.~\pageref{entry:symbol-registry-00447}} & P5.5 & ACA-zero strengthened by induction for every formula of the two-sorted language; used to calibrate the uniform Paris-Harrington principle.  \\
\(\mathsf{ACA}_0\) & arithmetical-comprehension theory & P5.7-P7.5 & \hyperref[entry:symbol-registry-00446]{p.~\pageref{entry:symbol-registry-00446}} & P5.7 & RCA-zero plus comprehension for every arithmetical formula; its first-order consequences are exactly those of PA.  \\
\(\mathsf{Act}(s),\sigma:\operatorname{Paths}_{\mathrm{fin}}(\mathcal M)\to\mathsf{Act},\operatorname{Cyl}(\pi)\) & enabled actions scheduler and finite-path cylinder & P9.5 & \hyperref[entry:symbol-registry-00659]{p.~\pageref{entry:symbol-registry-00659}} & P9.5 & Act-s is nonempty sigma chooses an enabled action from each finite history and Cyl-pi is the set of infinite paths extending pi; PCTL extrema range over exactly this scheduler class.  \\
\(\mathsf{AD},\mathsf{PD}\) & determinacy assertions & P6.8 & \hyperref[entry:symbol-registry-00339]{p.~\pageref{entry:symbol-registry-00339}} & P6.8 & AD quantifies over all payoff sets of reals; PD restricts to boldface projective payoff sets.  \\
\(\mathsf{AP}(A),\ell_A(s)\) & formula alphabet and induced state-label map & P9.4 & \hyperref[entry:symbol-registry-00649]{p.~\pageref{entry:symbol-registry-00649}} & P9.4 & AP(A) is the finite set of proposition letters occurring in A; ell-A maps a state to the subset true there under the predicate-indexed valuation V.  \\
\(\mathsf{ATR}_0\) & arithmetical-transfinite-recursion theory & P5.7-P7.5 & \hyperref[entry:symbol-registry-00448]{p.~\pageref{entry:symbol-registry-00448}} & P5.7 & RCA-zero plus the arithmetical-transfinite-recursion scheme along coded well-orders.  \\
\(\mathsf B_\ell,Z,L_\ell\) & binary countdown concepts in the ALC hardness reduction & P9.3 & \hyperref[entry:symbol-registry-00700]{p.~\pageref{entry:symbol-registry-00700}} & P9.3 & B-ell records one rank bit; Z and L-ell abbreviate conjunctions saying all bits or all lower bits are zero.  \\
\(\mathsf B(w,i,v)\) & bounded Goedel-beta graph formula & P4.3 & \hyperref[entry:symbol-registry-00251]{p.~\pageref{entry:symbol-registry-00251}} & P4.3 & Object-language arithmetic formula.  \\
\(\mathsf{BA},\mathsf{Stone}\) & categories in Stone duality & P10.2 & \hyperref[entry:symbol-registry-00724]{p.~\pageref{entry:symbol-registry-00724}} & P10.2 & Boolean algebras with homomorphisms preserving both bounds; compact Hausdorff spaces with clopen bases and continuous maps; trivial and empty objects are allowed.  \\
\(\mathsf{Base},\mathsf b\) & base-type set and base-type metavariable & P7.3-P7.4 & \hyperref[entry:symbol-registry-00401]{p.~\pageref{entry:symbol-registry-00401}} & P7.3 & Fixed generating set for the arrow-only STLC type grammar and one member of it.  \\
\(\mathsf C_1,\derives_{\mathsf C_1}\) & da Costa C1 calculus and derivability & P8.7 & \hyperref[entry:symbol-registry-00599]{p.~\pageref{entry:symbol-registry-00599}} & P8.7 & The source-fixed propositional calculus on sheet:hilbert-da-costa-c1; no claim is transferred to every member of the da Costa hierarchy.  \\
\(\mathsf{case}(L;x.M;y.N),\mathsf{abort}_{\alpha}(M)\) & sum and empty-type eliminators for the categorical STLC extension & P7.3 & \hyperref[entry:symbol-registry-00435]{p.~\pageref{entry:symbol-registry-00435}} & P7.3 & Typed eliminators used only by the explicitly defined bicartesian extension.  \\
\(\mathsf{CH},\mathsf{GCH}\) & continuum hypotheses & P6.3-P6.6 & \hyperref[entry:symbol-registry-00353]{p.~\pageref{entry:symbol-registry-00353}} & P6.3 & CH is the aleph-one equation at the continuum; GCH is its all-aleph scheme.  \\
\(\mathsf{CIC}\) & Calculus of Inductive Constructions identifier & P12.5 & \hyperref[entry:symbol-registry-00628]{p.~\pageref{entry:symbol-registry-00628}} & P12.5 & Calculus of Constructions extended by the separately described inductive families and eliminators.  \\
\(\mathsf{CK},\mathsf V,\mathsf{VW},\mathsf{VC},\mathsf{C1},\mathsf{C2}\) & conditional-logic identifiers & P8.4 & \hyperref[entry:symbol-registry-00594]{p.~\pageref{entry:symbol-registry-00594}} & P8.4 & Basic conditional logic and the separately stated Lewis and Stalnaker extensions; VC is also called C1 in the cited nomenclature.  \\
\(\mathsf{Comp}_{\chi}^{n},\mathsf{Choice}_{\chi}^{n}\) & matched second-order schema abbreviations & P2.3 & \hyperref[entry:symbol-registry-00192]{p.~\pageref{entry:symbol-registry-00192}} & P2.3 & Metalinguistic names for the universally closed eligible formulas displayed on sheet:sol.  \\
\(\mathsf{Comp}(e,x,s)\) & bounded accepting-history formula & P4.3 & \hyperref[entry:symbol-registry-00253]{p.~\pageref{entry:symbol-registry-00253}} & P4.3 & The bounded formula supplied by the inherited verification theorem; no canonical first formula is asserted.  \\
\(\mathsf{Con}_{\mathsf{Proof}_T}(T)\) & presentation-indexed consistency sentence & P5 & \hyperref[entry:symbol-registry-00064]{p.~\pageref{entry:symbol-registry-00064}} & P5.5 & Exposes the chosen proof predicate; for the fixed inherited predicate it is exactly \(\mathsf{Con}(T)\).  \\
\(\mathsf{Con}(\mathsf{ZFC})^\in\) & membership-language translation of the arithmetical ZFC-consistency sentence & P6.3 & \hyperref[entry:symbol-registry-00355]{p.~\pageref{entry:symbol-registry-00355}} & P6.3 & Obtained by translating the fixed arithmetical no-proof sentence along the canonical interpretation of arithmetic by the finite ordinals; the separate countable-transitive-model hypothesis is stronger and must not be conflated with this sentence.  \\
\(\mathsf{Con}(T)\) & consistency sentence & P5-P7.6 & \hyperref[entry:symbol-registry-00063]{p.~\pageref{entry:symbol-registry-00063}} & P5.2 & \(\neg\mathsf{Prov}_T(\overline{\ulcorner\bottom_A\urcorner})\); P7 names the theory argument explicitly.  \\
\(\mathsf{coNP}\) & complements of NP languages & P2.1-P7.8 & \hyperref[entry:symbol-registry-00146]{p.~\pageref{entry:symbol-registry-00146}} & P2.1 & Decision-language class under the declared encoding; reused in P4.6 and P7.8.  \\
\(\mathsf{CVP}\) & circuit value problem & P4.6 & \hyperref[entry:symbol-registry-00313]{p.~\pageref{entry:symbol-registry-00313}} & P4.6 & Valid acyclic Boolean circuit encodings with output value one.  \\
\(\mathsf D\) & effective calculus metavariable & P2.3 & \hyperref[entry:symbol-registry-00189]{p.~\pageref{entry:symbol-registry-00189}} & P2.3 & A calculus whose finite derivations are computably checkable; local to the non-enumerability theorem.  \\
\(\mathsf{DC}\) & dependent choice for a prescribed start & P10.3 & \hyperref[entry:symbol-registry-00761]{p.~\pageref{entry:symbol-registry-00761}} & P10.3 & A serial relation on a nonempty set admits an infinite path beginning at a prescribed point; used in the arbitrary-poset dense-set lemma.  \\
\(\mathsf{DC},\mathsf{AC}_{\omega},\mathsf{AC}_{\omega}^{\mathrm{fin}},\mathsf{BPI}\) & restricted or weaker choice principles & P6.2-P6.4 & \hyperref[entry:symbol-registry-00346]{p.~\pageref{entry:symbol-registry-00346}} & P6.2 & Dependent Choice, Countable Choice, Countable Choice for finite sets, and the Boolean Prime Ideal Theorem; implication arrows are stated separately.  \\
\(\mathsf{DPLL}\) & Davis--Putnam--Logemann--Loveland method identifier & P12.1 & \hyperref[entry:symbol-registry-00611]{p.~\pageref{entry:symbol-registry-00611}} & P12.1 & Exact recursive finite-clause search method defined before its correctness theorem and trace.  \\
\(\mathsf{DTIME}(f(n))\) & deterministic time-class constructor & P4.6 & \hyperref[entry:symbol-registry-00296]{p.~\pageref{entry:symbol-registry-00296}} & P4.6 & Under the fixed machine convention.  \\
\(\mathsf E,\mathsf A\) & CTL and CTL-star path quantifiers & P8.3-P9.5 & \hyperref[entry:symbol-registry-00526]{p.~\pageref{entry:symbol-registry-00526}} & P8.3 & Source-defined existential and universal path quantifiers; scoped away from P8.2 non-normal E.  \\
\(\mathsf{EC}\) & epsilon-free quantifier-free fragment & P7.7 & \hyperref[entry:symbol-registry-00474]{p.~\pageref{entry:symbol-registry-00474}} & P7.7 & Eliminable metalinguistic abbreviation used in the first epsilon theorem; not a second epsilon system.  \\
\(\mathsf{ECT}_0\) & extended Church-thesis scheme & P7.6 & \hyperref[entry:symbol-registry-00472]{p.~\pageref{entry:symbol-registry-00472}} & P7.6 & Universal closure of the displayed almost-negative-domain scheme; not an unrestricted Church-thesis axiom.  \\
\(\mathsf{Eq}_{\alpha},\forall_{\alpha}\) & typed equality and universal-quantifier constants & P2.3 & \hyperref[entry:symbol-registry-00206]{p.~\pageref{entry:symbol-registry-00206}} & P2.3 & Object-language typed constants; Eq-alpha is deliberately distinct from global metalinguistic alpha-equivalence.  \\
\(\mathsf{EUF}\) & equality with uninterpreted functions theory identifier & P12.2 & \hyperref[entry:symbol-registry-00616]{p.~\pageref{entry:symbol-registry-00616}} & P12.2 & Quantifier-free equality theory whose function symbols have no axioms beyond equality congruence.  \\
\(\mathsf{EX},\mathsf{AX},\mathsf{EF},\mathsf{AF}\) & CTL next and eventuality operators & P8.3-P9.5 & \hyperref[entry:symbol-registry-00527]{p.~\pageref{entry:symbol-registry-00527}} & P8.3 & Source-defined paired existential and universal next and eventuality operators.  \\
\(\mathsf{EXPSPACE}\) & exponential-space language class & P4.6 & \hyperref[entry:symbol-registry-00302]{p.~\pageref{entry:symbol-registry-00302}} & P4.6 & Decision-language class.  \\
\(\mathsf F\) & Church-style System-F calculus identifier & P7.4 & \hyperref[entry:symbol-registry-00416]{p.~\pageref{entry:symbol-registry-00416}} & P7.4 & Exact identifier on sheet:system-f; its scope is disjoint from temporal eventually in P8.3.  \\
\(\mathsf F\) & temporal eventually & P8.3-P9.5 & \hyperref[entry:symbol-registry-00096]{p.~\pageref{entry:symbol-registry-00096}} & P8.3 & Defined \(\mathsf FA:\equiv\top\mathbin{\mathsf U}A\), prefix unary at the negation rank.  \\
\(\mathsf f,\mathsf u,\mathsf t\) & neutral free-logic truth values & P8.4 & \hyperref[entry:symbol-registry-00585]{p.~\pageref{entry:symbol-registry-00585}} & P8.4 & Strong-Kleene false undefined and true values for the partial-term variant; only true is designated.  \\
\(\mathsf{Fin}_{\mathcal C}\) & c.e. set of finite-witness codes & P4.4 & \hyperref[entry:symbol-registry-00271]{p.~\pageref{entry:symbol-registry-00271}} & P4.4 & Rice-Shapiro finite positive witnesses.  \\
\(\mathsf{Fm}_{\mathsf{CPL}}\) & set of CPL formulas & P2.1 & \hyperref[entry:symbol-registry-00126]{p.~\pageref{entry:symbol-registry-00126}} & P2.1 & Least set generated by the declared propositional grammar.  \\
\(\mathsf{FOL}(Q)\) & generalized-quantifier extension placeholder & P2.4 & \hyperref[entry:symbol-registry-00223]{p.~\pageref{entry:symbol-registry-00223}} & P2.4 & Pointer notation only; Part 3 must define the supplied isomorphism-invariant quantifier family and its exact arity.  \\
\(\mathsf G\) & temporal always & P8.3-P9.5 & \hyperref[entry:symbol-registry-00097]{p.~\pageref{entry:symbol-registry-00097}} & P8.3 & Defined \(\mathsf GA:\equiv\neg\mathsf F\neg A\), prefix unary at the negation rank.  \\
\(\mathsf{GB},\mathsf{NBG},\mathsf{MK},\mathsf{KM}\) & set-and-class theory abbreviations & P6.2 & \hyperref[entry:symbol-registry-00342]{p.~\pageref{entry:symbol-registry-00342}} & P6.2 & The comparison table fixes GB as choice-free NBG and treats MK and KM as names for the declared impredicative class theory convention.  \\
\(\mathsf{GL},\mathsf{Grz}\) & normal modal logic identifiers & P8.2 & \hyperref[entry:symbol-registry-00497]{p.~\pageref{entry:symbol-registry-00497}} & P8.2 & Exact named provability and Grzegorczyk systems outside the ordinary modal-cube route.  \\
\(\mathsf{Glob}(\mathcal T)\) & ALC tableau global-GCI abbreviation & P9.3 & \hyperref[entry:symbol-registry-00642]{p.~\pageref{entry:symbol-registry-00642}} & P9.3 & Metalinguistic abbreviation for the NNF internalizations of every GCI in the active TBox; used in body initialization and on sheet:tableaux-alc.  \\
\(\mathsf H_{\mathsf{CC}}\) & matched comprehension-choice Henkin model class & P2.3 & \hyperref[entry:symbol-registry-00191]{p.~\pageref{entry:symbol-registry-00191}} & P2.3 & Models satisfy exactly the eligible comprehension and graph-choice schemas; distinct from arbitrary Henkin structures.  \\
\(\mathsf{HA}\) & Heyting arithmetic & P7.5-P7.6 & \hyperref[entry:symbol-registry-00440]{p.~\pageref{entry:symbol-registry-00440}} & P7.5 & First-order arithmetic with full induction over intuitionistic logic; not a semantic notion of constructive truth.  \\
\(\mathsf{HA},\mathsf{BA}\) & Heyting and Boolean classes in the common residual signature & P10.3 & \hyperref[entry:symbol-registry-00771]{p.~\pageref{entry:symbol-registry-00771}} & P10.3 & Classes of algebras for the intermediate-logic correspondence; HA here is not Heyting arithmetic from Part VII.  \\
\(\mathsf{Heyt}\) & category of Heyting algebras and homomorphisms & P10.9 & \hyperref[entry:symbol-registry-00815]{p.~\pageref{entry:symbol-registry-00815}} & P10.9 & Maps preserve both bounds, finite meets and joins, and Heyting implication; reindexing takes values in this category.  \\
\(\mathsf{HYP}\) & hyperarithmetical subsets of the naturals & P4.5 & \hyperref[entry:symbol-registry-00276]{p.~\pageref{entry:symbol-registry-00276}} & P4.5 & Lightface class.  \\
\(\mathsf I\Sigma_n\) & induction-fragment family & P6-P7.5 & \hyperref[entry:symbol-registry-00442]{p.~\pageref{entry:symbol-registry-00442}} & P6.6 & Robinson arithmetic plus induction exactly for Sigma-zero-n formulas, for n at least one; P6 first uses the n=1 conservativity calibration and P7.5 gives the systematic treatment.  \\
\(\mathsf{IF}\) & codes of ill-founded computable linear orders & P4.5 & \hyperref[entry:symbol-registry-00280]{p.~\pageref{entry:symbol-registry-00280}} & P4.5 & Complement of \(\mathsf{WF}\) within \(\mathsf{LO}\).  \\
\(\mathsf{IMLL},\mathsf{CLL}\) & linear-calculus identifiers & P8.6 & \hyperref[entry:symbol-registry-00562]{p.~\pageref{entry:symbol-registry-00562}} & P8.6 & Named multiplicative intuitionistic core and distinct full one-sided classical calculus.  The inherited four-rule bang extension is unnamed and therefore has no fabricated identifier.  \\
\(\mathsf{IPL},\mathsf{KC},\mathsf{LC}\) & intuitionistic and intermediate logic identifiers & P8.1 & \hyperref[entry:symbol-registry-00484]{p.~\pageref{entry:symbol-registry-00484}} & P8.1 & Exact IPL card and the named KC and Goedel--Dummett LC extensions; IQC is omitted because no IQC identifier is printed.  \\
\(\mathsf{KD45}_{I},\mathsf{KD45}_n\) & multi-agent doxastic calculus identifiers & P8.3 & \hyperref[entry:symbol-registry-00523]{p.~\pageref{entry:symbol-registry-00523}} & P8.3 & Exact agent-indexed KD45 system and its finite-cardinality alias.  \\
\(\mathsf L\) & associative Lambek-calculus identifier & P7.1-P8.6 & \hyperref[entry:symbol-registry-00390]{p.~\pageref{entry:symbol-registry-00390}} & P7.1 & The cited associative ordered nonempty-antecedent Lambek calculus; \(\mathsf L^*\) and \(\mathsf L_1\) are source-named contrast extensions rather than aliases.  \\
\(\mathsf L^*,\mathsf L_1\) & Lambek extension identifiers & P8.6 & \hyperref[entry:symbol-registry-00578]{p.~\pageref{entry:symbol-registry-00578}} & P8.6 & Contrast names for empty-antecedent and unit extensions; neither is the associative unit-free Lambek-L calculus displayed on its card.  \\
\(\mathsf L_n,\mathsf L_{\infty}\) & finite and infinite Lukasiewicz system identifiers & P8.5 & \hyperref[entry:symbol-registry-00597]{p.~\pageref{entry:symbol-registry-00597}} & P8.5 & Editorial system-card identifiers; the source symbols \(\text{\L}_n\) and \(\text{\L}\) name the finite chain and standard infinite logic.  \\
\(\mathsf{LCF}\) & LCF trusted-kernel architecture identifier & P12.5 & \hyperref[entry:symbol-registry-00620]{p.~\pageref{entry:symbol-registry-00620}} & P12.5 & Interactive-prover architecture in which an abstract theorem type can be constructed only by a small trusted kernel.  \\
\(\mathsf{LJ}\) & intuitionistic sequent-calculus identifier & P7.1-P8.1 & \hyperref[entry:symbol-registry-00389]{p.~\pageref{entry:symbol-registry-00389}} & P7.1 & Exact single-succedent calculus fixed by sheet:lj; reused for IPL and equality-free IQC with every turnstile subscripted.  \\
\(\mathsf{LK}_{\mathrm d},\mathsf{LK}_{=}\) & defined-negation and equality sequent-calculus identifiers & P7.1-P7.2 & \hyperref[entry:symbol-registry-00388]{p.~\pageref{entry:symbol-registry-00388}} & P7.1 & Exact registered presentations; neither identifier is an alias for canonical primitive-negation LK.  \\
\(\mathsf{LO}\) & valid computable linear-order codes & P4.5 & \hyperref[entry:symbol-registry-00277]{p.~\pageref{entry:symbol-registry-00277}} & P4.5 & Universe within which \(\mathsf{WF}\) and \(\mathsf{IF}\) are complementary.  \\
\(\mathsf{LTL},\mathsf{CTL},\mathsf{CTL}^{*},\mathsf{PAL}\) & temporal and dynamic calculus identifiers & P8.3 & \hyperref[entry:symbol-registry-00519]{p.~\pageref{entry:symbol-registry-00519}} & P8.3 & Exact registered LTL CTL full CTL-star and public-announcement presentations.  \\
\(\mathsf{MA}_\kappa,\mathsf{MA}\) & Martin axiom instance and scheme & P6.4 & \hyperref[entry:symbol-registry-00372]{p.~\pageref{entry:symbol-registry-00372}} & P6.4 & The instance meets at most kappa dense sets in a ccc poset; the scheme ranges below the continuum.  \\
\(\mathsf{MaxSAT}\) & maximum satisfiability optimization problem & P12.1 & \hyperref[entry:symbol-registry-00614]{p.~\pageref{entry:symbol-registry-00614}} & P12.1 & Optimization problem maximizing the number or declared weight of satisfied clauses under the local variant.  \\
\(\mathsf{MSFOL},\mathsf{SOL}_{\mathrm H}^{\mathsf{CC}}\) & translation-defined calculus identifiers & P2.3 & \hyperref[entry:symbol-registry-00198]{p.~\pageref{entry:symbol-registry-00198}} & P2.3 & Eliminable metalinguistic derivability abbreviations; neither names an independently generated native rule set.  \\
\(\mathsf{ND}_{\mathsf{CPL}},\mathsf{LK}_{\mathsf{CPL}},\mathsf H_{\mathsf{CPL}}\) & CPL calculus identifiers & P2.1 & \hyperref[entry:symbol-registry-00141]{p.~\pageref{entry:symbol-registry-00141}} & P2.1 & Names the registered natural-deduction sequent and Hilbert calculi; derivability still carries the calculus as a turnstile subscript.  \\
\(\mathsf{ND}_{\mathsf{FOL}},\mathsf{LK},\mathsf H_{\mathsf{FOL}}\) & first-order calculus identifiers & P2.2-P8.1 & \hyperref[entry:symbol-registry-00162]{p.~\pageref{entry:symbol-registry-00162}} & P2.2 & Names the registered natural-deduction sequent and equality-free Hilbert presentations; canonical LK recurs in P7 and the P8.1 negative-translation statements; turnstiles retain exact subscripts.  \\
\(\mathsf{NEXPTIME},\mathsf{coNEXPTIME},\mathsf{EXPTIME},2\mathsf{EXPTIME}\) & decision-problem complexity classes & P2.2-P4 & \hyperref[entry:symbol-registry-00180]{p.~\pageref{entry:symbol-registry-00180}} & P2.2 & Every bound applies only to the named decision problem.  \\
\(\mathsf{NF},\mathsf{TTT}\) & New Foundations and Tangled Type Theory abbreviations & P6.2 & \hyperref[entry:symbol-registry-00343]{p.~\pageref{entry:symbol-registry-00343}} & P6.2 & Distinct theories linked only by the cited consistency reduction; neither is presented as a subtheory of ZF.  \\
\(\mathsf{NJ}\) & intuitionistic natural-deduction calculus identifier & P7.1-P7.3 & \hyperref[entry:symbol-registry-00398]{p.~\pageref{entry:symbol-registry-00398}} & P7.1 & Exact calculus fixed by sheet:nj; excludes RAA, primitive negation, equality, and induction.  \\
\(\mathsf{not}\ A\) & logic-programming default negation & P9.7 & \hyperref[entry:symbol-registry-00679]{p.~\pageref{entry:symbol-registry-00679}} & P9.7 & Ground operational negation by finite failure in SLDNF; not classical object-language negation.  \\
\(\mathsf{NP}\) & nondeterministic polynomial-time language class & P2.1-P12.1 & \hyperref[entry:symbol-registry-00145]{p.~\pageref{entry:symbol-registry-00145}} & P2.1 & Equivalently polynomially witnessed languages under the declared encoding; reused in P4.6 P7.8 and the P12.1 proof-certificate comparison.  \\
\(\mathsf{NPSPACE}\) & nondeterministic polynomial-space language class & P4.6 & \hyperref[entry:symbol-registry-00301]{p.~\pageref{entry:symbol-registry-00301}} & P4.6 & Decision-language class.  \\
\(\mathsf{NSPACE}(f(n))\) & nondeterministic space-class constructor & P4.6 & \hyperref[entry:symbol-registry-00299]{p.~\pageref{entry:symbol-registry-00299}} & P4.6 & Counts work-tape cells.  \\
\(\mathsf{NTIME}(f(n))\) & nondeterministic time-class constructor & P4.6 & \hyperref[entry:symbol-registry-00297]{p.~\pageref{entry:symbol-registry-00297}} & P4.6 & Under the fixed machine convention.  \\
\(\#\mathsf P\) & counting complexity class & P12.1 & \hyperref[entry:symbol-registry-00613]{p.~\pageref{entry:symbol-registry-00613}} & P12.1 & Functions counting accepting witnesses of a polynomial-time decidable relation.  \\
\(\mathsf P\) & deterministic polynomial-time language class & P2.1 & \hyperref[entry:symbol-registry-00144]{p.~\pageref{entry:symbol-registry-00144}} & P2.1 & Local propositional-complexity standing.  \\
\(\mathsf P\) & deterministic polynomial-time language class & P4.6 & \hyperref[entry:symbol-registry-00293]{p.~\pageref{entry:symbol-registry-00293}} & P4.6 & Scope disjoint from the P4.3 pairing formula.  \\
\(\mathsf P\) & deterministic polynomial-time language class & P12.1 & \hyperref[entry:symbol-registry-00609]{p.~\pageref{entry:symbol-registry-00609}} & P12.1 & Complexity class reused in the P12 SAT and certificate comparison; disjoint from the P4.3 pairing-graph formula.  \\
\(\mathsf P_{\bowtie q}[\psi],\mathsf P^{\min}_{\bowtie q}[\psi],\mathsf P^{\max}_{\bowtie q}[\psi]\) & probability-threshold modalities & P9.5 & \hyperref[entry:symbol-registry-00661]{p.~\pageref{entry:symbol-registry-00661}} & P9.5 & The undecorated form uses the active DTMC or CTMC path measure; min and max forms take the MDP infimum and supremum over the declared scheduler class.  \\
\(\mathsf P_{\mathrm{pf}},\mathsf Q_{\mathrm{pf}}\) & Cook-Reckhow proof-system metavariables & P7.8 & \hyperref[entry:symbol-registry-00478]{p.~\pageref{entry:symbol-registry-00478}} & P7.8 & Polynomial-time encoded surjections onto TAUT; subscripts prevent collision with complexity and formula letters.  \\
\(\mathsf P,\mathsf H\) & unary past tense operators & P8.3 & \hyperref[entry:symbol-registry-00538]{p.~\pageref{entry:symbol-registry-00538}} & P8.3 & Prior-style sometime-past and always-past operators; their temporal standing is fixed locally.  \\
\(\mathsf P(x,y,z)\) & bounded arithmetic pairing-graph formula & P4.3 & \hyperref[entry:symbol-registry-00250]{p.~\pageref{entry:symbol-registry-00250}} & P4.3 & Scope is disjoint from complexity class \(\mathsf P\).  \\
\(\mathsf{PA}_2\) & full second-order Peano theory & P2.3 & \hyperref[entry:symbol-registry-00188]{p.~\pageref{entry:symbol-registry-00188}} & P2.3 & The finite source theory displayed in Section sec:second-order-logic-03; the expanded signature includes the four displayed recursion axioms.  \\
\(\mathsf{par}_n^m\) & parameter-specialization compiler & P4.2-P4.4 & \hyperref[entry:symbol-registry-00244]{p.~\pageref{entry:symbol-registry-00244}} & P4.2 & Total computable compiler for the fixed numbering.  \\
\(\mathsf{PC}_{\varepsilon}\) & equality-free epsilon-calculus identifier & P7.7 & \hyperref[entry:symbol-registry-00473]{p.~\pageref{entry:symbol-registry-00473}} & P7.7 & Exact Hilbert-FOL extension fixed by sheet:epsilon; primitive quantifiers are retained and equality is excluded.  \\
\(\mathsf{PDL}\) & propositional dynamic logic and calculus identifier & P9.6 & \hyperref[entry:symbol-registry-00666]{p.~\pageref{entry:symbol-registry-00666}} & P9.6 & The regular-program modal system whose exact grammar semantics and Hilbert sheet are registered in Part IX.  \\
\(\mathsf{PH}\) & polynomial hierarchy & P4.6 & \hyperref[entry:symbol-registry-00307]{p.~\pageref{entry:symbol-registry-00307}} & P4.6 & Union of its \(\Sigma^p_k\) levels.  \\
\(\mathsf{PHP}^{n}_{n+1},x_{i,j}\) & pigeonhole-principle CNF family and its variables & P7.8 & \hyperref[entry:symbol-registry-00482]{p.~\pageref{entry:symbol-registry-00482}} & P7.8 & Exact weak PHP encoding with n+1 pigeons and n holes; no clauses forbid one pigeon from occupying two holes.  \\
\(\mathsf{PRA}\) & primitive recursive arithmetic & P6-P7.5 & \hyperref[entry:symbol-registry-00441]{p.~\pageref{entry:symbol-registry-00441}} & P6.6 & Fixed for the master strength map by primitive-recursive function symbols, defining equations, and quantifier-free induction; developed systematically in P7.5.  \\
\(\mathsf{Proof}_T\) & proof predicate & P5 & \hyperref[entry:symbol-registry-00057]{p.~\pageref{entry:symbol-registry-00057}} & P5.1 & The fixed Sigma-zero-one formula \(\mathsf{Proof}_T(p,x)\) representing the primitive-recursive relation that p codes a T-proof with conclusion code x.  \\
\(\mathsf{Prop}\) & default propositional-atom set & P8 & \hyperref[entry:symbol-registry-00124]{p.~\pageref{entry:symbol-registry-00124}} & P8.1 & Used where the inherited or cited presentation declares this atom-set glyph.  A named source presentation may instead declare its own atom notation, which governs locally without house renaming.  \\
\(\mathsf{Prop}\) & set of propositional variables & P2.1 & \hyperref[entry:symbol-registry-00123]{p.~\pageref{entry:symbol-registry-00123}} & P2.1 & Exactly \(\{p_n:n\in\mathbb N\}\) in the fixed countable CPL language.  \\
\(\mathsf{Prov}_{\PA}\) & standard PA provability predicate & P5 & \hyperref[entry:symbol-registry-00059]{p.~\pageref{entry:symbol-registry-00059}} & P5.6 & The instance of the fixed inherited provability predicate used by the exact PA arithmetical realization in Solovay completeness.  \\
\(\mathsf{Prov}_T\) & provability predicate & P5 & \hyperref[entry:symbol-registry-00058]{p.~\pageref{entry:symbol-registry-00058}} & P5.1 & The exact inherited abbreviation \(\exists p\,\mathsf{Proof}_T(p,x)\); formula codes rather than formulas are its arguments.  \\
\(\mathsf{PSPACE}\) & polynomial-space language class & P4.6 & \hyperref[entry:symbol-registry-00300]{p.~\pageref{entry:symbol-registry-00300}} & P4.6 & Decision-language class.  \\
\(\mathsf{PTS}\) & pure type system identifier & P12.5 & \hyperref[entry:symbol-registry-00621]{p.~\pageref{entry:symbol-registry-00621}} & P12.5 & Parameterized typed calculus determined by its sorts axioms and allowed product triples.  \\
\(\mathsf Q_q,\mathsf H_i,\mathsf A_{i,\gamma}\) & configuration concept names in the ALC hardness reduction & P9.3 & \hyperref[entry:symbol-registry-00699]{p.~\pageref{entry:symbol-registry-00699}} & P9.3 & Unary concept names record machine control head position and one tape symbol per cell; their extensions are constrained by finite general inclusions.  \\
\(\mathsf{RCA}_0\) & recursive-comprehension base theory & P5.7-P7.5 & \hyperref[entry:symbol-registry-00444]{p.~\pageref{entry:symbol-registry-00444}} & P5.7 & The exact two-sorted base with Q number axioms, set extensionality, Sigma-zero-one induction, and Delta-zero-one comprehension is fixed in Part V.  \\
\(\mathsf{RProof}_T\) & Rosser proof formula & P5 & \hyperref[entry:symbol-registry-00061]{p.~\pageref{entry:symbol-registry-00061}} & P5.6 & A fixed Sigma-zero-one formula numeralwise representing the primitive-recursive Rosser proof relation.  \\
\(\mathsf{RProv}_T\) & Rosser provability predicate & P5 & \hyperref[entry:symbol-registry-00062]{p.~\pageref{entry:symbol-registry-00062}} & P5.6 & A proof of the coded sentence exists whose code is strictly smaller than every proof code for its negation; distinct from ordinary provability.  \\
\(\mathsf{Run}_{U_K}(p,x,t)\) & bounded prefix-machine computation formula & P4.5 & \hyperref[entry:symbol-registry-00289]{p.~\pageref{entry:symbol-registry-00289}} & P4.5 & Fixed-machine/string analogue of the inherited Comp formula; the finite trace codings are primitively recursively intertranslatable.  \\
\(\mathsf S_{\bowtie q}[\Phi]\) & CSL long-run threshold modality & P9.5 & \hyperref[entry:symbol-registry-00664]{p.~\pageref{entry:symbol-registry-00664}} & P9.5 & Compares the starting-state-dependent limiting probability of Phi-states with q; not the temporal since connective.  \\
\(\mathsf S,\mathsf K,\mathsf{SK}\) & S-K combinators and combinatory-calculus identifier & P7.4 & \hyperref[entry:symbol-registry-00413]{p.~\pageref{entry:symbol-registry-00413}} & P7.4 & Sheet:sk-combinatory fixes the two constants and exact calculus; scoped away from Post-clone and modal-system letters.  \\
\(\mathsf s,\mathsf t,\mathsf s_i;x^{\mathsf s},M_{\mathsf s}\) & sort metavariables sorted variables and carriers & P2.4 & \hyperref[entry:symbol-registry-00213]{p.~\pageref{entry:symbol-registry-00213}} & P2.4 & Sans-serif sort letters avoid the canonical assignment s; strict carriers are nonempty and pairwise disjoint.  \\
\(\mathsf{S5}_{I}^{C},\mathsf{S5}_{n}^{C}\) & common-knowledge calculus identifiers & P8.3 & \hyperref[entry:symbol-registry-00524]{p.~\pageref{entry:symbol-registry-00524}} & P8.3 & Exact common-knowledge extension and its finite-cardinality alias.  \\
\(\mathsf{S5}_{I},\mathsf{S5}_n\) & multi-agent epistemic calculus identifiers & P8.3 & \hyperref[entry:symbol-registry-00522]{p.~\pageref{entry:symbol-registry-00522}} & P8.3 & Exact agent-indexed S5 system and its finite-cardinality alias.  \\
\(\mathsf{SAT}\) & satisfiable propositional encodings & P2.1 & \hyperref[entry:symbol-registry-00147]{p.~\pageref{entry:symbol-registry-00147}} & P2.1 & Malformed strings are excluded by the local encoding convention.  \\
\(\mathsf{SAT}\) & satisfiable propositional encodings & P4.6 & \hyperref[entry:symbol-registry-00294]{p.~\pageref{entry:symbol-registry-00294}} & P4.6 & Cook-Levin target under the declared encoding.  \\
\(\mathsf{SAT}\) & satisfiable propositional encodings & P12.1 & \hyperref[entry:symbol-registry-00610]{p.~\pageref{entry:symbol-registry-00610}} & P12.1 & SAT instances under the exact clause-set and encoding conventions fixed in the automation chapter.  \\
\(\#\mathsf{SAT}\) & satisfying-assignment counting function & P12.1 & \hyperref[entry:symbol-registry-00612]{p.~\pageref{entry:symbol-registry-00612}} & P12.1 & Maps a propositional encoding to its number of satisfying total assignments; not a decision language.  \\
\(\mathsf{SCVP}\) & succinct circuit value problem & P4.6 & \hyperref[entry:symbol-registry-00314]{p.~\pageref{entry:symbol-registry-00314}} & P4.6 & Circuit-described exponentially large circuit value instances.  \\
\(\mathsf{SDL},O A,\mathsf P_OA\) & standard deontic identifier obligation and permission & P8.4 & \hyperref[entry:symbol-registry-00587]{p.~\pageref{entry:symbol-registry-00587}} & P8.4 & SDL is KD under the deontic reading; permission is the eliminable abbreviation \(\mathsf P_OA\defeq\neg O\neg A\).  \\
\(\mathsf{SH}\) & Suslin hypothesis & P6.4 & \hyperref[entry:symbol-registry-00369]{p.~\pageref{entry:symbol-registry-00369}} & P6.4 & Assertion that no Suslin tree exists under the chapter's tree convention.  \\
\(\mathsf{SMT}\) & satisfiability-modulo-theories identifier & P12.2 & \hyperref[entry:symbol-registry-00615]{p.~\pageref{entry:symbol-registry-00615}} & P12.2 & Decision problem for quantifier-free formulas interpreted in a fixed background theory.  \\
\(\mathsf{SPACE}(f(n))\) & deterministic space-class constructor & P4.6 & \hyperref[entry:symbol-registry-00298]{p.~\pageref{entry:symbol-registry-00298}} & P4.6 & Counts work-tape cells.  \\
\(\mathsf{STLC}\) & arrow-only simply typed lambda-calculus identifier & P7.3-P7.4 & \hyperref[entry:symbol-registry-00402]{p.~\pageref{entry:symbol-registry-00402}} & P7.3 & Exact calculus fixed by sheet:stlc and used in every corresponding typing turnstile.  \\
\(\mathsf{STLC}\) & arrow-only simply typed lambda-calculus identifier & P12.5 & \hyperref[entry:symbol-registry-00633]{p.~\pageref{entry:symbol-registry-00633}} & P12.5 & Exact calculus reused from sheet:stlc in the mechanization example.  \\
\(\mathsf{STQBF}\) & succinct true quantified Boolean formula problem & P4.6 & \hyperref[entry:symbol-registry-00315]{p.~\pageref{entry:symbol-registry-00315}} & P4.6 & Circuit-described exponentially long closed QBF instances.  \\
\(\mathsf T_0,\mathsf T_1,\mathsf M,\mathsf L,\mathsf S\) & Post maximal clones & P2.1 & \hyperref[entry:symbol-registry-00138]{p.~\pageref{entry:symbol-registry-00138}} & P2.1 & Respectively 0-preserving, 1-preserving, monotone, affine, and self-dual positive-arity Boolean functions.  \\
\(\mathsf T_k(e,\vec x,y)\) & primitive-recursive accepting-computation-code relation & P4.1 & \hyperref[entry:symbol-registry-00241]{p.~\pageref{entry:symbol-registry-00241}} & P4.1 & Kleene normal-form relation.  \\
\(\mathsf T_{\mathrm{qf}}\) & quantifier-free System-T target theory & P7.6 & \hyperref[entry:symbol-registry-00461]{p.~\pageref{entry:symbol-registry-00461}} & P7.6 & Separately named equality theory used as the Dialectica matrix target; not the term calculus alone.  \\
\(\mathsf T_{\mathrm{tm}}\) & System-T term calculus identifier & P7.6 & \hyperref[entry:symbol-registry-00460]{p.~\pageref{entry:symbol-registry-00460}} & P7.6 & Exact finite-type term calculus fixed by sheet:system-t.  \\
\(\mathsf{TAUT}\) & valid propositional encodings & P2.1-P7.8 & \hyperref[entry:symbol-registry-00148]{p.~\pageref{entry:symbol-registry-00148}} & P2.1 & Malformed strings are excluded by the fixed encoding convention; reused in P4.6 and P7.8.  \\
\(\mathsf{TOT}\) & totality index set & P4.3 & \hyperref[entry:symbol-registry-00254]{p.~\pageref{entry:symbol-registry-00254}} & P4.3 & Indices of programs accepting with valid output on every input.  \\
\(\mathsf{TQBF}\) & true closed quantified-Boolean-formula language & P4.6 & \hyperref[entry:symbol-registry-00308]{p.~\pageref{entry:symbol-registry-00308}} & P4.6 & Unbounded alternation.  \\
\(\mathsf U^1_{\Pi_n}\) & universal lightface Pi-one-n relation & P4.5 & \hyperref[entry:symbol-registry-00282]{p.~\pageref{entry:symbol-registry-00282}} & P4.5 & Complementary universal relation under the fixed coding.  \\
\(\mathsf U^1_{\Sigma_n}\) & universal lightface Sigma-one-n relation & P4.5 & \hyperref[entry:symbol-registry-00281]{p.~\pageref{entry:symbol-registry-00281}} & P4.5 & Pairs a valid definition code with an input satisfying that definition.  \\
\(\mathsf U^{\leq k},\mathsf X^{I},\mathsf U^{I}\) & bounded-step and real-time temporal operators & P9.5 & \hyperref[entry:symbol-registry-00662]{p.~\pageref{entry:symbol-registry-00662}} & P9.5 & The first bounds DTMC or MDP steps by natural k; X-I tests the first CTMC transition time and U-I bounds elapsed time by a declared real interval I.  \\
\(\mathsf U(y)\) & primitive-recursive normal-form output extractor & P4.1 & \hyperref[entry:symbol-registry-00242]{p.~\pageref{entry:symbol-registry-00242}} & P4.1 & Distinct from the universal program evaluator.  \\
\(\mathsf V\) & matrix value set & P8.5 & \hyperref[entry:symbol-registry-00111]{p.~\pageref{entry:symbol-registry-00111}} & P8.5 & Finite nonempty truth-value set for the generic matrix sheet; distinct from world valuation \(V\) and cumulative \(V_\alpha\).  \\
\(\mathsf V(L),\operatorname{Log}(\mathcal V)\) & order-reversing logic-variety assignments & P10.3 & \hyperref[entry:symbol-registry-00773]{p.~\pageref{entry:symbol-registry-00773}} & P10.3 & Heyting algebras validating L and formulas valid in V; inverse on the intervals IPC-to-CPL and BA-to-HA.  \\
\(\mathsf{Var}\) & first-order variable alphabet & global & \hyperref[entry:symbol-registry-00149]{p.~\pageref{entry:symbol-registry-00149}} & P2.2 & The fixed countably infinite object-language variable set; distinct from the occurrence-set operator \(\operatorname{Var}\).  \\
\(\mathsf{VarGuard}_{\mathrm{MS}}(\Lang_{\mathrm{MS}})\) & total-assignment sort-guard set & P2.3-P2.4 & \hyperref[entry:symbol-registry-00216]{p.~\pageref{entry:symbol-registry-00216}} & P2.3 & Metalinguistic one-sorted atoms constraining every object and relation variable to its declared sort; nonredundant with bound-variable guards and sort axioms, and choice-free for total assignments.  \\
\(\mathsf{WF}\) & codes of computable well-orders & P4.5 & \hyperref[entry:symbol-registry-00279]{p.~\pageref{entry:symbol-registry-00279}} & P4.5 & Subset of \(\mathsf{LO}\).  \\
\(\mathsf{WKL}_0\) & weak-Konig-lemma subsystem & P5.7-P7.5 & \hyperref[entry:symbol-registry-00445]{p.~\pageref{entry:symbol-registry-00445}} & P5.7 & RCA-zero plus the assertion that every infinite coded binary tree has an infinite path.  \\
\(\mathsf X\) & temporal next & P8.3-P9.4 & \hyperref[entry:symbol-registry-00094]{p.~\pageref{entry:symbol-registry-00094}} & P8.3 & Prefix unary at the negation rank.  \\
\(\mathsf Y A\) & temporal previous operator & P9.4 & \hyperref[entry:symbol-registry-00645]{p.~\pageref{entry:symbol-registry-00645}} & P9.4 & Prefix unary past operator on pointed words; false at position zero.  \\
\(\mathsf Z_2\) & full second-order arithmetic & P5.7-P7.5 & \hyperref[entry:symbol-registry-00443]{p.~\pageref{entry:symbol-registry-00443}} & P5.7 & RCA-zero augmented by full induction and full comprehension in the fixed two-sorted language; ambient theory for the displayed reverse-mathematics subsystems.  \\
\(\mathsf Z,\mathsf{ZC},\mathsf{KP},\mathsf{KP}^{\mathcal P}\) & set-theory system abbreviations & P6.2 & \hyperref[entry:symbol-registry-00341]{p.~\pageref{entry:symbol-registry-00341}} & P6.2 & Z omits Replacement from the stated ZF basis; ZC adds Choice; KP is the stated admissible-set basis and its power-set variant is explicitly superscripted.  \\
\(\mathsf{ZF},\mathsf{ZFC}\) & first-order set theories without and with Choice & P6 & \hyperref[entry:symbol-registry-00340]{p.~\pageref{entry:symbol-registry-00340}} & P6.1 & ZF is the displayed axiom basis without Choice; ZFC adds the displayed object-language Axiom of Choice.  \\
\(\mathsf{ZFA}\) & set theory with atoms & P6.3 & \hyperref[entry:symbol-registry-00344]{p.~\pageref{entry:symbol-registry-00344}} & P6.3 & Theory of the local permutation-model example; the cited transfer is required before drawing a pure-ZF conclusion.  \\
\(\models^{\ell}\) & local modal consequence & P8.2-P8.3 & \hyperref[entry:symbol-registry-00115]{p.~\pageref{entry:symbol-registry-00115}} & P8.2 & Truth preservation at one world in the declared modal or epistemic frame class.  \\
\(\models^{g}\) & global modal consequence & P8.2-P8.3 & \hyperref[entry:symbol-registry-00116]{p.~\pageref{entry:symbol-registry-00116}} & P8.2 & Truth preservation at all worlds in the declared modal or epistemic frame class.  \\
\(\models_{\mathrm{full}}\) & full second-order consequence & P2.3 & \hyperref[entry:symbol-registry-00076]{p.~\pageref{entry:symbol-registry-00076}} & P2.3 & Consequence over full standard structures.  \\
\(\models_{\mathrm H}\) & Henkin consequence & P2.3 & \hyperref[entry:symbol-registry-00077]{p.~\pageref{entry:symbol-registry-00077}} & P2.3 & Consequence over Henkin structures.  \\
\(\models_{\mathrm{neg}},\models_{\mathrm{pos}},\models_{\mathrm{neu}}\) & free-logic variant consequence relations & P8.4 & \hyperref[entry:symbol-registry-00586]{p.~\pageref{entry:symbol-registry-00586}} & P8.4 & Truth preservation in the negative and positive two-valued rows and designated-t preservation in the neutral strong-Kleene row; the subscript fixes the atomic and partial-term policy.  \\
\(\mu\) & monolith of a nontrivial algebra & P10.1 & \hyperref[entry:symbol-registry-00715]{p.~\pageref{entry:symbol-registry-00715}} & P10.1 & Least congruence strictly above equality when it exists; this condition characterizes nontrivial subdirect irreducibility.  \\
\(\mu g\) & unbounded minimization & P4.1 & \hyperref[entry:symbol-registry-00240]{p.~\pageref{entry:symbol-registry-00240}} & P4.1 & Definedness-sensitive least-zero operation; relational least-true use is the explicit zero-test abbreviation in P4.1.  \\
\(\multimap\) & linear implication and its residual interpretation & P8.6 & \hyperref[entry:symbol-registry-00103]{p.~\pageref{entry:symbol-registry-00103}} & P8.6 & Formula operands give IMLL or CLL linear implication; the algebraic and phase clauses interpret that same connective by the source-standard residual.  \\
\(N!\) & factorial & P4.3 & \hyperref[entry:symbol-registry-00256]{p.~\pageref{entry:symbol-registry-00256}} & P4.3 & Scope disjoint from the Part-8 exponential.  \\
\(N\) & normal worlds & P8.6 & \hyperref[entry:symbol-registry-00108]{p.~\pageref{entry:symbol-registry-00108}} & P8.6 & Nonempty distinguished subset of the broad inherited Routley--Meyer structure fixing its source-scoped validity notion.  \\
\(N,E_1,\ldots,E_q\) & semantic-clash nucleus and electron clauses & P12.3 & \hyperref[entry:symbol-registry-00618]{p.~\pageref{entry:symbol-registry-00618}} & P12.3 & Exact Slagle roles: \(N\) is the nucleus and the \(E_i\) are electrons tested against the fixed ground interpretation \(M\).  \\
\(N:W\to\mathcal P(\mathcal P(W))\) & neighborhood function & P8.2 & \hyperref[entry:symbol-registry-00505]{p.~\pageref{entry:symbol-registry-00505}} & P8.2 & The source's world-indexed neighborhood assignment.  Its function type distinguishes it from system EN and the later Routley--Meyer set of normal worlds.  \\
\(\neg\varphi\) & object-language negation & global & \hyperref[entry:symbol-registry-00006]{p.~\pageref{entry:symbol-registry-00006}} & P2.1 & Primitive unary connective; prefix and governs the immediately following formula.  \\
\(\nu\) & Heyting or open-set formula valuation & P8.1 & \hyperref[entry:symbol-registry-00491]{p.~\pageref{entry:symbol-registry-00491}} & P8.1 & Local semantic valuation extended homomorphically from propositional atoms; disjoint from the P6.3 forcing-name metavariable.  \\
\(\nu^G\) & name evaluation & P6.3 & \hyperref[entry:symbol-registry-00071]{p.~\pageref{entry:symbol-registry-00071}} & P6.3 & Interpretation by the generic filter \(G\) of a locally declared forcing name \(\nu,\nu_f\), or \(\zeta\).  \\
\(\nu,\nu_f,\zeta\) & forcing-name metavariables & P6.3 & \hyperref[entry:symbol-registry-00320]{p.~\pageref{entry:symbol-registry-00320}} & P6.3 & Members of the active P-name hierarchy; subscripts identify a name by its intended value and never turn a name into that value.  \\
\(O(A\mid B)\) & dyadic obligation & P8.4 & \hyperref[entry:symbol-registry-00588]{p.~\pageref{entry:symbol-registry-00588}} & P8.4 & Binary modality evaluated at the best accessible B-worlds; it is not the monadic formula \(O(B\to A)\).  \\
\(O(f(n))\) & asymptotic upper bound & P4.6 & \hyperref[entry:symbol-registry-00303]{p.~\pageref{entry:symbol-registry-00303}} & P4.6 & Resource-bound asymptotics.  \\
\(o(g(n))\) & strict asymptotic domination & P4.6 & \hyperref[entry:symbol-registry-00304]{p.~\pageref{entry:symbol-registry-00304}} & P4.6 & Hierarchy-gap hypothesis.  \\
\(\omega\) & least infinite ordinal and von Neumann natural-number set & P6-P7.5 & \hyperref[entry:symbol-registry-00380]{p.~\pageref{entry:symbol-registry-00380}} & P6.1 & Least inductive set under the displayed ZF construction; P7.5 uses only the ordinal standing.  \\
\(\omega_4\) & fourth uncountable initial ordinal & P6.5 & \hyperref[entry:symbol-registry-00693]{p.~\pageref{entry:symbol-registry-00693}} & P6.5 & Initial ordinal of cardinality aleph-four; used as the aleph index in Shelah's bound at aleph-omega.  \\
\(\Omega(c),\Omega(u),t(c)\) & sieves their restrictions and the maximal sieve & P10.9 & \hyperref[entry:symbol-registry-00850]{p.~\pageref{entry:symbol-registry-00850}} & P10.9 & A sieve is a set of arrows into \(c\) closed under precomposition; the maximal sieve contains the identity.  \\
\(\Omega,\mathsf{true}:1\to\Omega\) & subobject classifier and truth map & P10.9 & \hyperref[entry:symbol-registry-00847]{p.~\pageref{entry:symbol-registry-00847}} & P10.9 & Every monomorphism is a pullback along its unique characteristic arrow; the definition retains the full universal property.  \\
\(\omega_n\) & finite tower sequence below epsilon-zero & P7.5 & \hyperref[entry:symbol-registry-00456]{p.~\pageref{entry:symbol-registry-00456}} & P7.5 & Defined by omega-zero equals one and omega-(k+1) equals omega to the omega-k; used in the induction-fragment calibration.  \\
\(\omega^\omega\) & Baire function space & P6.6-P6.8 & \hyperref[entry:symbol-registry-00333]{p.~\pageref{entry:symbol-registry-00333}} & P6.6 & The set of functions from \(\omega\) to \(\omega\); Chapter 6.7 adds the product topology. Scope is disjoint from the ordinal-power role.  \\
\(\omega^\omega\) & ordinal exponentiation & P6.1 & \hyperref[entry:symbol-registry-00332]{p.~\pageref{entry:symbol-registry-00332}} & P6.1 & In the ordinal-arithmetic chapter this is the ordinal power obtained by right-recursion.  \\
\(\omega^\omega\) & ordinal exponentiation & P7.5 & \hyperref[entry:symbol-registry-00452]{p.~\pageref{entry:symbol-registry-00452}} & P7.5 & The ordinal power used in the PRA calibration; disjoint from the P6 Baire-space role.  \\
\(\Omega_{U_K}\) & halting probability of the fixed optimal prefix-free machine & P4.5 & \hyperref[entry:symbol-registry-00288]{p.~\pageref{entry:symbol-registry-00288}} & P4.5 & Not attached to an arbitrary plain machine.  \\
\(\operatorname{Add}(\omega,\kappa)\) & Cohen forcing adding kappa many reals & P6.3-P6.4 & \hyperref[entry:symbol-registry-00360]{p.~\pageref{entry:symbol-registry-00360}} & P6.3 & The explicitly displayed Fn poset; preservation and continuum claims retain their separate hypotheses.  \\
\(\operatorname{Alg}(L)\) & algebra class validating a normal modal logic & P10.3 & \hyperref[entry:symbol-registry-00743]{p.~\pageref{entry:symbol-registry-00743}} & P10.3 & Normal modal algebras validating every member of L under all algebra valuations.  \\
\(\operatorname{Alg}_\Lambda(L)\) & Boolean-operator algebras validating a finitary logic & P10.3 & \hyperref[entry:symbol-registry-00751]{p.~\pageref{entry:symbol-registry-00751}} & P10.3 & All Lambda-BAOs validating L under every algebra valuation, including trivial algebras.  \\
\(\operatorname{Alg}_{\mathrm{all}}(\Sigma,E),\operatorname{Alg}_{\mathrm{ne}}(\Sigma,E)\) & all-carrier and nonempty-carrier equational categories & P10.7 & \hyperref[entry:symbol-registry-00800]{p.~\pageref{entry:symbol-registry-00800}} & P10.7 & Both use carriers in the fixed universe and all operation-preserving functions; the latter is the full subcategory with nonempty carrier.  \\
\(\operatorname{ar}_{\mathcal F},\operatorname{ar}_{\mathcal R}\) & signature arity maps & global & \hyperref[entry:symbol-registry-00151]{p.~\pageref{entry:symbol-registry-00151}} & P2.2 & Maps the function- and relation-symbol families to their finite arities.  \\
\(\operatorname{BB}_{1}(n)\) & Busy Beaver one-count function & P4.5 & \hyperref[entry:symbol-registry-00291]{p.~\pageref{entry:symbol-registry-00291}} & P4.5 & Maximum number of one-filled cells at halt under the displayed convention.  \\
\(\operatorname{BB}_{\mathrm{step}}(n)\) & Busy Beaver step function & P4.5 & \hyperref[entry:symbol-registry-00290]{p.~\pageref{entry:symbol-registry-00290}} & P4.5 & Maximum halting transition count under the displayed machine convention.  \\
\(\operatorname{bin}_m(i)\) & fixed-width binary address & P4.6 & \hyperref[entry:symbol-registry-00235]{p.~\pageref{entry:symbol-registry-00235}} & P4.6 & Length-m expansion padded with leading zeros for succinct-instance bit addressing.  \\
\(\operatorname{bin}(n)\) & canonical binary word & P4.1-P4.2 & \hyperref[entry:symbol-registry-00234]{p.~\pageref{entry:symbol-registry-00234}} & P4.1 & Zero is represented by the one-symbol word zero.  \\
\(\operatorname{Bool}_{\mathbb B}(\varphi)\) & ambient Boolean truth value of a formula & P6.3 & \hyperref[entry:symbol-registry-00323]{p.~\pageref{entry:symbol-registry-00323}} & P6.3 & Value in an ambient-complete Boolean algebra; this is not the two-valued satisfaction relation.  \\
\(\operatorname{Bor}(X)\) & Borel sigma-algebra of a Polish space & P6.7 & \hyperref[entry:symbol-registry-00336]{p.~\pageref{entry:symbol-registry-00336}} & P6.7 & The least sigma-algebra containing the open subsets of \(X\).  \\
\(\operatorname{cf}(\alpha),\operatorname{cf}(P,\leq)\) & cofinality of an ordinal or ordered set & P6.3-P6.8 & \hyperref[entry:symbol-registry-00352]{p.~\pageref{entry:symbol-registry-00352}} & P6.3 & For an ordinal: least cofinal order type. For a nonempty ordered set: least cardinality of a cofinal family. Regularity is cf(kappa)=kappa.  \\
\(\operatorname{Cg}_A(a,b)\) & principal congruence generated by a pair & P10.1 & \hyperref[entry:symbol-registry-00712]{p.~\pageref{entry:symbol-registry-00712}} & P10.1 & Intersection of every congruence on A containing the pair (a,b).  \\
\(\operatorname{Clop}(g)\) & Boolean dual of a continuous map & P10.2 & \hyperref[entry:symbol-registry-00723]{p.~\pageref{entry:symbol-registry-00723}} & P10.2 & For g from X to Y sends a clopen subset of Y to its inverse image in X.  \\
\(\operatorname{Clop}(X)\) & Boolean algebra of clopen subsets & P10.2 & \hyperref[entry:symbol-registry-00719]{p.~\pageref{entry:symbol-registry-00719}} & P10.2 & Intersection union and relative complement act on the clopens of X; the bounds are the empty set and X.  \\
\(\operatorname{ClopUp}(X)\) & lattice of clopen upsets & P10.3 & \hyperref[entry:symbol-registry-00757]{p.~\pageref{entry:symbol-registry-00757}} & P10.3 & Bounded distributive lattice for a Priestley space; carries the residual complement-of-downclosure operation for an Esakia space.  \\
\(\operatorname{Cn}(T)\) & semantic consequence closure & global & \hyperref[entry:symbol-registry-00158]{p.~\pageref{entry:symbol-registry-00158}} & P2.2 & Sentences of the active signature semantically entailed by theory \(T\).  \\
\(\operatorname{code}_k(\vec n)\) & fixed-arity tuple code & P4.1-P4.3 & \hyperref[entry:symbol-registry-00237]{p.~\pageref{entry:symbol-registry-00237}} & P4.1 & Defined recursively from the numerical pairing function.  \\
\(\operatorname{code}(L)\) & propositional-literal ordering code & P2.1 & \hyperref[entry:symbol-registry-00133]{p.~\pageref{entry:symbol-registry-00133}} & P2.1 & \(\operatorname{code}(p_n)=2n\) and \(\operatorname{code}(\neg p_n)=2n+1\); fixes one formula string for each clause.  \\
\(\operatorname{Coll}(\kappa,\lambda)\) & Levy collapse forcing & P6.3 & \hyperref[entry:symbol-registry-00361]{p.~\pageref{entry:symbol-registry-00361}} & P6.3 & Partial functions from kappa to lambda of size below kappa under the catalog convention; collapse claims are cited.  \\
\(\operatorname{Comp}(P)\) & Clark completion of a normal program & P9.7 & \hyperref[entry:symbol-registry-00680]{p.~\pageref{entry:symbol-registry-00680}} & P9.7 & Predicate-defining biconditionals plus the explicitly stated signature-relative free-term equality axioms; distinct from least and stable models.  \\
\(\operatorname{Con}A,\Delta_A,\nabla_A\) & congruence lattice and its bounds & P10.1 & \hyperref[entry:symbol-registry-00711]{p.~\pageref{entry:symbol-registry-00711}} & P10.1 & Congruences ordered by inclusion; Delta is equality and nabla is the universal relation; meets are intersections.  \\
\(\operatorname{crit}(j)\) & critical point of an elementary embedding & P6.6-P6.8 & \hyperref[entry:symbol-registry-00337]{p.~\pageref{entry:symbol-registry-00337}} & P6.6 & Least ordinal moved by \(j\); subscripted embedding symbols such as \(j_U\) use the same operation.  \\
\(\operatorname{CSP}(B)\) & fixed-template constraint-satisfaction problem & P9.7 & \hyperref[entry:symbol-registry-00685]{p.~\pageref{entry:symbol-registry-00685}} & P9.7 & B is the fixed finite relational target and the input asks for a homomorphism from A to B.  \\
\(\operatorname{Cut}_A(D_0,D_1),\operatorname{Cut}(D),A_c\) & Cut node Cut-node set and node-indexed cut formula & P7.5 & \hyperref[entry:symbol-registry-00437]{p.~\pageref{entry:symbol-registry-00437}} & P7.5 & Metalinguistic notation for one Cut with premises D0 D1, the finite set of Cut nodes of D, and the formula attached to node c in the ordinal assignment.  \\
\(\operatorname{Def}(A)\) & definable subsets of a membership structure & P6.3 & \hyperref[entry:symbol-registry-00386]{p.~\pageref{entry:symbol-registry-00386}} & P6.3 & Subsets of A first-order definable over \((A,\in)\) with finitely many parameters from A.  \\
\(\operatorname{dev}(M)\) & complete development of a lambda term & P7.4 & \hyperref[entry:symbol-registry-00411]{p.~\pageref{entry:symbol-registry-00411}} & P7.4 & Metalinguistic recursion contracting all residuals exposed by the starting syntax; not a reduction relation.  \\
\(\operatorname{DGen}_\Lambda(L)\) & descriptive general frames validating a finitary logic & P10.3 & \hyperref[entry:symbol-registry-00754]{p.~\pageref{entry:symbol-registry-00754}} & P10.3 & The nonempty descriptive Lambda-general frames validating every theorem under all admissible valuations.  \\
\(\operatorname{Diag}(\mathcal C),=_{\mathcal C}\) & full true diagram and its equational theory & P10.8 & \hyperref[entry:symbol-registry-00792]{p.~\pageref{entry:symbol-registry-00792}} & P10.8 & Contains every true typed equation in every finite context of the named language; no effective enumeration is asserted.  \\
\(\operatorname{Diag}(\mathcal M)\) & atomic diagram & P3 & \hyperref[entry:symbol-registry-00044]{p.~\pageref{entry:symbol-registry-00044}} & P3.1 & Atomic sentences with parameters true in \(\mathcal M\).  \\
\(\operatorname{dom}X\) & domain operator & global & \hyperref[entry:symbol-registry-00023]{p.~\pageref{entry:symbol-registry-00023}} & P2.2 & Operand type fixes the structure-domain, function-domain, or name-domain reading; vertical bars are never used for domains.  \\
\(\operatorname{ElDiag}(\mathcal M)\) & elementary diagram & P2.2-P3 & \hyperref[entry:symbol-registry-00045]{p.~\pageref{entry:symbol-registry-00045}} & P2.2 & All parameter sentences true in \(\mathcal M\).  \\
\(\operatorname{Eq}_{\Delta}(J,J')\) & multi-sorted MLTT equality premise & P7.4 & \hyperref[entry:symbol-registry-00431]{p.~\pageref{entry:symbol-registry-00431}} & P7.4 & Eliminable metalinguistic abbreviation for type or term judgmental equality after transport to the common context Delta.  \\
\(\operatorname{ev}_c,\Lambda h\) & presheaf evaluation and transposition & P10.9 & \hyperref[entry:symbol-registry-00849]{p.~\pageref{entry:symbol-registry-00849}} & P10.9 & Evaluation uses the representing identity arrow; transposition uses every restriction of the parameter element.  \\
\(\operatorname{ext}_n(r),s^{\mathrm{ext}}\) & extensional-collapse map and induced assignment & P2.3 & \hyperref[entry:symbol-registry-00211]{p.~\pageref{entry:symbol-registry-00211}} & P2.3 & Maps a relation-sort object to its n-ary extension and transports relation-variable assignments through that map.  \\
\(\operatorname{Fm}_\Lambda(X)\) & formulas for a finitary modal signature & P10.3 & \hyperref[entry:symbol-registry-00748]{p.~\pageref{entry:symbol-registry-00748}} & P10.3 & Finite Boolean formulas over X with one constructor of arity n-f for every symbol f in Lambda.  \\
\(\operatorname{Fn}(I,J,{<}\delta)\) & bounded-size partial-function poset & P6.3 & \hyperref[entry:symbol-registry-00359]{p.~\pageref{entry:symbol-registry-00359}} & P6.3 & Partial functions from I to J of size below delta, ordered by reverse inclusion.  \\
\(\operatorname{form}_{\Lang}(C),\operatorname{ucl}_{\Lang}(C)\) & clause formula and universal closure & P2.2-P12.2 & \hyperref[entry:symbol-registry-00169]{p.~\pageref{entry:symbol-registry-00169}} & P2.2 & Fixed literal-order disjunction for clause \(C\) and its universal closure over every occurring variable.  \\
\(\operatorname{form}_{\mathsf{CPL}}(C)\) & formula interpretation of a clause & P2.1 & \hyperref[entry:symbol-registry-00134]{p.~\pageref{entry:symbol-registry-00134}} & P2.1 & Increasing-code disjunction of a finite clause; the empty clause maps to falsum.  \\
\(\operatorname{FV}(\varphi),\operatorname{FV}(\Gamma),\operatorname{FV}(M),\operatorname{FV}(\alpha)\) & free-variable operation & global & \hyperref[entry:symbol-registry-00034]{p.~\pageref{entry:symbol-registry-00034}} & P2.2 & Variables with a free occurrence in the locally declared formula, context, term, or type syntax; context action takes the union over members.  \\
\(\operatorname{Gen}(L),\operatorname{Fr}(L)\) & nonempty validating general-frame and ordinary-frame classes & P10.3 & \hyperref[entry:symbol-registry-00746]{p.~\pageref{entry:symbol-registry-00746}} & P10.3 & General frames validate L using admissible valuations; ordinary frames validate it using all subset valuations.  \\
\(\operatorname{gfp}(F)\) & greatest fixed point of a monotone map & P9.5-P9.7 & \hyperref[entry:symbol-registry-00654]{p.~\pageref{entry:symbol-registry-00654}} & P9.5 & The greatest set fixed by the displayed monotone operator; CTL formulas continue to use the object-language nu binder.  \\
\(\operatorname{ht}(D)\) & derivation height & P7.2 & \hyperref[entry:symbol-registry-00396]{p.~\pageref{entry:symbol-registry-00396}} & P7.2 & Greatest inference-step count on a leaf-to-root branch; distinct from tree-node height in P6.4.  \\
\(\operatorname{Id}_A(a,b)\) & MLTT intensional identity type & P7.4 & \hyperref[entry:symbol-registry-00427]{p.~\pageref{entry:symbol-registry-00427}} & P7.4 & Proposition-as-type of identifications between terms a and b of A; no judgmental equality reflection is assumed.  \\
\(\operatorname{Im}(m),q,i\) & componentwise image and its factorization & P10.9 & \hyperref[entry:symbol-registry-00846]{p.~\pageref{entry:symbol-registry-00846}} & P10.9 & The natural map factors through its image by a componentwise surjection followed by an inclusion; a monic map is isomorphic to that inclusion.  \\
\(\operatorname{Inf}(r)\) & infinitely recurring run states & P8.3-P9.4 & \hyperref[entry:symbol-registry-00541]{p.~\pageref{entry:symbol-registry-00541}} & P8.3 & States occurring infinitely often in run r; Buchi acceptance intersects F while a Rabin pair separately tests its E and F components.  \\
\(\operatorname{int}_{X}\) & topological interior operator & P8.1 & \hyperref[entry:symbol-registry-00492]{p.~\pageref{entry:symbol-registry-00492}} & P8.1 & Interior in the declared topological space X used by the open-set implication clause.  \\
\(\operatorname{len}(x)\) & bit length of a finite binary string & P4.5-P6.2 & \hyperref[entry:symbol-registry-00283]{p.~\pageref{entry:symbol-registry-00283}} & P4.5 & Replaces forbidden vertical-bar length notation.  \\
\([\operatorname{lfp}_{R,\bar x}\varphi](\bar t)\) & least fixed-point formula constructor & P4.6-P9 & \hyperref[entry:symbol-registry-00316]{p.~\pageref{entry:symbol-registry-00316}} & P4.6 & Requires positive occurrence of the relation variable and denotes the least monotone fixed point.  \\
\(\operatorname{Mod}(T)\) & model class & global & \hyperref[entry:symbol-registry-00035]{p.~\pageref{entry:symbol-registry-00035}} & P2.2 & Structures satisfying every sentence of \(T\).  \\
\(\operatorname{Mod}(T),U\) & Lawvere models and the underlying-set functor & P10.7 & \hyperref[entry:symbol-registry-00804]{p.~\pageref{entry:symbol-registry-00804}} & P10.7 & Finite-product-preserving functors to Set and all natural transformations; \(U(M)=M(1)\) and \(U(\alpha)=\alpha_1\).  \\
\(\operatorname{Nat}(F,G)\) & family of natural transformations & P10.7-P10.9 & \hyperref[entry:symbol-registry-00778]{p.~\pageref{entry:symbol-registry-00778}} & P10.7 & Component families in the containing ambient universe; Yoneda gives a small set of parameters for transformations from a representable.  \\
\(\operatorname{nnf},\operatorname{cl}\) & ALC normalization and finite closure operations & P9.3 & \hyperref[entry:symbol-registry-00641]{p.~\pageref{entry:symbol-registry-00641}} & P9.3 & Metalinguistic operations used by blocking completion and type elimination; they add no concept constructors.  \\
\(\operatorname{nnf}(\varphi),\operatorname{nnf}^{\neg}(\varphi)\) & first-order NNF translations & P2.2 & \hyperref[entry:symbol-registry-00166]{p.~\pageref{entry:symbol-registry-00166}} & P2.2 & Simultaneously recursive maps equivalent to \(\varphi\) and \(\neg\varphi\), respectively.  \\
\(\operatorname{Ord}\) & proper class of all ordinals & P6 & \hyperref[entry:symbol-registry-00379]{p.~\pageref{entry:symbol-registry-00379}} & P6.1 & Metalinguistic class notation and never a set-valued union index.  \\
\(\operatorname{otp}(A,<)\) & order type of a declared linear order & P5.4-P7 & \hyperref[entry:symbol-registry-00317]{p.~\pageref{entry:symbol-registry-00317}} & P5.4 & Metalinguistic isomorphism type of the displayed ordered set; not an ordinal unless the order is well founded.  \\
\(\operatorname{pack}_\Gamma,\operatorname{unpack}_\Gamma\) & inverse context packing arrows & P10.8 & \hyperref[entry:symbol-registry-00794]{p.~\pageref{entry:symbol-registry-00794}} & P10.8 & Compare \(\Gamma\) with the singleton named context for \(P_\Gamma\); the packing term and its context arrow have different codomains.  \\
\(\operatorname{pair}_{\N}(a,b)\) & numerical pairing function & P4.1-P4.3 & \hyperref[entry:symbol-registry-00236]{p.~\pageref{entry:symbol-registry-00236}} & P4.1 & Distinct from the set-theoretic ordered pair.  \\
\(\operatorname{pcf}(\mathfrak a)\) & possible cofinalities of a product & P6.5 & \hyperref[entry:symbol-registry-00692]{p.~\pageref{entry:symbol-registry-00692}} & P6.5 & Cofinalities of the ultrafilter quotients of the product of the infinite regular cardinals in a; principal ultrafilters are included.  \\
\(\operatorname{Pf}(H)\) & prime-filter order of a Heyting algebra & P10.3 & \hyperref[entry:symbol-registry-00767]{p.~\pageref{entry:symbol-registry-00767}} & P10.3 & All proper prime lattice filters ordered by inclusion; filters need not be maximal.  \\
\(\operatorname{Pre}_{\exists},\operatorname{Pre}_{\forall}\) & existential and universal predecessor maps & P9.5 & \hyperref[entry:symbol-registry-00652]{p.~\pageref{entry:symbol-registry-00652}} & P9.5 & Algorithmic maps from state sets to state sets in a finite transition structure; not CTL path quantifiers.  \\
\(\operatorname{Prf}_{\mathsf P_{\mathrm{pf}}}(\pi,\varphi)\) & encoded propositional proof relation & P7.8 & \hyperref[entry:symbol-registry-00479]{p.~\pageref{entry:symbol-registry-00479}} & P7.8 & Polynomial-time decidable relation whose range is exactly TAUT under the fixed malformed-input convention.  \\
\(\operatorname{ProofRel}_T,\operatorname{RProofRel}_T,\operatorname{NegCode}\) & metalinguistic primitive-recursive proof relations and negation-code function & P5 & \hyperref[entry:symbol-registry-00060]{p.~\pageref{entry:symbol-registry-00060}} & P5.6 & Relations and function on natural-number codes used before choosing object-language representatives; none is a term or predicate of the fixed arithmetic language.  \\
\(\operatorname{Reach}(a,b)\) & second-order reachability formula & P2.3 & \hyperref[entry:symbol-registry-00199]{p.~\pageref{entry:symbol-registry-00199}} & P2.3 & The full-semantic closure formula for finite directed reachability; its Henkin converse can fail.  \\
\(\operatorname{rect}(\varphi,V),\operatorname{pren}(\varphi)\) & rectification and prenex maps & P2.2 & \hyperref[entry:symbol-registry-00167]{p.~\pageref{entry:symbol-registry-00167}} & P2.2 & Metalinguistic constructions; the forbidden-variable parameter belongs only to rectification.  \\
\(\operatorname{rk}_{\mathrm{fm}}(\varphi)\) & formula rank & P7.2 & \hyperref[entry:symbol-registry-00395]{p.~\pageref{entry:symbol-registry-00395}} & P7.2 & Number of logical symbols in the displayed formula; the subscript separates it from set-theoretic rank.  \\
\(\operatorname{rk}(x)\) & set-theoretic rank & P6 & \hyperref[entry:symbol-registry-00066]{p.~\pageref{entry:symbol-registry-00066}} & P6.1 & Least \(\alpha\) with \(x\subseteq V_\alpha\).  \\
\(\operatorname{size}(a)\) & bit length of an encoded proof or formula object & P7.8 & \hyperref[entry:symbol-registry-00481]{p.~\pageref{entry:symbol-registry-00481}} & P7.8 & Defined from the fixed binary encoding; deliberately not vertical-bar cardinality notation.  \\
\(\operatorname{Sk}(\sigma)\) & Skolem form & P2.2 & \hyperref[entry:symbol-registry-00168]{p.~\pageref{entry:symbol-registry-00168}} & P2.2 & Universal Skolem sentence in the expanded signature; equisatisfiability and reduct directions retain their separate choice standing.  \\
\(\operatorname{ST}_x,P_p\) & modal standard translation and atom predicate & P8.2 & \hyperref[entry:symbol-registry-00501]{p.~\pageref{entry:symbol-registry-00501}} & P8.2 & ST at free world variable x maps each modal atom p to the corresponding unary first-order predicate \(P_p\).  \\
\(\operatorname{Step}_{M,x},\operatorname{Eq},\operatorname{Sel}_b\) & QBF reduction transition equality and selector formulas & P9.8 & \hyperref[entry:symbol-registry-00686]{p.~\pageref{entry:symbol-registry-00686}} & P9.8 & Metalinguistic polynomial-size Boolean formulas on configuration blocks; the selector permits one textual recursive occurrence.  \\
\(\operatorname{Sub}(\varphi)\) & subformula-set operation & P2.1 & \hyperref[entry:symbol-registry-00129]{p.~\pageref{entry:symbol-registry-00129}} & P2.1 & Set of subformulas fixed by the recursive CPL definition.  \\
\(\operatorname{Sub}(X)\) & ordered set of subpresheaves & P10.9 & \hyperref[entry:symbol-registry-00844]{p.~\pageref{entry:symbol-registry-00844}} & P10.9 & Restriction-stable families of component subsets; canonical images represent categorical subobjects and their factorization order.  \\
\(\operatorname{supp}(\Gamma)\) & support of a finite multiset & P2.1-P7.1 & \hyperref[entry:symbol-registry-00142]{p.~\pageref{entry:symbol-registry-00142}} & P2.1 & Formulas having positive multiplicity in the declared propositional or first-order finite-multiset context; reused in P7.1.  \\
\(\operatorname{supp}_{\mathrm{var}}(\theta)\) & variable support of a term substitution & P2.2-P12.3 & \hyperref[entry:symbol-registry-00170]{p.~\pageref{entry:symbol-registry-00170}} & P2.2 & Variables not fixed by \(\theta\); the subscript prevents collision with finite-multiset support.  \\
\(\operatorname{Sym}(A),\operatorname{Stab}_{\mathcal G}(x)\) & full permutation group and stabilizer & P6.3 & \hyperref[entry:symbol-registry-00331]{p.~\pageref{entry:symbol-registry-00331}} & P6.3 & The stabilizer is a subgroup of the declared permutation group; support and symmetry are defined from it.  \\
\(\operatorname{Term}(\Lang),\operatorname{Fm}(\Lang),\operatorname{Sent}(\Lang)\) & first-order term formula and sentence sets & global & \hyperref[entry:symbol-registry-00152]{p.~\pageref{entry:symbol-registry-00152}} & P2.2 & Least syntactic sets generated by the active first-order grammar; sentences are formulas with no free variables.  \\
\(\operatorname{Th}(\mathcal M)\) & theory of a structure & global & \hyperref[entry:symbol-registry-00036]{p.~\pageref{entry:symbol-registry-00036}} & P2.2 & Sentences true in \(\mathcal M\).  \\
\(\operatorname{tr}_{\mathrm{MS}},\mathsf{Ax}_{\mathrm{sort}}\) & many-sorted guard translation and sort theory & P2.3-P2.4 & \hyperref[entry:symbol-registry-00215]{p.~\pageref{entry:symbol-registry-00215}} & P2.3 & The recursive translation guards bound variables; the theory asserts sort nonemptiness typing closure and disjointness.  \\
\(\operatorname{tr}_{\mathsf H}(\varphi)\) & full-language-to-Hilbert translation & P2.1 & \hyperref[entry:symbol-registry-00143]{p.~\pageref{entry:symbol-registry-00143}} & P2.1 & Metalinguistic recursive map into the negation-implication fragment; avoids the globally reserved sentence metavariable \(\tau\).  \\
\(\operatorname{ucl}(\varphi)\) & formula-level universal closure & P2.3 & \hyperref[entry:symbol-registry-00194]{p.~\pageref{entry:symbol-registry-00194}} & P2.3 & Universally closes exactly the displayed free object and relation parameters; distinct from the clause-language subscripted closure.  \\
\(\operatorname{Ult}(A)\) & ultrafilters of the Boolean reduct & P10.3 & \hyperref[entry:symbol-registry-00752]{p.~\pageref{entry:symbol-registry-00752}} & P10.3 & The set of proper maximal Boolean filters, used as the canonical frame carrier of A.  \\
\(\operatorname{Ult}(V,\mathcal U)\) & universe ultrapower by a kappa-complete ultrafilter & P6.6 & \hyperref[entry:symbol-registry-00377]{p.~\pageref{entry:symbol-registry-00377}} & P6.6 & Class ultrapower whose well-foundedness and transitive collapse are cited.  \\
\(\operatorname{Up}(W)\) & Heyting algebra of upsets of a preorder & P10.3 & \hyperref[entry:symbol-registry-00765]{p.~\pageref{entry:symbol-registry-00765}} & P10.3 & Implication holds at w when every larger point in the antecedent upset belongs to the consequent upset.  \\
\(\operatorname{val}^{\mathcal M}_s(t)\) & term value & global & \hyperref[entry:symbol-registry-00156]{p.~\pageref{entry:symbol-registry-00156}} & P2.2 & Recursive denotation of the first-order term \(t\) in structure \(\mathcal M\) under assignment \(s\).  \\
\(\operatorname{Var}(e)\) & occurring-variable set & global & \hyperref[entry:symbol-registry-00153]{p.~\pageref{entry:symbol-registry-00153}} & P2.2 & All variables occurring in the declared term, formula, literal, clause, or finite clause set \(e\), whether free or bound; for a finite clause set the operator is the union of the occurrence sets of its clauses.  \\
\(\operatorname{Veb}_{\alpha}(\beta)\) & Veblen hierarchy of normal ordinal functions & P7.5 & \hyperref[entry:symbol-registry-00458]{p.~\pageref{entry:symbol-registry-00458}} & P7.5 & Veb-zero is ordinal exponentiation by omega; later levels enumerate common fixed points of all earlier levels.  \\
\(\operatorname{wp}(C,Q),\operatorname{wlp}(C,Q)\) & weakest and weakest liberal preconditions & P9.6 & \hyperref[entry:symbol-registry-00672]{p.~\pageref{entry:symbol-registry-00672}} & P9.6 & Predicate transformers distinguished by whether termination is required.  \\
\(\otimes\) & linear tensor and its homomorphic semantic interpretation & P8.6 & \hyperref[entry:symbol-registry-00102]{p.~\pageref{entry:symbol-registry-00102}} & P8.6 & Formula operands give the exact IMLL or CLL connective; the IMLL algebra and CLL phase clauses interpret that same connective with their source-standard tensor operation.  \\
\(\overline L\) & literal complement & P2.1-P7.8 & \hyperref[entry:symbol-registry-00132]{p.~\pageref{entry:symbol-registry-00132}} & P2.1 & Metalinguistic involution exchanging \(p_n\) and \(\neg p_n\); P7.8 reuses it on the extended-resolution sheet; disjoint from numeral bars.  \\
\(\overline{\vec a}\) & tuple of arithmetic numerals & P4.3-P5 & \hyperref[entry:symbol-registry-00229]{p.~\pageref{entry:symbol-registry-00229}} & P4.3 & Componentwise numeral formation from a declared tuple of standard naturals.  \\
\(p\Vdash_{M,\mathbb P}\varphi\) & imported recursive forcing assertion & P6.3 & \hyperref[entry:symbol-registry-00074]{p.~\pageref{entry:symbol-registry-00074}} & P6.3 & Formula-by-formula recursive assertion imported by the cited forcing theorem; equivalent there to the separately registered semantic relation. Bare \(p\Vdash\varphi\) is its local abbreviation only under the renewed ccc/name-counting convention.  \\
\(p\Vdash^{\mathrm{sem}}_{M,\mathbb P}\varphi\) & semantic set-theoretic forcing & P6.3 & \hyperref[entry:symbol-registry-00073]{p.~\pageref{entry:symbol-registry-00073}} & P6.3 & Exact external relation: every M-generic filter containing p yields an extension satisfying the formula at the interpreted names.  \\
\(P(B\mid A)\) & conditional probability & P8.4 & \hyperref[entry:symbol-registry-00595]{p.~\pageref{entry:symbol-registry-00595}} & P8.4 & Defined when \(P(A)>0\); probability P is distinct from the deontic permission abbreviation \(\mathsf P_O\).  \\
\(\{P\}\,C\,\{Q\}\) & Hoare partial or total correctness triple & P9.6 & \hyperref[entry:symbol-registry-00670]{p.~\pageref{entry:symbol-registry-00670}} & P9.6 & The active correctness tag determines whether termination is required; braces do not denote sets in this role.  \\
\(p_i^n,\langle f_i\rangle_{i<n},!_n\) & chosen powers and tuples in a Lawvere theory & P10.7 & \hyperref[entry:symbol-registry-00798]{p.~\pageref{entry:symbol-registry-00798}} & P10.7 & The arrows \(p_i^n:n\to1\) exhibit \(n=1^n\), with \(p_0^1=1_1\); the empty tuple is the unique arrow \(n\to0\).  \\
\(P(I),\top_I,\bot_I\) & predicate fibre and its bounds & P10.9 & \hyperref[entry:symbol-registry-00814]{p.~\pageref{entry:symbol-registry-00814}} & P10.9 & A small Heyting algebra of predicates in context object \(I\); its order is semantic entailment in that context.  \\
\(P^M\) & Gelfond-Lifschitz reduct & P9.7 & \hyperref[entry:symbol-registry-00681]{p.~\pageref{entry:symbol-registry-00681}} & P9.7 & Positive program obtained relative to atom set M by deleting defeated rules and remaining default-negated literals.  \\
\(P^n\) & predicate-constant metavariable & P2.3 & \hyperref[entry:symbol-registry-00186]{p.~\pageref{entry:symbol-registry-00186}} & P2.3 & An n-ary nonlogical predicate constant in the active relational second-order signature.  \\
\((p_n),\mathcal D\) & descending conditions and dense requirements & P10.3 & \hyperref[entry:symbol-registry-00762]{p.~\pageref{entry:symbol-registry-00762}} & P10.3 & Rasiowa-Sikorski construction has p-n-plus-one stronger than p-n and in D-n; its upward closure is generic for the given family.  \\
\(p_n; P,Q,R,S\) & object-language propositional variables & P2.1 & \hyperref[entry:symbol-registry-00125]{p.~\pageref{entry:symbol-registry-00125}} & P2.1 & \(p_n\) is the indexed variable schema; uppercase letters denote specified members of \(\mathsf{Prop}\) in examples and diagrams, never formula metavariables.  \\
\(P*Q\) & separating conjunction & P9.6 & \hyperref[entry:symbol-registry-00675]{p.~\pageref{entry:symbol-registry-00675}} & P9.6 & Assertion connective splitting one heap into disjoint subheaps satisfying P and Q; distinct from ordinary conjunction and arithmetic multiplication.  \\
\(p_U=\bigwedge e[U]\) & atom in a dense compact Boolean completion & P10.3 & \hyperref[entry:symbol-registry-00739]{p.~\pageref{entry:symbol-registry-00739}} & P10.3 & The meet of the images of all elements of an ultrafilter U; compactness makes it nonzero and density makes it an atom.  \\
\(P(X)=\mathcal P(X)\) & powerset instance of a predicate fibre & P10.9 & \hyperref[entry:symbol-registry-00824]{p.~\pageref{entry:symbol-registry-00824}} & P10.9 & Boolean algebra of all subsets of a small set \(X\), including the one-element fibre for empty \(X\).  \\
\(P(X)=\Omega^X\) & power object of a presheaf & P10.9 & \hyperref[entry:symbol-registry-00852]{p.~\pageref{entry:symbol-registry-00852}} & P10.9 & Here \(X\) is a presheaf and \(P(X)\) is an object of \(\widehat{\mathcal C}\), distinct from the earlier powerset-fibre usage for sets.  \\
\(\PA\) & Peano arithmetic & P3-P7.5 & \hyperref[entry:symbol-registry-00439]{p.~\pageref{entry:symbol-registry-00439}} & P3.3 & The fixed first-order theory defined in P5.1; earlier table and later metatheoretic uses denote that same theory.  \\
\(\perp\subseteq M\) & Girard phase-space pole & P8.6 & \hyperref[entry:symbol-registry-00570]{p.~\pageref{entry:symbol-registry-00570}} & P8.6 & Girard's exact perpendicular glyph, locally typed as the distinguished subset of the phase monoid rather than formula linear negation.  \\
\(\Phi\) & Fitting quantified-modal formula metavariable & P8.2 & \hyperref[entry:symbol-registry-00514]{p.~\pageref{entry:symbol-registry-00514}} & P8.2 & Ranges over formulas in the source presentation; locally typed apart from the filtration-set use of Phi.  \\
\(\Phi\) & ordered Lambek context metavariable & P8.6 & \hyperref[entry:symbol-registry-00580]{p.~\pageref{entry:symbol-registry-00580}} & P8.6 & Possibly empty context segment in rule notation subject to every displayed antecedent remaining nonempty.  \\
\(\Phi\vdash\psi[\Gamma]\) & contextual predicate sequent & P10.9 & \hyperref[entry:symbol-registry-00829]{p.~\pageref{entry:symbol-registry-00829}} & P10.9 & The conjunction of the finite assumptions is below the conclusion in \(P(I_\Gamma)\); unused variable declarations are retained.  \\
\(\Phi_{M,x}\) & closed QBF produced by the PSPACE reduction & P9.8 & \hyperref[entry:symbol-registry-00689]{p.~\pageref{entry:symbol-registry-00689}} & P9.8 & The polynomial-size formula R-m applied to the initial and canonical accepting configuration blocks.  \\
\(\Phi,\sim_{\Phi},[w]_{\Phi}\) & filtration set equivalence and quotient world & P8.2 & \hyperref[entry:symbol-registry-00504]{p.~\pageref{entry:symbol-registry-00504}} & P8.2 & Finite subformula-closed set Phi induces modal equivalence and its world classes in the declared filtration.  \\
\(\phi/\theta\) & quotient congruence above a kernel & P10.1 & \hyperref[entry:symbol-registry-00714]{p.~\pageref{entry:symbol-registry-00714}} & P10.1 & For theta contained in phi this relation on A/theta relates exactly the classes of phi-related representatives.  \\
\(\pi\) & encoded propositional proof metavariable & P7.8 & \hyperref[entry:symbol-registry-00480]{p.~\pageref{entry:symbol-registry-00480}} & P7.8 & Binary proof string for the fixed Cook-Reckhow encoding; no standing as a projection or permutation.  \\
\(\pi_0,\pi_1\) & primitive-recursive pairing projections & P7.6 & \hyperref[entry:symbol-registry-00471]{p.~\pageref{entry:symbol-registry-00471}} & P7.6 & First and second projections for the fixed numerical pairing function used in number realizability.  \\
\(\Pi^1_1\text{-}\mathsf{CA}_0\) & Pi-one-one comprehension theory & P5.7-P7.5 & \hyperref[entry:symbol-registry-00449]{p.~\pageref{entry:symbol-registry-00449}} & P5.7 & RCA-zero plus comprehension for formulas with universal set quantifiers over an arithmetical matrix.  \\
\(\Pi_D,p_a\) & classical encoding of prime lattice filters & P10.3 & \hyperref[entry:symbol-registry-00760]{p.~\pageref{entry:symbol-registry-00760}} & P10.3 & The proposition p-a records membership of a; Pi-D imposes only bounded-lattice homomorphism laws to two, not Heyting implication preservation.  \\
\(\pi_{I,A},f\times1_A\) & context projection and parameter substitution & P10.9 & \hyperref[entry:symbol-registry-00817]{p.~\pageref{entry:symbol-registry-00817}} & P10.9 & The projection \(I\times A\to I\) forgets one variable; a parameter map lifts to the displayed product substitution square.  \\
\(\pi_i,\iota_i,\langle f_i\rangle_i\) & product projections coproduct injections and pairing & P10.7-P10.9 & \hyperref[entry:symbol-registry-00781]{p.~\pageref{entry:symbol-registry-00781}} & P10.7 & Domains and codomains distinguish projections from injections; pairing is the unique arrow with its specified product components.  \\
\(\Pi_k\text{-}\mathsf{SAT}\) & universal-start fixed-alternation QBF language & P4.6 & \hyperref[entry:symbol-registry-00310]{p.~\pageref{entry:symbol-registry-00310}} & P4.6 & Fixed \(k\).  \\
\(\Pi^p_k\) & universal-start polynomial-hierarchy class & P4.6 & \hyperref[entry:symbol-registry-00306]{p.~\pageref{entry:symbol-registry-00306}} & P4.6 & Bounded alternating string quantifiers.  \\
\(\Pi_{x:A}B,\Sigma_{x:A}B\) & MLTT dependent product and dependent sum type formers & P7.4 & \hyperref[entry:symbol-registry-00426]{p.~\pageref{entry:symbol-registry-00426}} & P7.4 & Bind x in B on sheet:mltt-intensional; distinct from analytical classes and sequent-context metavariables.  \\
\(\Pr_s,\Pr_s^{\min},\Pr_s^{\max}\) & path and extremal scheduler probabilities & P9.5 & \hyperref[entry:symbol-registry-00660]{p.~\pageref{entry:symbol-registry-00660}} & P9.5 & The undecorated measure belongs to a DTMC or CTMC or to an MDP after fixing a scheduler; min and max quantify over MDP schedulers.  \\
\(\preceq_K,\min([A],\preceq_K)\) & AGM faithful preorder and minimal input valuations & P8.7 & \hyperref[entry:symbol-registry-00604]{p.~\pageref{entry:symbol-registry-00604}} & P8.7 & A total preorder assigned to prior belief set K; the minimum selects its most preferred A-valuations.  \\
\(\preceq_w,\operatorname{Best}_w(B)\) & world-relative ideality preorder and best-world selector & P8.4 & \hyperref[entry:symbol-registry-00589]{p.~\pageref{entry:symbol-registry-00589}} & P8.4 & A smooth total preorder on accessible worlds and the explicitly defined set of its minimal accessible B-worlds; the set is empty when there are no candidates.  \\
\(\prod\mathfrak a\) & product space of regular cardinals & P6.5 & \hyperref[entry:symbol-registry-00690]{p.~\pageref{entry:symbol-registry-00690}} & P6.5 & Functions with f(kappa) below kappa at every coordinate in the nonempty set a; the cardinality of this space is the cardinal product.  \\
\(\prod_{i\in I}\mathcal M_i/\mathcal U\) & ultraproduct & P3.3 & \hyperref[entry:symbol-registry-00047]{p.~\pageref{entry:symbol-registry-00047}} & P3.3 & Quotient of the direct product by equality on a set in the ultrafilter \(\mathcal U\).  \\
\(q_1\subseteq q_2\) & conjunctive-query containment & P9.7 & \hyperref[entry:symbol-registry-00684]{p.~\pageref{entry:symbol-registry-00684}} & P9.7 & Every answer to q1 is an answer to q2 on every finite database; distinct from set inclusion inside one instance.  \\
\(q(\bar x),D_q\) & conjunctive query and canonical database & P9.7 & \hyperref[entry:symbol-registry-00683]{p.~\pageref{entry:symbol-registry-00683}} & P9.7 & The canonical finite structure D-q carries the query variables constants atoms and distinguished free tuple.  \\
\((Q,\Sigma,\delta,q_0,\Omega)\) & deterministic parity word-automaton datum & P9.4 & \hyperref[entry:symbol-registry-00647]{p.~\pageref{entry:symbol-registry-00647}} & P9.4 & Omega assigns a finite priority to each state; the least priority occurring infinitely often determines acceptance.  \\
\(R\) & binary accessibility relation & P8.2-P8.3 & \hyperref[entry:symbol-registry-00107]{p.~\pageref{entry:symbol-registry-00107}} & P8.2 & Frame relation of modal, temporal, and epistemic structures; scope is disjoint from the P8.6 ternary relation.  \\
\(R\) & ternary accessibility relation & P8.6 & \hyperref[entry:symbol-registry-00109]{p.~\pageref{entry:symbol-registry-00109}} & P8.6 & Ternary relation in broad and matched Routley--Meyer structures; scope is disjoint from binary accessibility.  \\
\(r\doteq t\) & unification equation & P2.2-P12.3 & \hyperref[entry:symbol-registry-00173]{p.~\pageref{entry:symbol-registry-00173}} & P2.2 & Metalinguistic equation constraint processed by unification; not object-language identity.  \\
\(R_1,\ldots,R_{q+1}\) & semantic-clash intermediate resolvents & P12.3 & \hyperref[entry:symbol-registry-00619]{p.~\pageref{entry:symbol-registry-00619}} & P12.3 & Exact Slagle resolvent chain with \(R_1=N\) and final false \(M\)-resolvent \(R_{q+1}\); distinct from frame relations.  \\
\(R^2(ab)cd,R^2a(bc)d\) & composite ternary-frame relations & P8.6 & \hyperref[entry:symbol-registry-00560]{p.~\pageref{entry:symbol-registry-00560}} & P8.6 & The first two exact existential composites in the cited E R and RM frame presentation; parenthesization fixes the composition.  \\
\(R^2b(ac)d,R^2abbc\) & composite ternary-frame relations & P8.6 & \hyperref[entry:symbol-registry-00561]{p.~\pageref{entry:symbol-registry-00561}} & P8.6 & The remaining exact existential composites in the cited E R and RM frame presentation; parenthesization fixes the composition.  \\
\(r(A),h(c),o(D)\) & formula rank Cut height and derivation ordinal measure & P7.5 & \hyperref[entry:symbol-registry-00436]{p.~\pageref{entry:symbol-registry-00436}} & P7.5 & Metalinguistic measures for the finite-set propositional Cut strategy; o(D) is the Hessenberg natural sum displayed in its definition.  \\
\(R_\alpha\) & System-T finite-type recursor & P7.6 & \hyperref[entry:symbol-registry-00465]{p.~\pageref{entry:symbol-registry-00465}} & P7.6 & One primitive recursor at each finite type with argument order and computation equations fixed by sheet:system-t.  \\
\(R_e\) & computable relation decoded from e & P4.5 & \hyperref[entry:symbol-registry-00278]{p.~\pageref{entry:symbol-registry-00278}} & P4.5 & Scope disjoint from the P4.4 priority-requirement row.  \\
\(R_e\) & priority requirement against computing A from B & P4.4 & \hyperref[entry:symbol-registry-00265]{p.~\pageref{entry:symbol-registry-00265}} & P4.4 & Finite-injury requirement.  \\
\(R_f(U,V_1,\ldots,V_n)\) & canonical relation of a finitary Boolean operator & P10.3 & \hyperref[entry:symbol-registry-00733]{p.~\pageref{entry:symbol-registry-00733}} & P10.3 & Every tuple of elements from the successor ultrafilters has its operator value in U; for arity zero this is a unary predicate.  \\
\(R_G^{E},R_G^{C},R_G^{D}\) & group epistemic relations & P8.3 & \hyperref[entry:symbol-registry-00536]{p.~\pageref{entry:symbol-registry-00536}} & P8.3 & Union relation its reflexive-transitive closure and relation intersection for everyone common and distributed knowledge.  \\
\(R_i(u,v)\) & bounded configuration-reachability QBF family & P9.8 & \hyperref[entry:symbol-registry-00687]{p.~\pageref{entry:symbol-registry-00687}} & P9.8 & True exactly when v is reachable from u within two-to-the-i steps; a metalinguistic formula family rather than a new object-language predicate.  \\
\(R^{\mathcal M}_{\varphi}\) & relation defined by a formula in a structure & P2.2 & \hyperref[entry:symbol-registry-00161]{p.~\pageref{entry:symbol-registry-00161}} & P2.2 & The tuple relation cut out by \(\varphi\) in \(\mathcal M\); deliberately not \(R^*\), which is reserved for reflexive-transitive closure.  \\
\(R_n,\eta^n\) & representable finite free model and its generators & P10.7 & \hyperref[entry:symbol-registry-00811]{p.~\pageref{entry:symbol-registry-00811}} & P10.7 & The model \(R_n=T(n,-)\) is free on \(n\); its generator indexed by \(i\) is \(p_i^n:n\to1\).  \\
\(r,t\) & term metavariables & global & \hyperref[entry:symbol-registry-00004]{p.~\pageref{entry:symbol-registry-00004}} & P2.2 & Members of \(\mathrm{Term}(\mathcal L)\).  \\
\(\rho(x)\) & nonnegative rank of an ALC reduction model element & P9.3 & \hyperref[entry:symbol-registry-00701]{p.~\pageref{entry:symbol-registry-00701}} & P9.3 & Integer obtained from the concept-membership bits at x; every transition role edge decreases this value by exactly one.  \\
\(S\) & System-T successor constant & P7.6 & \hyperref[entry:symbol-registry-00464]{p.~\pageref{entry:symbol-registry-00464}} & P7.6 & Ground-type successor constant on sheet:system-t; scoped away from formula and Post-clone letters.  \\
\(s\) & variable assignment & global & \hyperref[entry:symbol-registry-00025]{p.~\pageref{entry:symbol-registry-00025}} & P2.2 & Map from first-order variables to \(\operatorname{dom}\mathcal M\).  \\
\(S_e\) & priority requirement against computing B from A & P4.4 & \hyperref[entry:symbol-registry-00266]{p.~\pageref{entry:symbol-registry-00266}} & P4.4 & Finite-injury requirement.  \\
\((S,\ell,\prec),\mu(X)\) & preferential-model datum and minimal-state operator & P8.7 & \hyperref[entry:symbol-registry-00602]{p.~\pageref{entry:symbol-registry-00602}} & P8.7 & The labeling map gives classical valuations; \(\mu(X)\) selects the \(\prec\)-minimal states of X.  \\
\(s^\flat,\varphi^{\rm alg}\) & translation between residual terms and formulas & P10.3 & \hyperref[entry:symbol-registry-00774]{p.~\pageref{entry:symbol-registry-00774}} & P10.3 & Converts algebraic zero one and residual into bottom top and implication or conversely; identities correspond to biconditionals with value one.  \\
\(s:I_\Delta\to I_\Gamma\) & interpreted simultaneous term substitution & P10.9 & \hyperref[entry:symbol-registry-00836]{p.~\pageref{entry:symbol-registry-00836}} & P10.9 & Its components interpret one substituting term per variable of \(\Gamma\); formula substitution is reindexing along this map.  \\
\(S_i,S\) & surviving ALC type sets & P9.3 & \hyperref[entry:symbol-registry-00698]{p.~\pageref{entry:symbol-registry-00698}} & P9.3 & S-zero contains all complete types; each stage removes types lacking existential witnesses; S is the stable survivor set.  \\
\(S=(\mathcal S,\mathcal A,\mathcal R);\lambda S\) & Barendregt pure-type-system specification and calculus & P12.5 & \hyperref[entry:symbol-registry-00622]{p.~\pageref{entry:symbol-registry-00622}} & P12.5 & The exact source datum consists of sorts sort-typing axioms and dependent-product triples; lambda S names the calculus it determines.  \\
\(S_n(T)\) & type space & P3 & \hyperref[entry:symbol-registry-00046]{p.~\pageref{entry:symbol-registry-00046}} & P3.2 & Complete \(n\)-types consistent with \(T\) under the clopen topology.  \\
\(S|u\) & restriction of a classified power-object element & P10.9 & \hyperref[entry:symbol-registry-00854]{p.~\pageref{entry:symbol-registry-00854}} & P10.9 & The pullback of \(S\subseteq\mathrm y c\times X\) along \(\mathrm y u\times1_X\).  \\
\(s[x\mapsto a]\) & assignment update & global & \hyperref[entry:symbol-registry-00155]{p.~\pageref{entry:symbol-registry-00155}} & P2.2 & Agrees with assignment \(s\) except that variable \(x\) is sent to \(a\).  \\
\([s],x|_n\) & binary cylinder and length-n initial segment & P10.2 & \hyperref[entry:symbol-registry-00726]{p.~\pageref{entry:symbol-registry-00726}} & P10.2 & The cylinder contains precisely the sequences extending the finite word s; the empty word gives all Cantor space.  \\
\(S_z\) & subpresheaf determined by a relation and a parameter element & P10.9 & \hyperref[entry:symbol-registry-00855]{p.~\pageref{entry:symbol-registry-00855}} & P10.9 & For \(z\in Z(c)\), its component consists of pairs \((u,x)\) with \((Z(u)(z),x)\) in the original relation.  \\
\(S4.2,S4.3,S5\) & Chellas normal modal logic identifiers & P8.2 & \hyperref[entry:symbol-registry-00496]{p.~\pageref{entry:symbol-registry-00496}} & P8.2 & Chellas's exact plain math-italic names for the displayed extensions.  \\
\(\sem A\) & IMLL algebraic interpretation & P8.6 & \hyperref[entry:symbol-registry-00571]{p.~\pageref{entry:symbol-registry-00571}} & P8.6 & The inherited source's semantic value of formula A in the declared commutative residuated partially ordered monoid.  \\
\(\Sigma^0_n,\Pi^0_n\) & syntactic arithmetic formula classes & P4-P7.5 & \hyperref[entry:symbol-registry-00017]{p.~\pageref{entry:symbol-registry-00017}} & P4.3 & Classes fixed by the declared normalization and unbounded-quantifier alternation; P7.5 uses the same hierarchy in induction-fragment names.  \\
\(\Sigma^1_n,\Pi^1_n,\Delta^1_n,\boldsymbol\Sigma^1_n,\boldsymbol\Pi^1_n\) & lightface and boldface analytical hierarchy classes & P2.4-P4.5 & \hyperref[entry:symbol-registry-00224]{p.~\pageref{entry:symbol-registry-00224}} & P2.4 & Lightface classes have effective matrices and \(\Delta^1_n\) is their intersection; boldface classes allow real or set parameters and must not be identified with them.  \\
\(\Sigma=\bigsqcup_{n\geq0}\Sigma_n,T_\Sigma\) & finite ranked alphabet and finite tree set & P9.4 & \hyperref[entry:symbol-registry-00650]{p.~\pageref{entry:symbol-registry-00650}} & P9.4 & Every symbol in Sigma-n has arity n and T-Sigma is the least finite-term set generated according to those ranks.  \\
\(\Sigma_k\text{-}\mathsf{SAT}\) & existential-start fixed-alternation QBF language & P4.6 & \hyperref[entry:symbol-registry-00309]{p.~\pageref{entry:symbol-registry-00309}} & P4.6 & Fixed \(k\).  \\
\(\Sigma_{\mathcal T},E_{\mathcal T},\sigma_f\) & named equational presentation of a theory & P10.7 & \hyperref[entry:symbol-registry-00807]{p.~\pageref{entry:symbol-registry-00807}} & P10.7 & One operation symbol for each arrow to the generator; projection and substitution equations recover precisely the original arrows.  \\
\(\Sigma_n^{N\mu},\Pi_n^{N\mu};\Sigma_n^{S\mu},\Pi_n^{S\mu};\Sigma_n^{EL\mu},\Pi_n^{EL\mu}\) & Niwiński simple and Emerson--Lei modal-mu alternation classes & P9.5 & \hyperref[entry:symbol-registry-00655]{p.~\pageref{entry:symbol-registry-00655}} & P9.5 & Source-specific positive-form syntactic classes; the variants differ by capture-avoiding substitution closure, closed substitution only, or no substitution closure.  \\
\(\Sigma^p_k\) & existential-start polynomial-hierarchy class & P4.6 & \hyperref[entry:symbol-registry-00305]{p.~\pageref{entry:symbol-registry-00305}} & P4.6 & Bounded alternating string quantifiers.  \\
\(\sigma,\tau\) & sentence metavariables & global & \hyperref[entry:symbol-registry-00010]{p.~\pageref{entry:symbol-registry-00010}} & P2.2 & Members of \(\mathrm{Sent}(\mathcal L)\).  \\
\(\sigma,\tau:\Delta\to\Gamma\) & MLTT substitutions with common domain and codomain & P7.4 & \hyperref[entry:symbol-registry-00430]{p.~\pageref{entry:symbol-registry-00430}} & P7.4 & Capture-avoiding substitutions used by typed constructor congruence; equality premises and conclusions live in Delta.  \\
\(\sigma:\Theta\to\Psi,M[\sigma]\) & typed simultaneous substitution and its action & P10.8 & \hyperref[entry:symbol-registry-00786]{p.~\pageref{entry:symbol-registry-00786}} & P10.8 & A tuple of Theta-terms with the types declared in Psi; substitutes into a Psi-term without resubstituting into inserted terms.  \\
\(\sim_{\Delta},[t],\mathcal M_{\Delta},s_0\) & Henkin quotient-model notation & P2.2 & \hyperref[entry:symbol-registry-00176]{p.~\pageref{entry:symbol-registry-00176}} & P2.2 & Provable-equality relation term class quotient term structure and canonical variable assignment for a maximal consistent witness set \(\Delta\).  \\
\(\sim_i\) & epistemic accessibility & P8.3 & \hyperref[entry:symbol-registry-00099]{p.~\pageref{entry:symbol-registry-00099}} & P8.3 & Agent-indexed accessibility; an equivalence relation for knowledge readings.  \\
\(\sim_T\) & Turing equivalence & P4.4 & \hyperref[entry:symbol-registry-00259]{p.~\pageref{entry:symbol-registry-00259}} & P4.4 & Mutual Turing reducibility.  \\
\(\simeq\) & Kleene equality & P4.2-P4.5 & \hyperref[entry:symbol-registry-00049]{p.~\pageref{entry:symbol-registry-00049}} & P4.2 & Both sides are defined with equal value or both are undefined.  \\
\(\sqcup\) & blank tape symbol & P4.1-P4.2 & \hyperref[entry:symbol-registry-00231]{p.~\pageref{entry:symbol-registry-00231}} & P4.1 & Distinguished tape-alphabet symbol.  \\
\(\square_\kappa\) & square principle & P6.4 & \hyperref[entry:symbol-registry-00371]{p.~\pageref{entry:symbol-registry-00371}} & P6.4 & Coherent club sequence with the explicit no-thread condition; distinct from the modal necessity glyph.  \\
\(\star,\langle M,N\rangle,\pi_iM,\mathsf{inl}(M),\mathsf{inr}(M)\) & unit product and sum constructors for the categorical STLC extension & P7.3 & \hyperref[entry:symbol-registry-00434]{p.~\pageref{entry:symbol-registry-00434}} & P7.3 & Term constructors governed by the exact typing table and beta-eta equations in the Curry-Howard entry.  \\
\(\sup A\) & ordinal supremum & P6-P7.5 & \hyperref[entry:symbol-registry-00382]{p.~\pageref{entry:symbol-registry-00382}} & P6.1 & Union of a set of ordinals and the least ordinal at least every member; extended to the P7.5 ordinal towers.  \\
\(t\leadsto_r u\) & role compatibility of ALC types & P9.3 & \hyperref[entry:symbol-registry-00697]{p.~\pageref{entry:symbol-registry-00697}} & P9.3 & Every universal r-restriction in t has its filler in u; this relation supplies all role edges in the canonical model.  \\
\(t\theta,\theta\eta\) & simultaneous substitution action and composition & P2.2-P12.3 & \hyperref[entry:symbol-registry-00171]{p.~\pageref{entry:symbol-registry-00171}} & P2.2 & Action extends homomorphically and composition satisfies \((\theta\eta)(x)=(\theta(x))\eta\).  \\
\(t\sigma,\widehat\sigma\) & simultaneous algebraic term substitution & P10.1 & \hyperref[entry:symbol-registry-00705]{p.~\pageref{entry:symbol-registry-00705}} & P10.1 & The map on variables extends homomorphically to terms; the extension may have a different generator alphabet as codomain.  \\
\(t^A(v)\) & evaluation of an algebraic term & P10.1 & \hyperref[entry:symbol-registry-00704]{p.~\pageref{entry:symbol-registry-00704}} & P10.1 & Term value in A under the variable assignment v; every assignment extends uniquely to a term-algebra homomorphism.  \\
\(T(\bar P,\bar Q,\bar Z),\bar P'\leq\bar P\) & circumscription theory and predicate-tuple inclusion & P8.7 & \hyperref[entry:symbol-registry-00607]{p.~\pageref{entry:symbol-registry-00607}} & P8.7 & P is minimized Q is fixed and Z may vary; tuple inclusion abbreviates componentwise universal closure.  \\
\(T_{\mathsf{CC}}^\sharp\) & translated matched-schema theory & P2.3 & \hyperref[entry:symbol-registry-00197]{p.~\pageref{entry:symbol-registry-00197}} & P2.3 & The many-sorted first-order translations of exactly all eligible comprehension and graph-choice instances.  \\
\(t,\operatorname{tp}_I(x)\) & complete ALC concept type and realized type & P9.3 & \hyperref[entry:symbol-registry-00696]{p.~\pageref{entry:symbol-registry-00696}} & P9.3 & A type is a subset of the fixed closure satisfying Boolean completeness and every internalized GCI; the realized type records truth at x in I.  \\
\(T_{\Sigma,E}(n,m)\) & term theory of an equational presentation & P10.7 & \hyperref[entry:symbol-registry-00802]{p.~\pageref{entry:symbol-registry-00802}} & P10.7 & Its arrows are \(m\)-tuples of congruence classes of \(n\)-variable terms; composition is substitution and object \(0\) is terminal.  \\
\(T_\tau(X)\) & term algebra on a generator set & P10.1 & \hyperref[entry:symbol-registry-00703]{p.~\pageref{entry:symbol-registry-00703}} & P10.1 & Finite terms over the finitary signature tau and variables X; a nonempty algebra only when variables or nullary symbols exist.  \\
\(T,U\) & theories & global & \hyperref[entry:symbol-registry-00021]{p.~\pageref{entry:symbol-registry-00021}} & P2.2 & Sets of sentences in the fixed signature.  \\
\(\tau_{\mathrm{fin}}\) & fixed finite relational signature & P9.1 & \hyperref[entry:symbol-registry-00635]{p.~\pageref{entry:symbol-registry-00635}} & P9.1 & The signature relative to which finite satisfiability and finite validity are defined; it is not the varying database signature of a combined-complexity input.  \\
\(\text{\L}\) & standard infinite-valued Lukasiewicz logic identifier & P8.5 & \hyperref[entry:symbol-registry-00544]{p.~\pageref{entry:symbol-registry-00544}} & P8.5 & Hajek's exact stroked-L identifier for the infinite-valued logic; its standard value algebra is the unit interval.  \\
\(\text{\L}_n\) & finite Lukasiewicz chain & P8.5 & \hyperref[entry:symbol-registry-00543]{p.~\pageref{entry:symbol-registry-00543}} & P8.5 & Cignoli--D'Ottaviano--Mundici's stroked-L name for the n-element chain itself; the corresponding n-valued system is described in source-bound prose.  \\
\(\Theta\vdash_{\beta\eta}M=N:\alpha\) & typed beta-eta equational consequence & P10.8 & \hyperref[entry:symbol-registry-00785]{p.~\pageref{entry:symbol-registry-00785}} & P10.8 & The equational metatheory includes alpha-renaming and typed beta and eta schemas; the registered reduction relation remains beta reduction.  \\
\(\theta\preceq_{\mathrm{inst}}\rho\) & substitution instance preorder & P2.2-P12.3 & \hyperref[entry:symbol-registry-00172]{p.~\pageref{entry:symbol-registry-00172}} & P2.2 & Every instance of \(\rho\) is an instance of \(\theta\); equivalently \(\rho=\theta\eta\) for some substitution \(\eta\).  \\
\(\theta_{a,b}\) & maximal congruence separating a chosen pair & P10.1 & \hyperref[entry:symbol-registry-00716]{p.~\pageref{entry:symbol-registry-00716}} & P10.1 & Maximal among congruences omitting (a,b); the quotient is subdirectly irreducible and still separates the chosen elements.  \\
\(\theta_\Gamma\) & congruence of provable identities & P10.1 & \hyperref[entry:symbol-registry-00708]{p.~\pageref{entry:symbol-registry-00708}} & P10.1 & The least fully invariant congruence on the nonempty term algebra containing all premise equation pairs.  \\
\(\theta_{\mathcal V,X}\) & semantic congruence defining a relative free algebra & P10.1 & \hyperref[entry:symbol-registry-00710]{p.~\pageref{entry:symbol-registry-00710}} & P10.1 & Two X-terms are related exactly when all evaluations into all algebras of the variety identify them.  \\
\(\Theta_{\mathrm{ty}}\) & term-variable typing context & P7.3-P7.6 & \hyperref[entry:symbol-registry-00403]{p.~\pageref{entry:symbol-registry-00403}} & P7.3 & Finite list of distinct term-variable declarations carrying well-formed types in the active typed calculus.  \\
\(\Theta_{\mathrm{ty}}\) & term-variable typing context & P12.5 & \hyperref[entry:symbol-registry-00632]{p.~\pageref{entry:symbol-registry-00632}} & P12.5 & Finite list of distinct term-variable declarations in the reused STLC example.  \\
\(\Theta_{\mathrm{ty}}\derives_{\mathsf{STLC}}M{:}\alpha\) & STLC typing judgment & P7.3-P7.4 & \hyperref[entry:symbol-registry-00405]{p.~\pageref{entry:symbol-registry-00405}} & P7.3 & The exact arrow-only typing judgment of sheet:stlc; the context contains distinct term-variable declarations.  \\
\(\Theta_{\mathrm{ty}}\derives_{\mathsf{STLC}}M{:}\alpha\) & STLC typing judgment & P12.5 & \hyperref[entry:symbol-registry-00634]{p.~\pageref{entry:symbol-registry-00634}} & P12.5 & Exact reused arrow-only typing judgment; no new P12 calculus is introduced.  \\
\(\Theta_{\mathrm{ty}}\derives_{\mathsf T_{\mathrm{tm}}}M{:}\alpha\) & System-T typing judgment & P7.6 & \hyperref[entry:symbol-registry-00462]{p.~\pageref{entry:symbol-registry-00462}} & P7.6 & Typing judgment of sheet:system-t with the exact finite-type term calculus in the turnstile subscript.  \\
\(\theta,\rho,\eta\) & simultaneous term substitutions & P2.2-P12.3 & \hyperref[entry:symbol-registry-00013]{p.~\pageref{entry:symbol-registry-00013}} & P2.2 & Finite-support variable-to-term maps with declared action and composition.  \\
\(\theta_X,Q_X\) & semantic term congruence and free equational algebra & P10.7 & \hyperref[entry:symbol-registry-00812]{p.~\pageref{entry:symbol-registry-00812}} & P10.7 & The congruence requires equality under every evaluation into an all-carrier model; its quotient includes the empty term case.  \\
\(U\Rightarrow_X V\) & Boolean implication between subsets & P10.9 & \hyperref[entry:symbol-registry-00825]{p.~\pageref{entry:symbol-registry-00825}} & P10.9 & The subset \((X\setminus U)\cup V\) is the Heyting residual in the powerset fibre.  \\
\(U\,A\,B,S\,A\,B\) & Kamp strict tense connectives & P8.3 & \hyperref[entry:symbol-registry-00537]{p.~\pageref{entry:symbol-registry-00537}} & P8.3 & Kamp's prefix source syntax: the first operand holds at the future or past witness and the second at every strictly intervening point.  \\
\(U_C\) & optimal plain decompressor & P4.5 & \hyperref[entry:symbol-registry-00286]{p.~\pageref{entry:symbol-registry-00286}} & P4.5 & Distinguished from theory metavariables.  \\
\(U(e,x)\) & universal partial computable function & P4.2-P4.4 & \hyperref[entry:symbol-registry-00243]{p.~\pageref{entry:symbol-registry-00243}} & P4.2 & Exact source notation; oracle superscript \(U^B\) gives relativization.  \\
\(U_K\) & optimal prefix-free machine & P4.5 & \hyperref[entry:symbol-registry-00287]{p.~\pageref{entry:symbol-registry-00287}} & P4.5 & Used throughout the Chaitin entries instead of bare \(U\).  \\
\(U_p\) & principal ultrafilter of a Boolean atom & P10.2 & \hyperref[entry:symbol-registry-00729]{p.~\pageref{entry:symbol-registry-00729}} & P10.2 & All elements above atom p; its singleton is the basic clopen hat-p and hence an isolated Stone point.  \\
\(u,v,w\) & finite words & P4.1 & \hyperref[entry:symbol-registry-00016]{p.~\pageref{entry:symbol-registry-00016}} & P4.1 & Finite words over the declared alphabet; later P4 chapters use the letters numerically.  \\
\(\ulcorner\varphi\urcorner\) & Goedel number & P5 & \hyperref[entry:symbol-registry-00055]{p.~\pageref{entry:symbol-registry-00055}} & P5.1 & Corner quotes around an expression denote its fixed effective code.  \\
\(\underline A,A_T,P_\Gamma\) & named type interpreted type and context object & P10.8 & \hyperref[entry:symbol-registry-00789]{p.~\pageref{entry:symbol-registry-00789}} & P10.8 & The named base type denotes \(A\); \(A_T\) interprets the type \(T\), and \(P_\Gamma\) is the recursively chosen product with leading terminal factor.  \\
\(\upsilon\) & uniform propositional substitution & P2.1 & \hyperref[entry:symbol-registry-00130]{p.~\pageref{entry:symbol-registry-00130}} & P2.1 & A metalinguistic map \(\mathsf{Prop}\to\mathsf{Fm}_{\mathsf{CPL}}\); not the globally registered sentence metavariables \(\sigma,\tau\).  \\
\(V\) & ambient universe of sets & P6 & \hyperref[entry:symbol-registry-00381]{p.~\pageref{entry:symbol-registry-00381}} & P6.1 & Metalinguistic class abbreviation for the union of the cumulative hierarchy; not a set or a valuation.  \\
\(v\) & valuation & global & \hyperref[entry:symbol-registry-00024]{p.~\pageref{entry:symbol-registry-00024}} & P2.1 & Map from propositional variables to the active truth-value set in propositional and finite-matrix semantics; its homomorphic extension \(\widehat v\) is fixed by the active semantics.  \\
\(V\) & world valuation & P8 & \hyperref[entry:symbol-registry-00120]{p.~\pageref{entry:symbol-registry-00120}} & P8.1 & \(V:\mathsf{Prop}\to\mathcal P(W)\) (temporal structures: \(V:\mathbb N\to\mathcal P(\mathsf{Prop})\)); distinct from matrix \(\mathsf V\) and cumulative \(V_\alpha\).  \\
\(v\models\varphi\) & valuation satisfaction & P2.1 & \hyperref[entry:symbol-registry-00121]{p.~\pageref{entry:symbol-registry-00121}} & P2.1 & The valuation's homomorphic extension designates \(\varphi\).  \\
\(v_0(p)=\widehat{[p]_L}\) & canonical modal valuation & P10.3 & \hyperref[entry:symbol-registry-00755]{p.~\pageref{entry:symbol-registry-00755}} & P10.3 & Assigns each proposition variable the basic clopen determined by its formula class.  \\
\(V_\alpha\) & cumulative-hierarchy level & P6 & \hyperref[entry:symbol-registry-00065]{p.~\pageref{entry:symbol-registry-00065}} & P6.1 & \(V=\bigcup_\alpha V_\alpha\) is class notation, registered with its level notation.  \\
\(V^{\mathbb B}_\alpha,V^{\mathbb B}\) & ambient Boolean-valued name hierarchy & P6.3 & \hyperref[entry:symbol-registry-00322]{p.~\pageref{entry:symbol-registry-00322}} & P6.3 & Names are Boolean-valued functions on earlier levels for an ambient-complete algebra; distinct from the P-name hierarchy and the ground-model-computed variant.  \\
\((V^{\mathbb B})^M,\operatorname{Bool}_{\mathbb B}^M,S_{\varphi,\bar{\mathbf v}}^M\) & ground-model-computed Boolean names truth and attained-value set & P6.3 & \hyperref[entry:symbol-registry-00324]{p.~\pageref{entry:symbol-registry-00324}} & P6.3 & All recursions, Separation, and joins are computed in M when M regards B as complete; ambient completeness is not assumed.  \\
\(V^{\mathbb P}\) & class of P-names & P6.3 & \hyperref[entry:symbol-registry-00069]{p.~\pageref{entry:symbol-registry-00069}} & P6.3 & Name metavariables \(\nu\) range over \(M\cap V^{\mathbb P}\).  \\
\(V(p,x)\) & Davies--Humberstone atomic valuation & P8.2 & \hyperref[entry:symbol-registry-00511]{p.~\pageref{entry:symbol-registry-00511}} & P8.2 & Exact source two-place truth-value notation for atom p at world x inside the model pair; locally typed apart from set-valued world valuations.  \\
\(\varepsilon_0\) & least epsilon ordinal & P5-P7 & \hyperref[entry:symbol-registry-00078]{p.~\pageref{entry:symbol-registry-00078}} & P5.4 & Least fixed point of ordinal exponentiation base \(\omega\); calibrates PA in P7.5.  \\
\(\varepsilon_{\mathcal T},\rho\) & isomorphism recovering a Lawvere theory & P10.7 & \hyperref[entry:symbol-registry-00809]{p.~\pageref{entry:symbol-registry-00809}} & P10.7 & Evaluation from the named term theory is identity on objects; \(\rho\) denotes this fixed isomorphism when transporting free models.  \\
\(\varepsilon_X\) & Stone-space reconstruction homeomorphism & P10.2 & \hyperref[entry:symbol-registry-00721]{p.~\pageref{entry:symbol-registry-00721}} & P10.2 & Maps x to the ultrafilter of all clopen subsets containing x.  \\
\(\varepsilon x\,\varphi(x)\) & Hilbert epsilon term former & P7.7 & \hyperref[entry:symbol-registry-00475]{p.~\pageref{entry:symbol-registry-00475}} & P7.7 & Capture-avoiding term binder on the equality-free epsilon sheet; distinct from epsilon-zero.  \\
\(\varepsilon_X:X\to X_{\operatorname{ClopUp}(X)}\) & ordered-space reconstruction by prime filters & P10.3 & \hyperref[entry:symbol-registry-00759]{p.~\pageref{entry:symbol-registry-00759}} & P10.3 & Maps a point to its membership filter of clopen upsets; an order isomorphism and homeomorphism.  \\
\(\varnothing^{(n)}\) & iterated jump of the empty set & P4.4-P5 & \hyperref[entry:symbol-registry-00054]{p.~\pageref{entry:symbol-registry-00054}} & P4.4 & First defined with the general jump iteration.  \\
\(\varphi\Leftrightarrow\psi\) & formula equivalence & global & \hyperref[entry:symbol-registry-00030]{p.~\pageref{entry:symbol-registry-00030}} & P2.1 & Metalinguistic biconditionals are written in prose as iff; \(\Leftrightarrow\) relates formulas and \(\equiv\) is reserved for elementary equivalence.  \\
\(\varphi\leftrightarrow\psi\) & object-language biconditional & global & \hyperref[entry:symbol-registry-00031]{p.~\pageref{entry:symbol-registry-00031}} & P2.1 & Object-language constructor; rank 7 of the extended precedence ladder.  \\
\(\varphi\land\psi\) & object-language conjunction & global & \hyperref[entry:symbol-registry-00007]{p.~\pageref{entry:symbol-registry-00007}} & P2.1 & Primitive or explicitly defined binary connective in the active language; left associative in handbook prose.  \\
\(\varphi\lor\psi\) & object-language disjunction & global & \hyperref[entry:symbol-registry-00008]{p.~\pageref{entry:symbol-registry-00008}} & P2.1 & Primitive or explicitly defined binary connective in the active language; left associative in handbook prose.  \\
\(\varphi\to\psi\) & object-language implication & global & \hyperref[entry:symbol-registry-00009]{p.~\pageref{entry:symbol-registry-00009}} & P2.1 & Primitive or explicitly defined binary connective in the active language; right associative in handbook prose.  \\
\(\varphi\equiv_L\psi,[\varphi]_L\) & normal-modal provable equivalence and its classes & P10.3 & \hyperref[entry:symbol-registry-00741]{p.~\pageref{entry:symbol-registry-00741}} & P10.3 & Equivalence means theoremhood of the biconditional in L and is a congruence for all formula constructors.  \\
\(\varphi\mathbin{\&}\psi,\varphi\Rightarrow\psi,\mathsf{false}\) & Pitts's propositional formula constructors & P10.9 & \hyperref[entry:symbol-registry-00831]{p.~\pageref{entry:symbol-registry-00831}} & P10.9 & Conjunction, implication, and falsity at the formula level; their denotations are the corresponding Heyting operations.  \\
\(\varphi=_{\alpha}\psi\) & alpha-equivalence & global & \hyperref[entry:symbol-registry-00154]{p.~\pageref{entry:symbol-registry-00154}} & P2.2 & Least constructor congruence generated by capture-avoiding bound-variable renaming; distinct from identity formula equivalence and elementary equivalence.  \\
\(\varphi_e\) & partial computable function with index e & P4-P5 & \hyperref[entry:symbol-registry-00052]{p.~\pageref{entry:symbol-registry-00052}} & P4.2 & Oracle form \(\varphi^B_e\) is relative to \(B\); finite-string form \(\varphi_e^\tau\) and stage form \(\varphi^{B_s}_{e,s}\) are defined in the two priority proofs.  \\
\(\varphi^k\) & repeated Mix-formula occurrence & P7.1-P7.2 & \hyperref[entry:symbol-registry-00397]{p.~\pageref{entry:symbol-registry-00397}} & P7.1 & Exactly k comma-separated occurrences on the registered Mix sheets; not exponentiation or iteration.  \\
\(\varphi^{\mathrm D}\) & Dialectica translation of a formula & P7.6 & \hyperref[entry:symbol-registry-00467]{p.~\pageref{entry:symbol-registry-00467}} & P7.6 & Translation into an existential witness block, universal challenge block, and quantifier-free target matrix.  \\
\(\varphi_{\mathrm D}(\vec x,\vec y)\) & Dialectica quantifier-free matrix & P7.6 & \hyperref[entry:symbol-registry-00468]{p.~\pageref{entry:symbol-registry-00468}} & P7.6 & Matrix associated with the locally fixed Dialectica translation; typed witness and challenge tuples remain explicit.  \\
\(\varphi,\psi,\chi\) & formula metavariables & global & \hyperref[entry:symbol-registry-00005]{p.~\pageref{entry:symbol-registry-00005}} & P2.1 & Members of the active formula set.  A source-bound presentation may instead declare its own formula metavariables locally; the source declaration governs that presentation.  \\
\(\varphi[t/x]\) & capture-avoiding substitution & global & \hyperref[entry:symbol-registry-00033]{p.~\pageref{entry:symbol-registry-00033}} & P2.2 & Substitution of \(t\) for free occurrences of \(x\) in \(\varphi\).  \\
\(\varphi[U^n/X^n]\) & capture-avoiding relation-variable substitution & P2.3 & \hyperref[entry:symbol-registry-00187]{p.~\pageref{entry:symbol-registry-00187}} & P2.3 & Metalinguistic substitution of a relation variable or predicate constant of the same arity; distinct from term substitution.  \\
\(\varphi[\upsilon]\) & uniform-substitution instance & P2.1 & \hyperref[entry:symbol-registry-00131]{p.~\pageref{entry:symbol-registry-00131}} & P2.1 & Recursive action of the declared uniform substitution on a CPL formula.  \\
\(\vec a,\vec n\) & finite tuples of natural-number parameters & P4-P5 & \hyperref[entry:symbol-registry-00227]{p.~\pageref{entry:symbol-registry-00227}} & P4.1 & Vector arrows distinguish tuples from numeral bars.  \\
\(\vec f,\vec g,\vec t\) & finite tuples of System-T functionals and terms & P7.6 & \hyperref[entry:symbol-registry-00470]{p.~\pageref{entry:symbol-registry-00470}} & P7.6 & Typed functional witnesses in translation clauses and the extracted closed-term tuple in Dialectica soundness.  \\
\(\vec x,\vec y,\vec u,\vec v\) & Dialectica witness and challenge tuples & P7.6 & \hyperref[entry:symbol-registry-00469]{p.~\pageref{entry:symbol-registry-00469}} & P7.6 & Finite tuples of finite-type variables; distinct from the P4-P5 tuples of numerical variables.  \\
\(\vec x,\vec y,\vec z,\vec u,\vec v\) & finite tuples of numerical variables & P4-P5 & \hyperref[entry:symbol-registry-00228]{p.~\pageref{entry:symbol-registry-00228}} & P4.1 & Vector arrows distinguish variable blocks from numeral bars.  \\
\(w^*\) & Routley star & P8.6 & \hyperref[entry:symbol-registry-00110]{p.~\pageref{entry:symbol-registry-00110}} & P8.6 & Postfix unary operation on relevance worlds used by the negation clause; distinct from closure stars and the extension marker in \(\mathsf L^*\).  \\
\(W\) & set of worlds & P8 & \hyperref[entry:symbol-registry-00106]{p.~\pageref{entry:symbol-registry-00106}} & P8.1 & Nonempty world carrier for intuitionistic, modal, temporal, epistemic, and broad Routley--Meyer semantics; world letters are declared locally.  \\
\(w\Vdash A\) & Kripke forcing & P8.1 & \hyperref[entry:symbol-registry-00090]{p.~\pageref{entry:symbol-registry-00090}} & P8.1 & The negated partner is \(\nVdash\) (macro \texttt{\detokenize{\nforces}}); modal satisfaction continues to use \(\models\).  \\
\((W,D),\Gamma_E,U_k\) & default theory grounded closure and stages & P8.7 & \hyperref[entry:symbol-registry-00606]{p.~\pageref{entry:symbol-registry-00606}} & P8.7 & Gamma-E is the least closure grounded through the stages U-k while justification consistency is tested against candidate E.  \\
\(W_e\) & c.e. set induced by program index e & P4.4 & \hyperref[entry:symbol-registry-00267]{p.~\pageref{entry:symbol-registry-00267}} & P4.4 & Domain under the acceptable numbering.  \\
\(W_i,R,P,\operatorname{rd}(M,F)\) & state copies reachable overapproximation interpolant and reverse depth & P9.5 & \hyperref[entry:symbol-registry-00657]{p.~\pageref{entry:symbol-registry-00657}} & P9.5 & W-i are renamed state blocks R overapproximates reachable states P is the one-state interpolant and rd is the maximum finite shortest distance to the bad set.  \\
\((W,\mathcal F,\mu,\pi),w(\varphi)\) & atemporal probability structure and likelihood term & P9.5 & \hyperref[entry:symbol-registry-00665]{p.~\pageref{entry:symbol-registry-00665}} & P9.5 & Mu is the event measure and w-phi its value on the phi-event; pi is locally the world-valuation map and not the temporal path metavariable.  \\
\(w{:}\varphi\) & world-labelled formula & P7.7 & \hyperref[entry:symbol-registry-00476]{p.~\pageref{entry:symbol-registry-00476}} & P7.7 & Labelled-sequent syntax only; distinct from type ascription and not a semantic satisfaction assertion.  \\
\(\widehat{\mathcal C}\) & presheaf category on a small base & P10.9 & \hyperref[entry:symbol-registry-00841]{p.~\pageref{entry:symbol-registry-00841}} & P10.9 & The category of functors \(\mathcal C^{\mathrm{op}}\to\mathsf{Set}\) and natural transformations, under the declared size convention.  \\
\(\widehat y,b(y)\) & ALC retained-edge target and direct blocker & P9.3 & \hyperref[entry:symbol-registry-00644]{p.~\pageref{entry:symbol-registry-00644}} & P9.3 & Defined only for a target y of an edge retained from an unblocked source; widehat-y is y when unblocked and otherwise a directly blocking unblocked ancestor b(y).  \\
\(wRv\) & relational atom in labelled-sequent syntax & P7.7 & \hyperref[entry:symbol-registry-00477]{p.~\pageref{entry:symbol-registry-00477}} & P7.7 & Syntactic frame-relation atom inside a labelled derivation language; distinct from semantic accessibility in P8.  \\
\(|X|\) & cardinality & global & \hyperref[entry:symbol-registry-00041]{p.~\pageref{entry:symbol-registry-00041}} & P1.3 & Vertical bars mean cardinality only; domains use \(\operatorname{dom}\).  \\
\(X^*\) & finite iteration or closure & global & \hyperref[entry:symbol-registry-00037]{p.~\pageref{entry:symbol-registry-00037}} & P1.3 & For an alphabet this is finite-word Kleene star; for a relation it is reflexive-transitive closure. Operand type fixes the reading.  \\
\(x\mathbin{\ast_{\text{\L}}}y,x\mathbin{\ast_{\Pi}}y\) & Lukasiewicz and Product t-norms & P8.5 & \hyperref[entry:symbol-registry-00553]{p.~\pageref{entry:symbol-registry-00553}} & P8.5 & Hajek's source subscripts use stroked L for the Lukasiewicz max-zero-sum-minus-one t-norm and Pi for the Product multiplication t-norm.  \\
\(X\mathbin{\upharpoonright}Y\) & restriction or reduct operator & global & \hyperref[entry:symbol-registry-00157]{p.~\pageref{entry:symbol-registry-00157}} & P2.2 & Operand types distinguish function or set restriction, a structure reduct, and the P8.3 public-announcement submodel induced by the surviving worlds.  \\
\(x\ast y\) & t-norm operation & P8.5 & \hyperref[entry:symbol-registry-00554]{p.~\pageref{entry:symbol-registry-00554}} & P8.5 & Hajek's generic associative commutative monotone binary operation on the unit interval with unit one.  \\
\(x\Rightarrow y\) & t-norm residuum and Lukasiewicz implication & P8.5 & \hyperref[entry:symbol-registry-00552]{p.~\pageref{entry:symbol-registry-00552}} & P8.5 & Hajek's plain source glyph.  For the Lukasiewicz t-norm it is min of one and one-minus-x-plus-y; numeric operands type it as a truth operation rather than a sequent arrow.  \\
\(X_B,\widehat a\) & Stone space and its basic clopen sets & P10.2 & \hyperref[entry:symbol-registry-00718]{p.~\pageref{entry:symbol-registry-00718}} & P10.2 & X-B consists of the ultrafilters of B; hat-a consists of those containing a; trivial B gives the empty space.  \\
\(X_D,\widehat a\) & Priestley spectrum and its basic clopen upsets & P10.3 & \hyperref[entry:symbol-registry-00756]{p.~\pageref{entry:symbol-registry-00756}} & P10.3 & Proper prime lattice filters ordered by inclusion; hat-a consists of those containing a, and the topology also uses their complements.  \\
\((x)e,\varphi(M)\) & metalevel formula abstraction and instantiation & P10.9 & \hyperref[entry:symbol-registry-00833]{p.~\pageref{entry:symbol-registry-00833}} & P10.9 & If \(\varphi=(x)e\), then \(\varphi(M)=e[M/x]\); this is binding notation, not an additional first-order connective.  \\
\([X]^\ell\) & ell-element subsets of a set & P6.4-P6.6 & \hyperref[entry:symbol-registry-00364]{p.~\pageref{entry:symbol-registry-00364}} & P6.4 & Domain of finite-color partition relations; ell is a positive finite integer.  \\
\(X(h)\) & continuous dual of a Boolean homomorphism & P10.2 & \hyperref[entry:symbol-registry-00722]{p.~\pageref{entry:symbol-registry-00722}} & P10.2 & For h from B to C sends an ultrafilter of C to its inverse image under h, an ultrafilter of B.  \\
\(X^n,Y^n,U^n,F^{n+1}\) & relation-variable metavariables & P2.3 & \hyperref[entry:symbol-registry-00185]{p.~\pageref{entry:symbol-registry-00185}} & P2.3 & Superscripts record arity; Y and F carry freshness conditions and U is an admissible relation-variable or predicate-constant instance.  \\
\(X^\perp\) & phase orthogonal of a subset & P8.6 & \hyperref[entry:symbol-registry-00569]{p.~\pageref{entry:symbol-registry-00569}} & P8.6 & Girard's source notation for subset orthogonality relative to the phase-space pole.  Subset operands distinguish it from formula linear negation without altering the glyph.  \\
\(X(u),\alpha_c\) & presheaf restriction and natural-transformation component & P10.9 & \hyperref[entry:symbol-registry-00843]{p.~\pageref{entry:symbol-registry-00843}} & P10.9 & For \(u:d\to c\), restriction maps \(X(c)\) to \(X(d)\); the component \(\alpha_c\) maps \(X(c)\) to \(Y(c)\).  \\
\(x^+,x^-\) & positive and negative support coordinates & P8.5 & \hyperref[entry:symbol-registry-00548]{p.~\pageref{entry:symbol-registry-00548}} & P8.5 & Boolean support indicators of an FDE value; the two coordinates are required to distinguish both from neither.  \\
\((x)(x')e,\psi(M,N)\) & binary formula abstraction and simultaneous instantiation & P10.9 & \hyperref[entry:symbol-registry-00840]{p.~\pageref{entry:symbol-registry-00840}} & P10.9 & Application replaces the two displayed argument variables simultaneously in the body, avoiding capture.  \\
\(X,Y\) & phase subsets and facts & P8.6 & \hyperref[entry:symbol-registry-00577]{p.~\pageref{entry:symbol-registry-00577}} & P8.6 & Girard uses the same source metavariables for arbitrary phase subsets and for biorthogonally closed facts; the locally stated closure property supplies the fact role.  \\
\(x,y,z\) & object-language variables & global & \hyperref[entry:symbol-registry-00003]{p.~\pageref{entry:symbol-registry-00003}} & P2.2 & Variables occurring in terms and formulas.  \\
\(\xRightarrow{\theta}\) & SLD transition labelled by an mgu & P9.7 & \hyperref[entry:symbol-registry-00678]{p.~\pageref{entry:symbol-registry-00678}} & P9.7 & Applies the displayed most general unifier to the entire resolvent on sheet:sld-resolution.  \\
\((Y^X)(c)\) & component of a presheaf exponential & P10.9 & \hyperref[entry:symbol-registry-00848]{p.~\pageref{entry:symbol-registry-00848}} & P10.9 & The set \(\operatorname{Nat}(\mathrm y c\times X,Y)\), with restriction by precomposition; it need not be the pointwise external function set.  \\
\(Z\) & modal bisimulation relation & P8.2 & \hyperref[entry:symbol-registry-00503]{p.~\pageref{entry:symbol-registry-00503}} & P8.2 & Relation between worlds of two pointed Kripke structures satisfying atom zig and zag clauses; disjoint from P8.3 fixpoint variables.  \\
\(Z,Y,\mu Z.A,\nu Z.A\) & modal fixed-point variables and binders & P8.3-P9.5 & \hyperref[entry:symbol-registry-00528]{p.~\pageref{entry:symbol-registry-00528}} & P8.3 & Z and Y are Kozen's fixed-point variables; mu is the primitive least-fixed-point binder and nu is its capture-avoiding defined dual in the exact L-mu grammar.  \\
\(\zeta,\operatorname{Ext}_{\zeta}(A)\) & modal fixed-point environment and extension map & P8.3-P9.5 & \hyperref[entry:symbol-registry-00529]{p.~\pageref{entry:symbol-registry-00529}} & P8.3 & Zeta maps fixed-point variables to state sets and Ext gives the set denoted by A; the P8.3 role is scoped away from the P6.3 forcing-name use.  \\
